\chapter[从集合的“大小”到测度：构造、正则性与传输]{从集合的“大小”到测度：\\构造、正则性与传输}\label{ch:02}

一个盒子的体积可以由边长相乘得到；有限个不交盒的总体积也没有疑问。真正的困难从
“有限”两个字消失时开始。一个稠密可数集能否仍然没有体积？用可数多个盒从外面覆盖一个集合时，
覆盖的最小总成本能否成为可数可加的测度？若两个测度只在区间或矩形上给出相同答案，怎样知道
它们在所有可测集上相同？又怎样把质量沿映射推到另一个空间，或者把两个空间的质量组合起来？

这些问题有一个共同形式：我们只能直接核算有限而简单的观测，却希望由它们构造、传输并唯一识别
一个对可数极限稳定的对象。本章从 Riemann 与 Jordan 理论的局限出发。外测度先给每个集合
一个外部价格，\Caratheodory{} 条件再筛出能够被无损切割的集合。生成类定理负责唯一性，Radon
正则性把抽象测度重新交还给开集和紧集；推前、乘积与 Stieltjes 构造则展示测度怎样运动、组合和
被一维函数编码。最后，我们回到 Euclidean 空间，建立换元公式所需的局部线性化、线性体积缩放和
Lipschitz 零集控制。

本章停留在集合、覆盖、有限和与有限维线性化的层面。一般测度积分从下一章开始；推前积分公式、
Radon--Nikodym 分解与非线性 Jacobian 公式将在相应工具建立后证明。Rademacher 定理以及 area、
coarea 公式则属于更深入的几何测度论。

记号与前置概念均沿用第~1~章：可测空间、可测映射与 Borel $\sigma$-代数见
\cref{def:1-4-2-measurable}；局部紧 Hausdorff 空间、函数支撑以及
$C_c(X),C_b(X),C_0(X)$ 见\cref{def:1-3-1-1,def:1-3-1-3}。本章首次调用这些记号时不重新定义，
但新增的测度、完成化、Radon 性与乘积测度均在使用前定义。

\section{从历史缺陷到可数可加性}\label{sec:02-01}

\subsection{Riemann 与 Jordan 理论的边界}\label{subsec:2-1-1}
\index{Riemann 积分}\index{Jordan 内容}

Riemann 积分把一个区间切成有限多段，在每段上用一个矩形估计函数图像。这个办法的力量和
限制来自同一个词：有限。有限网格善于处理边界规整的对象，却没有任何机制把可数多次微小误差
合并起来。下面的例子分别提出“哪些函数可以积分”“哪些集合可以测量”和
“哪些极限可以穿过积分号”。

十九世纪末到二十世纪初出现了两条互补路线。一条从集合长度出发，先扩大可测集合，再把函数拆成水平集；
Lebesgue 理论沿这条路把可数可加性放在中心。另一条由 Young、Daniell 等人的工作推动，先研究连续函数
或简单函数上的正线性泛函，再从单调极限恢复更大的积分。下面暂不采用第二条路线的抽象定理，但两条路线
共同指出：有限分割本身不足以控制可数极限。

\begin{definition}[Darboux 上下和]\label{def:2-1-1-1}
设 $f:[a,b]\to\R$ 有界，分割
$P=\{a=x_0<x_1<\cdots <x_N=b\}$。记
\[
 M_j=\sup_{x\in[x_{j-1},x_j]}f(x),\qquad
 m_j=\inf_{x\in[x_{j-1},x_j]}f(x),
\]
并定义
\[
 U(f,P)=\sum_{j=1}^N M_j(x_j-x_{j-1}),\qquad
 L(f,P)=\sum_{j=1}^N m_j(x_j-x_{j-1}).
\]
上积分是所有 $U(f,P)$ 的下确界，下积分是所有 $L(f,P)$ 的上确界。二者相等时，称
$f$ Riemann 可积，其公共值记为 $\int_a^b f(x)\,\dd x$。
\end{definition}

\begin{example}[Dirichlet 函数]\label{ex:2-1-1-1}
Dirichlet 函数 $d=\1_{\Q\cap[0,1]}$ 在每个非退化区间内同时取到 $0$ 和 $1$。所以对每个分割
$P$，
\[
 U(d,P)=1,\qquad L(d,P)=0.
\]
它不 Riemann 可积。另一方面，$d$ 取值 $1$ 的点只有可数多个；若一种面积理论允许可数个点
仍然没有面积，那么它理应给 $d$ 积分 $0$。
\end{example}

\begin{example}[Thomae 函数]\label{ex:2-1-1-2}
在 $[0,1]$ 上定义
\[
 t(x)=
 \begin{cases}
 q^{-1},&x=p/q\text{，其中 }p\in\Z, q\in\N_{>0},\ \gcd(|p|,q)=1,\\
 0,&x\notin\Q.
 \end{cases}
\]
无理数在每个区间中稠密，故每个下和为 $0$。给定 $\varepsilon>0$，取 $N$ 使
$1/N<\varepsilon/2$。分母不超过 $N$ 的既约有理数在 $[0,1]$ 中只有有限多个，记为
$r_1,\dots,r_s$。对每个 $j$ 取 $r_j\in(\alpha_j,\beta_j)$，并使
\[
 \sum_{j=1}^s(\beta_j-\alpha_j)<\frac{\varepsilon}{2}.
\]
令 $P$ 包含 $0,1$ 以及所有落在 $[0,1]$ 内的 $\alpha_j,\beta_j$。把 $P$ 的小区间分成两类：
内部与 $\bigcup_j[\alpha_j,\beta_j]$ 相交于正长度的那些，和其余小区间。第一类小区间的长度和
不超过 $\sum_j(\beta_j-\alpha_j)$；第二类小区间的内部不含任何 $r_j$，故在其上
$t\le 1/N<\varepsilon/2$（分割端点不影响上确界，因为每个 $r_j$ 位于所选区间的内部）。因此
\[
 U(t,P)
 \le \sum_{j=1}^s(\beta_j-\alpha_j)
      +\frac1N\sum_{I\in P}|I|
 <\frac{\varepsilon}{2}+\frac{\varepsilon}{2}=\varepsilon.
\]
因此 $t$ Riemann 可积且积分为 $0$。它在每个有理点不连续、在每个无理点连续；与 Dirichlet
函数比较可见，坏点是否稠密并不是决定因素，坏点集合的“大小”才是。

最后两项连续性也可直接核对。若 $x\notin\Q$ 且 $\eta>0$，取 $N$ 使 $1/N<\eta$。分母不超过
$N$ 的既约有理数只有有限多个，且都不等于 $x$；取 $\delta>0$ 小于 $x$ 到这有限集的距离。于是
$|y-x|<\delta$ 时，要么 $y$ 无理而 $t(y)=0$，要么 $y$ 的既约分母大于 $N$ 而
$0<t(y)<\eta$，故 $t(y)\to t(x)=0$。若 $x=p/q$ 为既约有理数，取无理数列
$y_k\to x$，则 $t(y_k)=0$，而 $t(x)=1/q>0$，所以 $t$ 在 $x$ 不连续。
\end{example}

为了把集合本身纳入这项比较，我们先把有限盒计算说清楚。以下各盒一律采用半开约定
$\prod_{j=1}^n[a_j,b_j)$；改变有限个端面不影响本节的结论。

\begin{lemma}[初等集体积的良定义性]\label{lem:jordan-elementary-volume-well-defined}
若同一个有界集合 $A\subset\R^n$ 分别写成有限个两两不交半开盒的并
$A=\bigsqcup_{i=1}^r Q_i=\bigsqcup_{j=1}^s R_j$，则
\[
 \sum_{i=1}^r |Q_i|=\sum_{j=1}^s|R_j|,
 \qquad |\textstyle\prod_k[a_k,b_k)|:=\prod_k(b_k-a_k).
\]
\end{lemma}

\begin{proof}
收集所有 $Q_i,R_j$ 在第 $k$ 个坐标方向出现的有限多个端点，并按这些端点切分一条包含全部盒的
大区间。各坐标切分的笛卡尔积给出有限共同网格。每个 $Q_i$ 和 $R_j$ 都是若干网格小盒的
不交并。记共同网格胞腔为 $(G_\alpha)_\alpha$。一维恒等式
\[
 b-a=\sum_{l=1}^N(c_l-c_{l-1})
 \quad(a=c_0<\cdots<c_N=b)
\]
逐坐标相乘并展开，给出
\[
 |Q_i|=\sum_{G_\alpha\subset Q_i}|G_\alpha|,
 \qquad
 |R_j|=\sum_{G_\alpha\subset R_j}|G_\alpha|.
\]
两种不交表示所包含的网格胞腔都恰为满足 $G_\alpha\subset A$ 的那些，故
\[
 \sum_i|Q_i|
 =\sum_{G_\alpha\subset A}|G_\alpha|
 =\sum_j|R_j|.
\]
\end{proof}

\begin{definition}[Jordan 可测与内容]\label{def:2-1-1-2}
有限个两两不交半开盒的并称为\emph{初等集}，其体积由上一引理定义。对有界
$E\subset\R^n$，令
\[
 J^*(E)=\inf\{|A|:A\supset E,\ A\text{ 为初等集}\},\qquad
 J_*(E)=\sup\{|B|:B\subset E,\ B\text{ 为初等集}\}.
\]
若 $J^*(E)=J_*(E)$，称 $E$ Jordan 可测，公共值记作 $J(E)$。
\end{definition}

\begin{example}[可数并超出 Jordan 定义域]\label{ex:2-1-1-3}
单点 $\{x\}$ 可放在边长任意小的盒中，所以 $J^*(\{x\})=0$。然而
$\Q\cap[0,1]$ 是这些单点的可数并，并且它的闭包及其补集在 $[0,1]$ 中的闭包都是 $[0,1]$。
下面的边界判据将给出 $\Q\cap[0,1]$ 不 Jordan 可测。失败的不是单点计算，而是 Jordan 理论的
定义域不承受可数并。
\end{example}

\begin{example}[移动尖峰]\label{ex:2-1-1-4}
例如取归一化帐篷函数
\[
 \phi(u)=
 \begin{cases}
  16(u-\tfrac14),&\tfrac14\le u\le\tfrac12,\\
  16(\tfrac34-u),&\tfrac12<u\le\tfrac34,\\
  0,&\text{其余}.
 \end{cases}
\]
它非负、属于 $C_c((0,1))$，并且两个三角形各有底 $1/4$、高 $4$，故
$\int_0^1\phi(u)\,\dd u=2\cdot\frac12\cdot\frac14\cdot4=1$。在 $[0,1]$ 上令
$f_n(x)=n\phi(nx)$，并把 $\phi$ 在 $\R\setminus(0,1)$ 上延为 $0$。每个 $f_n$ 连续。若
$x>0$，当 $n>1/x$ 时 $nx>1$，故 $f_n(x)=0$；在 $x=0$ 也有 $f_n(0)=0$。所以
$f_n(x)\to0$ 处处成立，而换元 $u=nx$ 给出
\[
 \int_0^1 f_n(x)\,\dd x=\int_0^n\phi(u)\,\dd u=1.
\]
逐点收敛本身不足以交换极限和积分，而且失败即使在连续函数列中也会发生。
\end{example}

\begin{remark}[可数稳定性的几个侧面]\label{def:2-1-1-3}
上述例子把新理论所需的稳定性分成几个彼此独立的方面：非负单调极限与适当受控极限下的连续性，
不交可数分解上的可加性，平移和线性变换下的协变性，以及由连续函数或几何简单集给出的逼近。
这些性质将在不同定理中分别建立；它们并非一组等价条件。
\end{remark}

\begin{proposition}[Jordan 边界判据]\label{proposition:2-1-1-2}
有界集合 $E\subset\R^n$ Jordan 可测，当且仅当 $J^*(\partial E)=0$。此时，对任一包含
$E$ 的闭盒 $Q$，指标函数 $\1_E:Q\to\R$ Riemann 可积，且
\[
 J(E)=\int_Q\1_E(x)\,\dd x.
\]
\end{proposition}

\begin{proof}
先设 $J^*(\partial E)=0$。若 $\partial E=\varnothing$，则不存在同时遇到 $E$ 与 $E^c$ 的线段：
否则对这样一条线段的参数集合
\[
 T=\{t\in[0,1]:(1-t)x+ty\in E\}\qquad(x\in E,\ y\notin E)
\]
取 $T$ 与其补集的分界点，就会得到 $\partial E$ 中的点。因此任意网格上 $\1_E$ 都没有振荡胞腔，
以下结论立即成立。现在设 $\partial E\ne\varnothing$。给定 $\varepsilon>0$，由外内容为零，先取
有限半开盒 $Q_1,\dots,Q_r$ 覆盖 $\partial E$，使
\[
 \sum_{j=1}^r|Q_j|<\frac{\varepsilon}{2}.
\]
把每个 $Q_j$ 略微外扩成开盒 $U_j\supset Q_j$，并逐个控制
$|U_j|<|Q_j|+\varepsilon/(2r)$，于是
\[
 \partial E\subset O:=\bigcup_{j=1}^rU_j,
 \qquad \sum_{j=1}^r|U_j|<\varepsilon.
\]
$E$ 有界，故 $\partial E$ 紧；闭集 $O^c$ 与紧集 $\partial E$ 不交，所以
\[
 \delta:=\dist(\partial E,O^c)>0.
\]
把 $Q$ 切成直径小于 $\delta$ 的有限网格。若一个闭网格小盒 $C$ 同时遇到 $E$ 与
$Q\setminus E$，连接两点的线段包含某个 $z\in\partial E\cap C$。于是对每个 $w\in C$，
$|w-z|<\delta$，故 $w\notin O^c$，即 $C\subset O$。因此所有使 $\1_E$ 振幅为 $1$ 的小盒的并
包含于 $O$；用共同端点网格比较有限盒体积，得到这些小盒的总体积不超过
$\sum_j|U_j|$。于是
\[
 U(\1_E,P)-L(\1_E,P)\le\sum_j|U_j|<\varepsilon.
\]
于是 $\1_E$ Riemann 可积。更显式地，令 $\mathscr P$ 为该有限网格的胞腔族，并置
\[
 B_P=\bigcup_{C\in\mathscr P:\,\inf_C\1_E=1}C,
 \qquad
 A_P=\bigcup_{C\in\mathscr P:\,\sup_C\1_E=1}C.
\]
网格面的有限重叠为零内容，不影响这些初等集的体积，因而
\[
 B_P\subset E\subset A_P,qquad
 |B_P|=L(\1_E,P),\qquad |A_P|=U(\1_E,P).
\]
所以
\[
 0\le |A_P|-|B_P|=U(\1_E,P)-L(\1_E,P)<\varepsilon.
\]
这逐式给出 Jordan 内外逼近。由 Jordan 内容的定义以及 Darboux 可积性，分别有
\[
 |B_P|\le J(E)\le |A_P|,
 \qquad
 |B_P|=L(\1_E,P)\le\int_Q\1_E\le U(\1_E,P)=|A_P|.
\]
所以
\[
 \left|J(E)-\int_Q\1_E\right|
 \le |A_P|-|B_P|<\varepsilon;
\]
令 $\varepsilon\downarrow0$ 即得 $J(E)=\int_Q\1_E$。

反过来，设 $E$ Jordan 可测。取初等集 $B\subset E\subset A$ 且
$|A|-|B|<\varepsilon$，并取两者的共同网格。若 $x\in\partial E$ 落在某个网格小盒的内部，那么该
小盒若包含于 $B$，则 $x$ 是 $E$ 的内点；若不包含于 $A$，则 $x$ 是 $E^c$ 的内点。因此该小盒
必包含在 $A\setminus B$。这些胞腔两两不交，其总体积不超过 $|A|-|B|<\varepsilon$。网格边界是
有限个集合
\[
 \{x\in Q:x_i=c\}\qquad(1\le i\le n, c\text{ 为某个网格坐标}),
\]
每个都可用法向厚度为 $2\delta$ 的有限薄盒覆盖；有限求和后其总成本随 $\delta\downarrow0$ 趋于
零。若 $H$ 表示网格边界、$C_1,\dots,C_s$ 表示上述位于 $A\setminus B$ 的胞腔，则
\[
 \partial E\subset H\cup\bigcup_{l=1}^sC_l,
 \qquad
 J^*(\partial E)
 \le J^*(H)+\sum_{l=1}^s|C_l|
 \le J^*(H)+|A|-|B|.
\]
令 $\delta\downarrow0$ 得 $J^*(\partial E)\le\varepsilon$，再令 $\varepsilon\downarrow0$，得到
$J^*(\partial E)=0$。
\end{proof}

\begin{proposition}[Jordan 内容的有限性边界]\label{proposition:2-1-1-3}
Jordan 内容对有限个两两不交 Jordan 可测集可加，但 Jordan 可测集不构成 $\sigma$-代数。
具体地，每个单点内容为零，而 $\Q\cap[0,1]$ 的边界是 $[0,1]$，因而不 Jordan 可测。
\end{proposition}

\begin{proof}
若 $E_1,\dots,E_r$ 两两不交且 Jordan 可测，固定误差 $\varepsilon>0$，分别以初等集
$B_j\subset E_j\subset A_j$ 夹逼，使 $|A_j|-|B_j|<\varepsilon/r$。共同网格细分后，
$\bigsqcup_jB_j\subset\bigsqcup_jE_j\subset\bigcup_jA_j$；外并的体积至多各 $|A_j|$ 之和，
内并的体积等于各 $|B_j|$ 之和。于是
\[
 \sum_{j=1}^r|B_j|
 \le J_*\!\left(\bigsqcup_{j=1}^rE_j\right)
 \le J^*\!\left(\bigsqcup_{j=1}^rE_j\right)
 \le\left|\bigcup_{j=1}^rA_j\right|
 \le\sum_{j=1}^r|A_j|.
\]
两端之差小于 $\varepsilon$；令 $\varepsilon\downarrow0$ 得
$J(\bigsqcup_jE_j)=\sum_jJ(E_j)$。第二个断言由单点计算和上一命题得到。
\end{proof}

\begin{proposition}[单调极限迫使零集可数稳定]\label{proposition:2-1-1-4}
设 $I$ 定义在一个非负函数锥上，该锥含
$\1_{E_k},\1_{\cup_{k\le n}E_k},\1_{\cup_kE_k}$，对有限非负线性组合和有界递增极限封闭。
若 $I$ 可加、非负齐次、保持次序，并满足 $g_n\uparrow g$ 时 $I(g_n)\uparrow I(g)$，则
\[
 I(\1_{E_k})=0\quad(k\ge1)
 \quad\Longrightarrow\quad I(\1_{\cup_kE_k})=0.
\]
\end{proposition}

\begin{proof}
对每个 $n$，点态不等式
$\1_{\cup_{k\le n}E_k}\le\sum_{k=1}^n\1_{E_k}$ 与正性、有限线性给出
\[
 0\le I(\1_{\cup_{k\le n}E_k})
 \le\sum_{k=1}^nI(\1_{E_k})=0.
\]
左端指标函数随 $n$ 递增至 $\1_{\cup_kE_k}$，再用单调极限性质得到结论。
\end{proof}

\begin{theorem}[Lebesgue 的 Riemann 判据]\label{theorem:2-1-1-1}
有界函数 $f:[a,b]\to\R$ Riemann 可积，当且仅当它的不连续点集合具有 Lebesgue 测度零。
\end{theorem}

Lebesgue 测度将在\cref{subsec:2-1-3}中构造。判据的证明还需要函数振幅与微分理论，见
\cref{theorem:4-1-2-3}。

\begin{figure}[htbp]
\centering
\begin{tikzpicture}[node distance=13mm and 22mm,
  box/.style={draw,rounded corners,align=center,minimum width=31mm,minimum height=12mm},
  bad/.style={draw,dashed,rounded corners,align=center,inner sep=4pt},
  arr/.style={-{Latex[length=2mm]},thick}]
 \node[box] (j) {有限分割\\Jordan 内容};
 \node[box,right=of j] (l) {可数集合运算\\Lebesgue 测度};
 \node[box,right=of l] (d) {单调极限\\积分};
 \node[bad,below=of j] (q) {稠密可数并};
 \node[bad,below=of l] (dir) {Dirichlet 函数};
 \node[bad,below=of d] (sp) {移动尖峰};
 \draw[arr] (q)--(j);
 \draw[arr] (dir)--(l);
 \draw[arr] (sp)--(d);
 \draw[arr] (j)--node[midway,above,font=\small,fill=white,inner sep=1pt]{扩大定义域}(l);
 \draw[arr] (l)--node[midway,above,font=\small,align=center,fill=white,inner sep=1pt]{建立\\极限定理}(d);
\end{tikzpicture}
\caption{有限分割、可数集合运算与极限过程之间的三种失配。}
\label{fig:jordan-lebesgue-daniell-failures}
\end{figure}

\begin{remark}
Jordan 证明中真正失效的步骤是“取一个共同有限网格”。可数多个单点各自拥有小盒，不会自动产生
一个仍可由有限网格控制的覆盖。下一小节允许可数覆盖，并把近优覆盖的误差写成可求和的序列。
移动尖峰则只排除了“无控制的逐点收敛足以交换积分与极限”；一致收敛是一种修补，却不是唯一修补，第~3~章
还会建立单调收敛和受控收敛的不同条件。
\end{remark}

这些失败并不限于教材中的反例。Fourier 级数逼近和概率中的随机变量极限都天然先给逐点或几乎处处
信息；若积分要在这些极限下稳定，就必须补上可核查的控制条件。Young--Daniell 路线把这种压力写成
正泛函的单调连续性，集合路线则要求零集对可数并稳定。\cref{proposition:2-1-1-4} 已证明两者之间
最基本的联系；完整积分定理仍留到第~3~章。

\begin{exercise}
把上述 Thomae 函数证明改写成 Darboux 振幅判据的形式，并证明它恰在无理点连续。
\end{exercise}
\begin{exercise}
证明有限集合 Jordan 内容为零，并用边界判据再次证明 $\Q\cap[0,1]$ 不 Jordan 可测。
\end{exercise}
\begin{exercise}
构造一列连续非负函数，逐点趋于零但积分恒为给定常数 $c>0$。
\end{exercise}
\begin{exercise}
构造 Riemann 可积函数列，其逐点极限为 Dirichlet 函数。
\end{exercise}

有限网格已经无法同时追踪可数覆盖。下一小节保留“简单集合的覆盖成本”，但把覆盖数目放宽为可数，并
用一个切割条件从定义在所有子集上的外部价格中筛出真正可加的集合。

\subsection{外测度与 \Caratheodory{} 可测性}\label{subsec:2-1-2}
\index{外测度}\index{Caratheodory@\Caratheodory{} 可测性}

Jordan 理论试图从集合内部和外部同时有限夹逼。现在换一个次序：先不问集合是否可测，只允许用
简单集合从外面覆盖，并把所有覆盖成本的下确界当作外部价格。这个价格天然服从可数次可加，却
通常不能在所有集合上相加；可测性将由“切开任何测试集合都不损失成本”产生。

Lebesgue 的原始构造围绕实线上区间长度展开；\Caratheodory{} 的形式把其中真正使用的机制抽成外测度与
无损切割条件。这样，Euclidean 体积、Hausdorff 尺度和由预测度生成的测度可以共享一套证明。可测类
不是事先指定的全部子集，而是由“对每个测试集切割都不增加成本”挑出的、与该外测度相容的自然定义域。

在引入外部价格以前，先固定“测度”本身的含义。上一章已经定义了 $\sigma$-代数；现在增加的要求是
不交可数分解上的精确相加。

\begin{definition}[测度与测度空间]\label{def:2-1-2-measure-space}
设 $\mathcal A$ 是集合 $X$ 上的 $\sigma$-代数。函数
$\mu:\mathcal A\to[0,\infty]$ 称为\emph{测度}，若 $\mu(\varnothing)=0$，并且对任意两两不交的
$E_k\in\mathcal A$ 都有
\[
 \mu\Bigl(\bigsqcup_{k=1}^{\infty}E_k\Bigr)
 =\sum_{k=1}^{\infty}\mu(E_k).
\]
三元组 $(X,\mathcal A,\mu)$ 称为测度空间。若 $\mu(X)<\infty$，称 $\mu$ 为有限测度；若
$X=\bigcup_kX_k$ 且 $X_k\in\mathcal A$、$\mu(X_k)<\infty$，称 $\mu$ 为 $\sigma$-有限测度；
若 $\mu(X)=1$，称其为概率测度。拓扑空间上定义于 $\Borel(X)$ 的测度称为 Borel 测度。
性质 $P(x)$ 称为在 $\mu$-\emph{几乎处处}成立，若存在 $N\in\mathcal A$ 满足 $\mu(N)=0$，且
$P(x)$ 对每个 $x\in X\setminus N$ 成立；在未完成的测度空间中，这一定义不要求例外点集本身可测。
\end{definition}

可数可加性不只计算不交并；它还控制单调集合列的极限。后文的正则逼近和测度唯一性都会直接使用
下面两个公式。

\begin{proposition}[测度的单调连续性]\label{proposition:2-1-2-measure-continuity}
设 $(X,\mathcal A,\mu)$ 为测度空间。
\begin{enumerate}
 \item 若 $E_n\uparrow E$，则 $\mu(E_n)\uparrow\mu(E)$；
 \item 若 $E_n\downarrow E$ 且 $\mu(E_1)<\infty$，则 $\mu(E_n)\downarrow\mu(E)$。
\end{enumerate}
\end{proposition}

\begin{proof}
对递增列令 $D_1=E_1$、$D_n=E_n\setminus E_{n-1}$（$n\ge2$）。这些集合两两不交，且
$E=\bigsqcup_nD_n$、$E_N=\bigsqcup_{n\le N}D_n$，所以
\[
 \mu(E_N)=\sum_{n\le N}\mu(D_n)
 \uparrow\sum_{n=1}^{\infty}\mu(D_n)=\mu(E).
\]
若 $E_n\downarrow E$ 且 $\mu(E_1)<\infty$，则 $E_1\setminus E_n\uparrow E_1\setminus E$。
由第一项和有限可加性，
\[
 \mu(E_1)-\mu(E_n)=\mu(E_1\setminus E_n)
 \uparrow\mu(E_1\setminus E)=\mu(E_1)-\mu(E).
\]
这里所有被减数都不超过有限数 $\mu(E_1)$，故没有出现 $\infty-\infty$；移项即得第二项。
\end{proof}

\begin{example}[有限原子上的直接核算]\label{ex:2-1-2-finite-measure}
在 $X=\{1,\ldots,r\}$ 上取 $\mathcal A=2^X$，给定权重 $a_j\ge0$，并置
\[
 \mu(E)=\sum_{j\in E}a_j.
\]
任意不交可数族中至多有 $r$ 个非空成员，故可数可加性退化为有限和的重排，$\mu$ 是有限测度；当
$\sum_ja_j=1$ 时它是概率测度。这个模型说明测度公理怎样把有限权重相加，而外测度接下来要解决的
问题是：尚未给出可测 $\sigma$-代数时，怎样先为任意子集规定一致的外部价格。
\end{example}

\begin{definition}[外测度]\label{def:2-1-2-1}
集合 $X$ 上的函数 $\mu^*:2^X\to[0,\infty]$ 称为外测度，如果
\begin{enumerate}
 \item $\mu^*(\varnothing)=0$；
 \item $A\subset B$ 蕴含 $\mu^*(A)\le\mu^*(B)$；
 \item 对任意集合列 $(E_k)$，
 \(
 \mu^*(\bigcup_kE_k)\le\sum_k\mu^*(E_k).
 \)
\end{enumerate}
\end{definition}

\begin{definition}[覆盖生成的外测度]\label{def:2-1-2-3}
给定简单集合族 $\mathcal C\subset2^X$ 和成本 $c:\mathcal C\to[0,\infty]$，允许有限或可数覆盖，
并约定空族以成本 $0$ 覆盖空集，定义
\[
 \mu^*(E)=\inf\left\{\sum_{j=1}^{N}c(C_j):
 E\subset\bigcup_{j=1}^{N}C_j,\ C_j\in\mathcal C,\ 1\le N\le\infty\right\}.
\]
若不存在这样的覆盖，下确界取 $\infty$。
\end{definition}

\begin{lemma}[覆盖下确界给出外测度]\label{lemma:2-1-2-1}
若成本非负且空集有零成本空覆盖，则上一式定义外测度。
\end{lemma}

\begin{proof}
空集条件由约定得到。若 $A\subset B$，每个覆盖 $B$ 的族也覆盖 $A$，故下确界满足单调性。
只需证明可数次可加。若某个 $\mu^*(E_k)=\infty$，所需不等式的右端为无穷。否则，给定
$\varepsilon>0$，为每个 $k$ 选择一个有限或可数覆盖
$(C_{k,j})_{1\le j\le N_k}$（$N_k\le\infty$），使
\[
 \sum_{j=1}^{N_k}c(C_{k,j})<\mu^*(E_k)+\varepsilon2^{-k}.
\]
当 $E_k=\varnothing$ 时允许取空覆盖。所有这些有限或可数族的并仍是可数族，并覆盖
$\bigcup_kE_k$。于是
\[
 \mu^*\Bigl(\bigcup_kE_k\Bigr)
 \le\sum_{k,j}c(C_{k,j})
 \le\sum_k\mu^*(E_k)+\varepsilon.
\]
令 $\varepsilon\downarrow0$。这里同时为可数多个 $k$ 选择近优覆盖使用上一章
定义~\ref{def:1-4-3-1} 中的可数选择；证明并未假设任何下确界达到。
\end{proof}

\begin{example}[三种外部价格]\label{ex:2-1-2-1}
在 $\R^n$ 中以半开盒为 $\mathcal C$、以盒体积为成本，所得外测度记作 $m^*$。若
$E=\{x_1,x_2,\dots\}$ 可数，给第 $k$ 个点取体积小于 $\varepsilon2^{-k}$ 的盒，便有
$m^*(E)\le\varepsilon$，所以 $m^*(E)=0$。

另一个例子是计数外测度
\[
 \#^*(E)=\begin{cases}|E|,&E\text{ 有限},\\ \infty,&E\text{ 无限}.
 \end{cases}
\]
它稍后会成为所有集合均可测的情形。若改以直径的 $s$ 次幂为小尺度覆盖成本，则得到
Hausdorff 外测度；下一例说明这种成本如何反映集合的尺度。
\end{example}

\begin{example}[Hausdorff 外测度的覆盖动机]\label{ex:2-1-2-2}
若只用直径不超过 $\delta$ 的集合覆盖 $E$，并把成本取作
$\sum_j(\diam U_j)^s$，把一条长度为 \(r\) 的线段等分为 \(N\) 段便得到总成本
\[
 N\left(\frac rN\right)^s=r^sN^{1-s}.
\]
在 \(s=1\) 时该成本恒为 \(r\)，在 \(s>1\) 时则随 \(N\to\infty\) 趋于零；单点在每个
\(s>0\) 时也有零成本。先固定 $\delta$ 取覆盖下确界，再令 $\delta\downarrow0$，便得到能分辨尺度的
外部价格。\Cref{subsec:2-4-3} 将给出正式定义及 Lipschitz 映射下的覆盖畸变不等式。
\end{example}

\begin{definition}[\Caratheodory{} 可测]\label{def:2-1-2-2}
给定外测度 $\mu^*$，集合 $E\subset X$ 称为 $\mu^*$-可测，如果对每个 $A\subset X$，
\[
 \mu^*(A)=\mu^*(A\cap E)+\mu^*(A\setminus E).
 \label{eq:caratheodory-cut}
\]
\end{definition}

外测度次可加已经给出等式中“$\le$”的一边；定义要求切开后的总成本不超过整体最优成本。

\begin{example}[可测性依赖外测度]\label{ex:2-1-2-3}
对计数外测度，若 $A$ 有限，等式 \eqref{eq:caratheodory-cut} 是有限集合的基数分解；若 $A$ 无限，
右侧至少一项为无穷，等式仍成立。因此每个集合都可测。
\end{example}

\begin{example}[不是每个外测度都测量所有集合]\label{ex:2-1-2-4}
设 $X$ 至少有两个点，令 $\mu^*(\varnothing)=0$，每个非空集合的外测度为 $1$。它满足外测度
三条公理。若 $E$ 为真非空子集，以 $A=X$ 测试，则
\[
 1=\mu^*(X)\ne\mu^*(E)+\mu^*(X\setminus E)=2.
\]
所以 $E$ 不可测。这说明 \Caratheodory{} 可测性不是集合脱离外测度后的绝对属性。
\end{example}

\begin{proposition}[只需有限成本测试]\label{proposition:2-1-2-4}
检验 $E$ 的 \Caratheodory{} 条件时，只需考虑 $\mu^*(A)<\infty$ 的测试集。
\end{proposition}

\begin{proof}
当 $\mu^*(A)=\infty$ 时，次可加性给
$\infty\le\mu^*(A\cap E)+\mu^*(A\setminus E)$，所以右端也是无穷，等式自动成立。
\end{proof}

\begin{theorem}[\Caratheodory{} 定理]\label{theorem:2-1-2-2}
所有 $\mu^*$-可测集构成一个 $\sigma$-代数 $\mathcal M$，并且
$\mu=\mu^*|_{\mathcal M}$ 是完备测度。特别地，若 $N\in\mathcal M$ 且 $\mu(N)=0$，则
$N$ 的每个子集都属于 $\mathcal M$ 且测度为零。
\end{theorem}

\begin{proof}
空集满足切割等式，补集稳定来自等式对 $E$ 和 $E^c$ 的对称性。设 $E,F$ 可测。对任意测试集
$A$，先以 $E$ 切割，再以 $F$ 切割 $A\setminus E$，得到
\[
\begin{aligned}
 \mu^*(A)
 &=\mu^*(A\cap E)+\mu^*(A\setminus E)\\
 &=\mu^*(A\cap E)+\mu^*((A\setminus E)\cap F)
   +\mu^*(A\setminus(E\cup F))\\
 &\ge \mu^*(A\cap(E\cup F))+\mu^*(A\setminus(E\cup F)).
\end{aligned}
\]
最后一步使用
$A\cap(E\cup F)=(A\cap E)\cup((A\setminus E)\cap F)$ 的次可加性。反向不等式由把 $A$ 分成
这两块的次可加性得到，所以 $E\cup F$ 可测；归纳得到有限并稳定。

现在设 $E_k\in\mathcal M$，令
$D_1=E_1$、$D_k=E_k\setminus\bigcup_{j<k}E_j$。有限并和补集稳定保证 $D_k$ 可测且两两
不交。记 $D=\bigcup_kD_k$。对任意测试集 $A$，逐次用 $D_1,\dots,D_n$ 切割，得到
\[
 \mu^*(A)\ge
 \sum_{k=1}^n\mu^*(A\cap D_k)
 +\mu^*\Bigl(A\setminus\bigcup_{k=1}^nD_k\Bigr).
 \label{eq:caratheodory-iterate}
\]
最后一项不小于 $\mu^*(A\setminus D)$。令 $n\to\infty$，
\[
 \mu^*(A)\ge\sum_{k=1}^{\infty}\mu^*(A\cap D_k)+\mu^*(A\setminus D)
 \ge\mu^*(A\cap D)+\mu^*(A\setminus D),
\]
第二个不等式是外测度次可加。反向不等式同样来自次可加，故 $D$ 可测。这证明 $\mathcal M$
为 $\sigma$-代数。

在 \eqref{eq:caratheodory-iterate} 中取 $A=D$ 并令 $n\to\infty$，与次可加的反向估计合并，
得到 $\mu(D)=\sum_k\mu(D_k)$，所以限制是测度。

最后，设 $N$ 可测且 $\mu^*(N)=0$，并取 $S\subset N$。对任意 $A$，单调性给
$\mu^*(A\cap S)=0$。又因为
$A=(A\setminus S)\cup(A\cap S)$，
\[
 \mu^*(A)\le\mu^*(A\setminus S)+0\le\mu^*(A).
\]
所以 $S$ 满足 \Caratheodory{} 等式，并且零测。
\end{proof}

\begin{definition}[度量外测度]\label{def:2-1-2-4}
度量空间上的外测度 $\mu^*$ 若在 $\dist(A,B)>0$ 时满足
$\mu^*(A\cup B)=\mu^*(A)+\mu^*(B)$，称为\emph{度量外测度}（metric outer measure）。
\end{definition}

\begin{proposition}[度量外测度的 Borel 性]\label{proposition:2-1-2-3}
度量外测度的每个闭集都 \Caratheodory{} 可测，因而每个 Borel 集可测。
\end{proposition}

\begin{proof}
设 $F$ 闭。由 \cref{proposition:2-1-2-4}，只需取 $\mu^*(A)<\infty$。令
\[
 C_j=A\cap\{2^{-j}\le\dist(x,F)<2^{-j+1}\}\quad(j\ge1),
\]
并把 $\dist(x,F)\ge1$ 的部分记为 $C_0$。若 $k\ge j+2$，则对
$x\in C_j,y\in C_k$ 有
\[
 \dist(x,F)-\dist(y,F)
 \ge2^{-j}-2^{-k+1}\ge2^{-j-1}.
\]
距离函数是 $1$-Lipschitz，故 $d(x,y)\ge2^{-j-1}>0$。因此同为奇指标或同为偶指标的任意有限
子族彼此可逐项使用度量可加性，给出
\[
 \sum_{j\le N,\ j\text{ odd}}\mu^*(C_j)\le\mu^*(A),\qquad
 \sum_{j\le N,\ j\text{ even}}\mu^*(C_j)\le\mu^*(A).
\]
令 $N\to\infty$ 后相加，得到
\[
 \sum_{j=1}^{\infty}\mu^*(C_j)\le2\mu^*(A)<\infty,
\]
所以该级数收敛。

令 $B_n=A\cap\{\dist(x,F)\ge2^{-n}\}$。由
\[
 A\setminus F\subset B_n\cup\bigcup_{j>n}C_j
\]
和外测度次可加性，
\[
 \mu^*(A\setminus F)
 \le \mu^*(B_n)+\sum_{j>n}\mu^*(C_j).
 \label{eq:metric-outer-ring-tail}
\]
另一方面，$B_n$ 与 $A\cap F$ 正距离，故
\[
 \mu^*(A)\ge\mu^*(B_n\cup(A\cap F))
 =\mu^*(B_n)+\mu^*(A\cap F).
\]
集合 $B_n$ 递增，故数列 $b_n=\mu^*(B_n)$ 递增；又 $B_n\subset A$ 且 $\mu^*(A)<\infty$，
所以 $b_n\uparrow b<\infty$。在 \eqref{eq:metric-outer-ring-tail} 中令 $n\to\infty$，环带级数的
尾和趋零，得到 $\mu^*(A\setminus F)\le b$。另一方面，上式的正距离加性对每个 $n$ 给
$b_n+\mu^*(A\cap F)\le\mu^*(A)$；令 $n\to\infty$ 后合并两式，得到
\[
 \mu^*(A)\ge\mu^*(A\setminus F)+\mu^*(A\cap F).
\]
反向由次可加，故 $F$ 可测。闭集生成 Borel $\sigma$-代数，结论随之成立。
\end{proof}

\begin{definition}[测度空间的完成化]\label{def:measure-completion}
设 $(X,\mathcal A,\mu)$ 是测度空间，令
\[
 \begin{aligned}
 \overline{\mathcal A}^{\mu}
 &=\{E\subset X:E\mathbin\triangle A\subset Z
       \text{ 对某些 }A,Z\in\mathcal A,\ \mu(Z)=0\}\\
 &=\{A\cup N:A\in\mathcal A,\ N\subset Z\in\mathcal A,\ \mu(Z)=0\}.
 \end{aligned}
\]
若 $E\mathbin\triangle A$ 包含于可测零集，置 $\bar\mu(E)=\mu(A)$，称所得测度空间为原空间的
完成化。
\end{definition}

\begin{proposition}[完成化的代表]\label{prop:measure-completion-representative}
$\overline{\mathcal A}^{\mu}$ 是 $\sigma$-代数，上述 $\bar\mu$ 良定义且完备。每个完成化可测集
$E$ 都有 $A\in\mathcal A$ 使 $E\mathbin\triangle A$ 包含于一个 $\mu$-零集。
\end{proposition}

\begin{proof}
若 $E=A\cup N$，则 $E\mathbin\triangle A=N\setminus A$ 包含于可测零集。反过来，若
$E\mathbin\triangle A\subset Z$ 且 $Z$ 可测零，则
$E=(A\setminus Z)\cup(E\cap Z)$，所以定义中的两种表示等价。补集和可数并时，可测代表分别取补集和
可数并，误差仍包含于可数个可测零集之并，故该类为 $\sigma$-代数。若两个表示给出代表 $A,B$，
则 $A\mathbin\triangle B$ 包含于可测零集，于是
$\mu(A\setminus B)=\mu(B\setminus A)=0$，从可加性得到 $\mu(A)=\mu(B)$；故 $\bar\mu$ 良定义。

设两两不交的 $E_i$ 有可测代表 $A_i$，且 $E_i\mathbin\triangle A_i\subset Z_i$，其中
$\mu(Z_i)=0$。同时选取可数族 $(A_i,Z_i)$ 是一次可数选择；在通常的 \textsf{ZFC} 背景下这是
允许的。令
\[
 B_1=A_1,\qquad B_i=A_i\setminus\bigcup_{j<i}A_j\quad(i\ge2).
\]
则 $B_i$ 可测且两两不交。由于 $E_i\cap E_j=\varnothing$，有
$A_i\cap A_j\subset Z_i\cup Z_j$；从而 $A_i\mathbin\triangle B_i$ 包含于有限个零集之并，
故 $B_i$ 仍是 $E_i$ 的代表。因此
\[
 \left(\bigcup_iE_i\right)\mathbin\triangle
 \left(\bigcup_iB_i\right)
 \subset\bigcup_i(E_i\mathbin\triangle B_i)
\]
包含于可测零集 $\bigcup_iZ_i'$，其中每个 $Z_i'$ 是相应的有限零集并；故
$\bigsqcup_iB_i$ 确实是 $\bigsqcup_iE_i$ 的可测代表。于是
\[
 \bar\mu\Bigl(\bigsqcup_iE_i\Bigr)
 =\mu\Bigl(\bigsqcup_iB_i\Bigr)
 =\sum_i\mu(B_i)=\sum_i\bar\mu(E_i).
\]
这证明可数可加性。若 $\bar\mu(E)=0$，取零测可测代表 $A$；则 $E$ 包含于 $A$ 与表示误差零集之并，
所以 $E$ 的任意子集仍属于完成化且测度为零。最后一句已包含在上述对称差表示中。
\end{proof}

\begin{proposition}[完成化可测函数的代表]\label{prop:completion-measurable-function-representative}
设 $\bar\mu$ 是 $(X,\mathcal A,\mu)$ 的完成化。若
$f:X\to[-\infty,\infty]$ 满足
$f^{-1}(B)\in\overline{\mathcal A}^{\mu}$ 对每个 Borel 集
$B\subset[-\infty,\infty]$ 成立，则存在具有
$g^{-1}(B)\in\mathcal A$ 的 $g:X\to[-\infty,\infty]$，使 $f=g$ 在某个 $\mu$-零集之外成立。
\end{proposition}

\begin{proof}
先设 $0\le f\le1$。令
\[
 s_n=2^{-n}\lfloor 2^nf\rfloor,
\]
并在 $f=1$ 处令 $s_n=1$。$s_n$ 只取有限多个值；它的每个水平集属于完成化，因而有
$\mathcal A$ 中的代表。具体地，把 $s_n$ 的不同值列为 $v_{n,0},\dots,v_{n,r_n}$，令
$L_{n,j}=\{s_n=v_{n,j}\}$，并取 $A_{n,j}\in\mathcal A$ 与可测零集 $Z_{n,j}$，使
\[
 L_{n,j}\mathbin\triangle A_{n,j}\subset Z_{n,j}.
\]
指标对 $(n,j)$ 的集合可数；同时作这些选择使用可数选择。按固定次序不交化：
\[
 B_{n,0}=A_{n,0},\qquad
 B_{n,j}=A_{n,j}\setminus\bigcup_{l<j}A_{n,l}\quad(1\le j\le r_n),
\]
并在余集 $X\setminus\bigcup_jB_{n,j}$ 上指定默认值 $v_{n,0}$。令 $t_n=v_{n,j}$ 于
$B_{n,j}$。这是 $\mathcal A$-可测有限值函数。若
$x\notin Z_n:=\bigcup_jZ_{n,j}$，则 $x$ 对所有 $A_{n,j}$ 的隶属关系与对 $L_{n,j}$ 的隶属关系
相同；水平集 $(L_{n,j})_j$ 恰好分割 $X$，故此时恰有一个代表含 $x$，且
$t_n(x)=s_n(x)$。

置 $Z=\bigcup_nZ_n$ 和 $g=\limsup_nt_n$。为显式核对可测性，对每个 $a\in\R$ 有
\[
 \{g>a\}
 =\bigcup_{\substack{q\in\Q\\q>a}}
   \bigcap_{N=1}^{\infty}\bigcup_{n\ge N}\{t_n>q\}\in\mathcal A.
\]
这些半直线生成扩展实线的 Borel $\sigma$-代数，故 $g$ 是 $\mathcal A$-可测的。在
$X\setminus Z$ 上，
$t_n=s_n\to f$，所以 $g=f$。

一般情形取把扩展实线双射到 $[0,1]$ 的严格递增连续函数，例如在有限点置
$\psi(x)=\tfrac12+\pi^{-1}\arctan x$，并令 $\psi(-\infty)=0,\psi(\infty)=1$。对
$\psi\circ f$ 使用上一步，再与 Borel 可测的 $\psi^{-1}$ 复合。
\end{proof}

\begin{figure}[htbp]
\centering
\begin{tikzpicture}[scale=.95]
 \path[fill=BookAccent!10]
   (-3,-1.6)--(0,-1.6)..controls(-1,-.5)and(1,.5)..(0,1.6)--(-3,1.6)--cycle;
 \path[fill=BookWarm!12]
   (0,-1.6)--(3,-1.6)--(3,1.6)--(0,1.6)..controls(1,.5)and(-1,-.5)..(0,-1.6)--cycle;
 \draw[rounded corners,thick] (-3,-1.6) rectangle (3,1.6);
 \draw[BookInk,thick] (0,-1.6)..controls(-1,-.5)and(1,.5)..(0,1.6);
 \node at (-1.65,0) {$A\cap E$};
 \node at (1.65,0) {$A\setminus E$};
 \node[above] at (0,1.6) {测试集 $A$};
 \draw[-{Latex[length=2mm]},thick] (0,-2.15)--(0,-1.55);
 \node[below,align=center] at (0,-2.15)
 {$\mu^*(A)=\mu^*(A\cap E)+\mu^*(A\setminus E)$\\表示分别覆盖不增加最优成本};
\end{tikzpicture}
\caption{\Caratheodory{} 条件把任意测试集无损切成两块。}
\label{fig:caratheodory-lossless-cut}
\end{figure}

\begin{remark}
外测度与测度不可混同。外测度定义在所有子集上，却通常不加；真正的测度是它在可测 $\sigma$-代数
上的限制。相反，完成化是在已经有测度之后，把可测零集的所有子集补入定义域。
\end{remark}

覆盖外测度将在本章产生 Lebesgue 测度，在\cref{subsec:2-4-3}产生 Hausdorff 尺度，也可用于路径
空间或乘积空间上的外部构造。另一条存在性路线是先在环上给预测度，再作 \Caratheodory{} 扩张；它在
\cref{subsec:2-2-1}正式建立。两条路线都用可数覆盖，但“从所有子集中筛选”与“从简单集合向外延拓”
承担不同任务。

\begin{exercise}
从覆盖定义逐项验证 Lebesgue 外测度的平移不变性。
\end{exercise}
\begin{exercise}
证明零外测度集合满足 \Caratheodory{} 条件。
\end{exercise}
\begin{exercise}
在 $|X|\ge2$ 时给出一个外测度，使恰有 $\varnothing$ 与 $X$ 可测。
\end{exercise}
\begin{exercise}
补写度量外测度命题中环带内缩和奇偶尾估计的全部细节，并据此证明开集可测。
\end{exercise}

外测度给出了抽象筛选机制；要把它变成 Euclidean 体积，还需从盒覆盖中证明平移、缩放以及开紧逼近，
而不能把这些几何性质当作定义的一部分。

\subsection{Lebesgue 测度：Euclidean 空间中的正则体积}\label{subsec:2-1-3}
\index{Lebesgue 测度}\index{正则性}

现在把覆盖成本取成 Euclidean 盒体积。仅有可数可加性还不足以刻画体积：它还应尊重
平移与缩放，并允许用开集从外、紧集从内逼近。这些性质来自盒几何、紧性和
Euclidean 空间的可数拓扑骨架。

Lebesgue 理论的关键进展，是先建立一个对可数集合运算稳定的 Borel 测度，
再把 Borel 零集的所有子集补入完成化。完成化扩大了可测对象，却保留几何体积；Jordan 内容和 Riemann
积分则应当作为规整边界上的旧理论被重新识别，而不是与新理论平行共存、互不相干。

\begin{definition}[Lebesgue 外测度与测度]\label{def:2-1-3-1}
对半开盒 $Q=\prod_{j=1}^n[a_j,b_j)$ 置 $|Q|=\prod_j(b_j-a_j)$，并令
\[
 m^*(E)=\inf\left\{\sum_k|Q_k|:E\subset\bigcup_kQ_k\right\}.
\]
把 $m^*$ 的 \Caratheodory{} 可测类暂记为 $\mathcal L_C$，限制测度暂记为 $m_C$。
\end{definition}

\begin{definition}[Borel 与 Lebesgue $\sigma$-代数]\label{def:2-1-3-2}
$\Borel(\R^n)$ 是开集生成的 $\sigma$-代数。下文将先证明
$\Borel(\R^n)\subset\mathcal L_C$；把 $m_C$ 限制到 $\Borel(\R^n)$ 得 Borel Lebesgue 测度
$m_B$。$m_B$ 的完成化记为 $\mathcal L$，称为 Lebesgue $\sigma$-代数，其测度记作 $m$。
\Cref{lem:lebesgue-completion-characterization} 将证明 $\mathcal L_C=\mathcal L$ 且
$m_C=m$；在那以前两种构造的记号保持分开。
\end{definition}

\begin{definition}[正则性]\label{def:2-1-3-3}
拓扑空间上的测度 $\mu$ 称外正则，如果
\[
 \mu(E)=\inf\{\mu(U):E\subset U,\ U\text{ 开}\};
\]
称在 $E$ 上内正则，如果
\[
 \mu(E)=\sup\{\mu(K):K\subset E,\ K\text{ 紧}\}.
\]
\end{definition}

\begin{lemma}[盒的外测度]\label{lem:lebesgue-box-lower-bound}
每个半开盒 $Q$ 满足 $m^*(Q)=|Q|$。
\end{lemma}

\begin{proof}
单盒覆盖给 $m^*(Q)\le|Q|$。若某条边长为零，则 $Q=\varnothing$（按半开盒约定），等式成立。
以下设各边长 $d_i=b_i-a_i>0$。取
\[
 K_\eta=\prod_{i=1}^n[a_i+\eta,b_i-\eta]
 \subset Q,
 \qquad0<\eta<\tfrac12\min_i d_i.
\]
给定可数半开盒
覆盖 $(Q_j)$；若 $\sum_j|Q_j|=\infty$，所需下界平凡。以下设总成本有限。给定 $\varepsilon>0$，
把 $Q_j$ 向外略微放大为开盒 $U_j$，使 $|U_j|<|Q_j|+\varepsilon2^{-j}$。紧集 $K_\eta$ 有有限子覆盖
$U_{j_1},\dots,U_{j_r}$。把 $K_\eta$ 与这些有限开盒在各坐标方向的所有端点组成共同网格。每个正体积
网格胞腔 $C$ 若包含于 $K_\eta$，取 $z\in C^\circ$。某个 $U_{j_l}$ 包含 $z$；因为网格包含
$U_{j_l}$ 的全部端点，所以 $C^\circ\subset U_{j_l}$。把每个胞腔分配给一个这样的 $l$。
胞腔边界位于有限个坐标超平面中；用法向总厚度小于 $\tau$ 的有限薄盒覆盖这些超平面后，有限盒
体积的可加性给
\[
 |K_\eta|\le\sum_{l=1}^r|U_{j_l}|+\tau
 \le\sum_j|Q_j|+\varepsilon+\tau.
\]
先令 $\tau\downarrow0$、再令
$\varepsilon\downarrow0$，得到
\[
 |K_\eta|=\prod_{i=1}^n(d_i-2\eta)\le\sum_j|Q_j|.
\]
最后令 $\eta\downarrow0$，得到
$|Q|\le\sum_j|Q_j|$。对所有覆盖
取下确界即得结论。
\end{proof}

\begin{lemma}[任意集合的开外逼近]\label{lem:lebesgue-outer-open-approximation}
对任意 $A\subset\R^n$ 与 $\varepsilon>0$，存在开集 $U\supset A$ 使
\[
 m^*(U)\le m^*(A)+\varepsilon.
\]
当 $m^*(A)<\infty$ 时还可取严格不等号。
\end{lemma}

\begin{proof}
若 $m^*(A)=\infty$，取 $U=\R^n$ 即可。以下设 $m^*(A)<\infty$。取半开盒覆盖
$A\subset\bigcup_iQ_i$，使
\[
 \sum_i|Q_i|<m^*(A)+\frac{\varepsilon}{2}.
\]
对每个 $i$ 选开盒 $O_i$ 与半开盒 $R_i$，使
\[
 Q_i\subset O_i\subset R_i,
 \qquad |R_i|<|Q_i|+\varepsilon2^{-i-1}.
\]
令 $U=\bigcup_iO_i$。则 $U$ 开、$A\subset U\subset\bigcup_iR_i$，从而
\[
 m^*(U)\le\sum_i|R_i|
 <\sum_i|Q_i|+\frac{\varepsilon}{2}
 <m^*(A)+\varepsilon.
\]
\end{proof}

\begin{proposition}[外测度可由任意细立方体计算]\label{proposition:2-1-3-1a}
给定 $\delta>0$，$m^*(E)$ 等于所有边长不超过 $\delta$ 的可数半开立方体覆盖的总体积下确界。
\end{proposition}

\begin{proof}
立方体是盒，所以允许所有盒所得下确界不大于所述下确界。若 $m^*(E)=\infty$，反向不等式平凡。
以下设 $m^*(E)<\infty$，并取总成本小于 $m^*(E)+\varepsilon/2$ 的盒覆盖 $(R_k)$。对每个 $k$，
选 $0<h_k\le\delta$ 足够小，并用网格 $h_k\Z^n$ 中所有与 $R_k$ 相交的半开立方体覆盖 $R_k$。
若 $R_k$ 的边长为 $d_{k,1},\ldots,d_{k,n}$，这些立方体的总体积至多
\[
 \prod_{j=1}^n(d_{k,j}+2h_k),
\]
它在 $h_k\downarrow0$ 时趋于 $|R_k|$。故可令该总量小于
$|R_k|+\varepsilon2^{-k-1}$。合并所有覆盖得到总成本小于 $m^*(E)+\varepsilon$，再令
$\varepsilon\downarrow0$。
\end{proof}

\begin{remark}
薄盒 $[0,1]\times[0,\varepsilon]$ 面积趋零而直径趋近 $1$。因此以后用 Lipschitz 映射控制像的
直径时，不能直接传送任意盒覆盖；必须先使用上一命题把覆盖立方体化。
\end{remark}

\begin{lemma}[Lebesgue 外测度的度量可加性]\label{lem:lebesgue-metric-outer}
若 $\dist(A,B)>0$，则
\[
 m^*(A\cup B)=m^*(A)+m^*(B).
\]
因此 $m^*$ 是度量外测度，所有 Borel 集都属于 $\mathcal L_C$；特别地，每个半开盒
\Caratheodory{} 可测。
\end{lemma}

\begin{proof}
次可加性给“$\le$”。设 $\delta=\dist(A,B)>0$，取边长小于 $\delta/\sqrt n$ 的固定网格。
给定 $A\cup B$ 的可数半开盒覆盖，把每个覆盖盒沿该网格切成至多可数个不交半开小盒；小盒总体积
等于原盒体积，每个小盒的直径小于 $\delta$，故不可能同时遇到 $A$ 与 $B$。收集所有遇到 $A$ 的
小盒得到 $A$ 的覆盖，收集所有遇到 $B$ 的小盒得到 $B$ 的覆盖；两族的总体积之和不超过原覆盖的
总成本。因此
\[
 m^*(A)+m^*(B)\le\sum_k|Q_k|
\]
对 $A\cup B$ 的每个盒覆盖 $(Q_k)$ 成立。取下确界得到“$\ge$”。于是 $m^*$ 是 metric outer
measure，\cref{proposition:2-1-2-3} 给出 Borel 可测性。半开盒是 Borel 集，最后一句随之成立。
\end{proof}

\begin{theorem}[Lebesgue 测度的基本几何性质]\label{theorem:2-1-3-1}
Lebesgue 测度在 $\R^n$ 上平移不变；对 $r>0$ 和 Lebesgue 可测 $E$，
\[
 m(rE)=r^n m(E).
\]
每个半开盒 $Q$ 满足 $m(Q)=|Q|$。
\end{theorem}

\begin{proof}
平移一个盒不改变成本，所以逐项平移覆盖给 $m^*(E+x)=m^*(E)$；用 $-x$ 得反向不等式。
对缩放，任一盒覆盖 $(Q_k)$ 被送到盒覆盖 $(rQ_k)$，其成本为
$\sum_k|rQ_k|=r^n\sum_k|Q_k|$；对逆缩放再做一次，故
$m^*(rE)=r^nm^*(E)$。这些等式表明平移和缩放保持 \Caratheodory{} 等式。它们还是
Euclidean 同胚，因而保持 Borel 集；外测度公式又说明它们把 Borel 零集送到 Borel 零集，所以也
保持完成化 $\mathcal L$。盒公式由\cref{lem:lebesgue-box-lower-bound,lem:lebesgue-metric-outer}
得到。
\end{proof}

\begin{theorem}[正则性与 $\sigma$-有限性]\label{theorem:2-1-3-2}
$\R^n=\bigcup_{k\ge1}[-k,k]^n$，每个耗尽盒有限测。Borel 集外正则，开集内正则；Lebesgue
完成化对每个可测集外正则，并对每个有限测度可测集内正则。紧集具有有限测度。
\end{theorem}

\begin{proof}
若 $U$ 开，对每个 $x\in U$ 可取有理端点闭盒 $K\subset U$ 且 $x\in\interior K$。
这样的盒只有可数多个，故 $U=\bigcup_jK_j$。令 $L_N=\bigcup_{j\le N}K_j$；它紧且
$L_N\uparrow U$。测度的可数可加性蕴含从下连续，故
$m(U)=\sup_Nm(L_N)$，得到开集内正则。

设 $E$ 为 Borel 集。若 $m(E)=\infty$，则每个包含 $E$ 的开集也有无穷测度，外正则等式平凡。
以下设 $m(E)<\infty$。由引理 \ref{lem:lebesgue-outer-open-approximation}，对每个 $j$ 可取
开集 $U_j\supset E$，使
\[
 m(U_j)=m^*(U_j)<m^*(E)+2^{-j}=m(E)+2^{-j}.
\]
令 $j\to\infty$，得到 Borel 集外正则。

若 $E$ 是有限测度 Borel 集，由 $[-R,R]^n\uparrow\R^n$ 的从下连续性，先取 $R$ 使
$m(E\setminus[-R,R]^n)<\varepsilon/2$。对 Borel 集 $[-R,R]^n\setminus E$ 作外逼近，取开集
$V$ 包含它并满足
\[
 m\bigl(V\setminus([-R,R]^n\setminus E)\bigr)<\varepsilon/2.
\]
令 $F=[-R,R]^n\setminus V$。则 $F$ 是包含于 $E$ 的紧集，而且
\[
 E\setminus F
 \subset (E\setminus[-R,R]^n)
 \cup\bigl(V\setminus([-R,R]^n\setminus E)\bigr),
\]
故 $m(E\setminus F)<\varepsilon$。这给有限 Borel 集内正则。

完成化可测集 $E$ 有 Borel 代表 $B$ 且 $E\mathbin\triangle B$ 零测。若 $m(E)=\infty$，外正则
等式平凡。若 $m(E)<\infty$，取 Borel 零集 $Z\supset E\mathbin\triangle B$。给定
$\varepsilon>0$，由上一引理取开集 $O\supset Z$ 使 $m(O)<\varepsilon/2$，并由 Borel 外正则性
取开集 $V\supset B$ 使 $m(V)<m(B)+\varepsilon/2$。则 $E\subset V\cup O$ 且
\[
 m(V\cup O)\le m(V)+m(O)<m(E)+\varepsilon,
\]
故 $E$ 外正则。继续使用上述 Borel 零集 $Z$。先取紧集 $K\subset B$，使
$m(B\setminus K)<\varepsilon/2$；
再取开集 $O\supset Z$，使 $m(O)<\varepsilon/2$。于是 $K\setminus O$ 是包含于 $E$ 的紧集，并且
\[
 m\bigl(E\setminus(K\setminus O)\bigr)
 \le m(E\setminus B)+m(B\setminus K)+m(K\cap O)<\varepsilon.
\]
这证明完成化的有限测度集合内正则。对于无限测度集合，内逼近须先在有限盒中局部化；不能把两个无穷
测度作差来制造误差估计。
\end{proof}

\begin{proposition}[Borel 代表]\label{proposition:2-1-3-3}
每个 Lebesgue 可测集 $E$ 存在 $G_\delta$ 集 $G\supset E$ 和 $F_\sigma$ 集 $F\subset E$，使
$m(G\setminus E)=m(E\setminus F)=0$。
\end{proposition}

\begin{proof}
令 $Q_k=[-k,k]^n$，并置 $D_1=Q_1$、$D_k=Q_k\setminus Q_{k-1}$（$k\ge2$）。这些 Borel 层
两两不交、覆盖 $\R^n$，且每层有限测。对每个 $j,k\ge1$，由外正则性取开集
$U_{j,k}\supset E\cap D_k$，使
\[
 m\bigl(U_{j,k}\setminus(E\cap D_k)\bigr)<2^{-j-k}.
\]
令 $U_j=\bigcup_kU_{j,k}$、$G=\bigcap_jU_j$。则 $G$ 是 $G_\delta$ 集且包含 $E$。又因为
\[
 U_j\setminus E
 \subset\bigcup_k\bigl(U_{j,k}\setminus(E\cap D_k)\bigr),
\]
所以 $m(U_j\setminus E)\le2^{-j}$，进而 $m(G\setminus E)=0$。

再由有限测度集合的内正则性，对每个 $k$ 取紧集 $K_k\subset E\cap Q_k$，使
$m((E\cap Q_k)\setminus K_k)<2^{-k}$，并令 $F=\bigcup_kK_k$。这是包含于 $E$ 的
$F_\sigma$ 集。固定 $k$ 后，对每个 $\ell\ge k$ 都有
\[
 (E\cap Q_k)\setminus F
 \subset (E\cap Q_\ell)\setminus K_\ell,
\]
故其测度不超过 $2^{-\ell}$；令 $\ell\to\infty$，得到
$m((E\cap Q_k)\setminus F)=0$。最后对 $k$ 取可数并，便有 $m(E\setminus F)=0$。
\end{proof}

\begin{proposition}[两种 Lebesgue 可测类相同]\label{lem:lebesgue-completion-characterization}
$\mathcal L_C=\mathcal L$，并且 $m_C=m$。
\end{proposition}

\begin{proof}
完成化包含于 $\mathcal L_C$，因为该类含 Borel 集且 $m_C$ 完备。反过来，设
$E\in\mathcal L_C$，把空间分为不交有限测度盒层 $D_k$。对每个 $j,k$，先取开集
$U_{j,k}\supset E\cap D_k$；由引理 \ref{lem:lebesgue-outer-open-approximation} 可使
$m^*(U_{j,k})<m^*(E\cap D_k)+2^{-j-k}$。集合 $E\cap D_k$ 是 \Caratheodory{} 可测的且
外测度有限；以它切割可测开集 $U_{j,k}$ 得
\[
 m^*(U_{j,k})=m^*(E\cap D_k)+m^*(U_{j,k}\setminus(E\cap D_k)),
\]
所以差集外测度小于 $2^{-j-k}$。令
$U_j=\bigcup_kU_{j,k}$ 和 $G=\bigcap_jU_j$。则 $G$ 为 Borel 集，$E\subset G$，并且
$m^*(G\setminus E)\le2^{-j}$ 对每个 $j$ 成立，故差集外测度为零。令 $N=G\setminus E$。
对每个 $r\ge1$，由 $m^*(N)=0$ 取开集 $O_r\supset N$ 且 $m(O_r)<2^{-r}$，再置
\[
 Z=\bigcap_{r=1}^{\infty}O_r.
\]
则 $Z$ 是包含 $N$ 的 Borel 集，并且对每个 $r$ 都有
$0\le m(Z)\le m(O_r)<2^{-r}$，故 $m(Z)=0$。因此 $E$ 属于
Borel 完成化，且两种限制测度都等于 $m^*(E)$。
\end{proof}

\begin{example}[可数集零测度]\label{ex:2-1-3-1}
前面的逐点覆盖证明给出每个可数集零测。特别地，Dirichlet 函数取值 $1$ 的集合虽在每个区间
稠密，却在 Lebesgue 意义下可忽略。
\end{example}

\begin{example}[Cantor 集]\label{ex:2-1-3-2}
三分 Cantor 集 $C$ 在第 $k$ 阶包含于 $2^k$ 个长度 $3^{-k}$ 的闭区间并中，故
$m(C)\le2^k3^{-k}=(2/3)^k$。令 $k\to\infty$ 得 $m(C)=0$。它却不可数、紧且无内点，说明基数、
拓扑大小和测度大小彼此不同。
\end{example}

\begin{example}[完成化确实增加集合]\label{ex:2-1-3-3}
Euclidean 空间的 Borel 集由一个可数基生成。把“取补集”和“取可数并”按可数序数逐层编码，每个
Borel 集可由一个可数码描述；所有可数码的基数至多为
$|\N^\N|=|\R|$，故 Borel 集总数至多为连续统。另一方面，映射
\[
 (\varepsilon_j)_{j\ge1}\longmapsto
 \sum_{j=1}^{\infty}2\varepsilon_j3^{-j},
 \qquad \varepsilon_j\in\{0,1\},
\]
把 $\{0,1\}^{\N}$ 双射到 Cantor 集 $C$（端点采用只含 $0,2$ 的三进制展开），所以
$|C|=|\R|$。Cantor 定理给 $|2^C|>|\R|$，故 $C$ 至少有一个非 Borel 子集 $A$。由于
$A\subset C$ 且
$m(C)=0$，完成化使 $A$ Lebesgue 可测且零测。
\end{example}

\begingroup
\raggedbottom
\begin{proposition}[Jordan 理论嵌入]\label{proposition:2-1-3-4}
有界集合 $E$ Jordan 可测，当且仅当 $m(\partial E)=0$；此时 $E$ Lebesgue 可测且
$J(E)=m(E)$。
\end{proposition}

\begin{proof}
若 $E$ Jordan 可测，\cref{proposition:2-1-1-2} 给出边界的有限任意小盒覆盖，故
$m(\partial E)=0$。反过来，$\partial E$ 是有界闭集，因而紧。给定 $\varepsilon>0$，先以总体积
小于 $\varepsilon/2$ 的可数半开盒 $(Q_j)$ 覆盖它，再把第 $j$ 个盒略微外扩为开盒 $O_j$，使
\[
 |O_j|<|Q_j|+\varepsilon2^{-j-2}.
\]
紧性选出有限子覆盖 $O_{j_1},\dots,O_{j_r}$。对每个选中开盒再取包含它的半开盒 $R_l$，并把
有限个增量分配成
\[
 \sum_{l=1}^r(|R_l|-|O_{j_l}|)<\frac{\varepsilon}{4}.
\]
于是
\[
 \sum_{l=1}^r|R_l|
 <\sum_{j=1}^{\infty}|Q_j|+\frac{\varepsilon}{4}+\frac{\varepsilon}{4}
 <\varepsilon.
\]
故 $J^*(\partial E)=0$，Jordan 边界判据给出
Jordan 可测性。

在任一方向，只要 $m(\partial E)=0$，就有
$E=\interior{E}\cup(E\setminus\interior{E})$，其中第一项为 Borel 集，第二项包含于零测闭集
$\partial E$；Lebesgue 测度的完备性因此给出 $E$ 可测。最后，给定 $\varepsilon>0$，取初等集
$B\subset E\subset A$ 使 $|A|-|B|<\varepsilon$。初等集的 Lebesgue 测度等于其体积，故
\[
 |B|=m(B)\le m(E)\le m(A)=|A|,
 \qquad
 |B|\le J(E)\le |A|.
\]
于是 $|m(E)-J(E)|\le |A|-|B|<\varepsilon$；令 $\varepsilon\downarrow0$ 即得
$J(E)=m(E)$。
\end{proof}

\begin{figure}[htbp]
\centering
\begin{tikzpicture}[scale=1]
 \draw[BookAccent,thick,rounded corners] (-3,-1.5) rectangle (3,1.5) node[above left] {$U_j$};
 \draw[BookInk,thick] plot[smooth cycle] coordinates {(-2,-.7)(-1.4,1)(-.2,.55)(.9,1.05)(2.1,.15)(1.5,-1)(.2,-.7)(-1.1,-1.1)};
 \node at (0,.05) {$E$};
 \draw[BookWarm!90,thick,rounded corners] (-1.2,-.55) rectangle (1.05,.55) node[below right] {$K_j$};
 \draw[<->,thin] (2.2,-1.45)--(2.2,-1.0) node[midway,right] {$\to0$};
 \draw[<->,thin] (-1.2,-.55)--(-1.2,-.95) node[midway,left] {$\to0$};
\end{tikzpicture}
\caption{Lebesgue 正则性：开集从外、紧集从内逼近。箭头标记的是测度误差，而非 Euclidean 距离。}
\label{fig:lebesgue-inner-outer-regularity}
\end{figure}

\begin{remark}
平移不变和单位盒正规化不能在 $2^{\R^n}$ 上产生可数可加体积；\Caratheodory{} 构造必须舍弃某些
子集。另一个常见混淆来自完成化：连续映射自动拉回 Borel 集，却不自动拉回任意完成化集合。
若要拉回完成化，需另证目标零测 Borel 集的逆像在源中零测，或者在拉回后重新完成。
\end{remark}

这一体积将成为第~3~章积分、换元与卷积的底层测度，并在 Fourier 分析中提供平移不变的积分背景。
正则性还会把集合逼近转化为 $C_c(\R^n)$ 的函数逼近；概率密度则是在这个测度之上描述分布的一种
方式。相应的积分、卷积与密度理论从下一章开始建立。

\begin{exercise}
证明中间三分 Cantor 集不可数、无内点并且零测。
\end{exercise}
\begin{exercise}
证明每个开集是可数个两两仅在边界相交的半开 dyadic 盒之并。
\end{exercise}
\begin{exercise}
从正则性重新证明每个 Lebesgue 可测集有 Borel 代表。
\end{exercise}
\begin{exercise}
先对任意 $A\subset\R^n$ 从覆盖定义证明
$m^*(A+x)=m^*(A)$ 与 $m^*(rA)=r^nm^*(A)$；再证明这些映射保持可测性，并对可测 $A$ 写成相应的
$m$ 等式。
\end{exercise}

下一节不再依赖 Euclidean 盒。我们将把“存在一个扩张”和“扩张由简单测试唯一确定”分开处理，
并把拓扑正则性、有符号抵消纳入同一语言。
\clearpage
\endgroup

\section{Borel 测度、Radon 正则性与有符号扩张}\label{sec:02-02}

\subsection{生成类、单调类与测度唯一性}\label{subsec:2-2-1}
\index{Dynkin 系统}\index{单调类}\index{预测度}

测度常常只在少数集合上可以直接计算：实线上的区间、Euclidean 空间中的矩形、坐标空间中的
柱集。假设两个有限测度在这样一族集合上相同。把“相等的集合”收集起来，看上去可以证明它们
组成 $\sigma$-代数；尝试补集时却立刻出现
\[
 \mu(A^c)=\mu(X)-\mu(A).
\]
这个式子要求先知道总质量相等且有限。处理不交并比处理任意并自然得多。与其强迫这个集合族
满足过多运算，不如准确记录证明真正保持的闭包。

\begin{definition}[半环、环与代数]\label{def:2-2-1-0a}
集合族 $\mathcal S\subset2^X$ 称为半环，如果 $\varnothing\in\mathcal S$，它对有限交封闭，
并且 $A\setminus B$ 可写成有限个两两不交的 $\mathcal S$ 中集合之并。对有限并和差封闭的集合族
称为环；含 $X$ 的环称为代数。半环中集合的有限不交并构成它生成的环。
\end{definition}

\begin{definition}[预测度]\label{def:2-2-1-0b}
环 $\mathcal R$ 上的函数 $\mu_0:\mathcal R\to[0,\infty]$ 称为预测度，如果
$\mu_0(\varnothing)=0$，并且每当 $A=\bigsqcup_{k\ge1}A_k$ 且
$A,A_k\in\mathcal R$ 时，
\[
 \mu_0(A)=\sum_{k=1}^{\infty}\mu_0(A_k).
\]
\end{definition}

\begin{lemma}[半环数据延拓到环]\label{lem:premeasure-semiring-to-ring}
若 $\mu_0$ 在半环 $\mathcal S$ 上满足上述不交可数分解条件，则
\[
 \widetilde\mu_0\Bigl(\bigsqcup_{j=1}^rS_j\Bigr)=\sum_{j=1}^r\mu_0(S_j)
 \label{eq:semiring-ring-extension}
\]
在 $\mathcal S$ 生成的环上良定义，并且是预测度。
\end{lemma}

\begin{proof}
若同一集合还有不交表示 $\bigsqcup_{k=1}^tT_k$，则
$S_j=\bigsqcup_k(S_j\cap T_k)$ 和 $T_k=\bigsqcup_j(S_j\cap T_k)$，其中所有交仍在半环中。
半环上的有限可加性给
\[
 \sum_j\mu_0(S_j)=\sum_{j,k}\mu_0(S_j\cap T_k)=\sum_k\mu_0(T_k),
\]
故 \eqref{eq:semiring-ring-extension} 良定义。

若环中 $A=\bigsqcup_{n\ge1}A_n$，写成不交半环分解
\[
 A=\bigsqcup_{j=1}^rS_j,
 \qquad A_n=\bigsqcup_{k=1}^{r_n}T_{n,k}.
\]
因为 $A_n\subset A$ 且各 $A_n$ 不交，对固定的 $n,k$ 有
$T_{n,k}=\bigsqcup_{j=1}^r(T_{n,k}\cap S_j)$；对固定的 $j$ 又有可数不交分解
$S_j=\bigsqcup_{n,k}(S_j\cap T_{n,k})$。于是半环上的分解条件和非负级数的换序给
\[
 \sum_{n=1}^{\infty}\widetilde\mu_0(A_n)
 =\sum_{n,k,j}\mu_0(T_{n,k}\cap S_j)
 =\sum_{j=1}^r\mu_0(S_j)
 =\widetilde\mu_0(A).
\]
所以 $\widetilde\mu_0$ 是预测度。
\end{proof}

\begin{example}[半开盒的扩张数据]\label{ex:2-2-1-0}
$\R^n$ 中的半开盒连同空集构成半环：两个盒的交仍是盒，而盒之差可沿双方端点切成有限个不交
盒。盒体积满足不交分解可加性。若 $Q=\bigsqcup_{k\ge1}Q_k$，取紧内缩盒 $K_\eta\subset Q$。
把每个 $Q_k$ 外扩为开盒 $U_k$，并分配误差使
\[
 |U_k|<|Q_k|+\varepsilon2^{-k}.
\]
紧性给有限子覆盖 $K_\eta\subset\bigcup_{k\in F}U_k$；共同端点网格比较给
\[
 |K_\eta|\le\sum_{k\in F}|U_k|
 \le\sum_{k=1}^{\infty}|Q_k|+\varepsilon.
\]
依次令 $\varepsilon\downarrow0$、$\eta\downarrow0$，得到
$|Q|\le\sum_k|Q_k|$。反向地，对每个 $N$，有限共同网格给
\[
 \sum_{k=1}^N|Q_k|
 =\left|\bigsqcup_{k=1}^NQ_k\right|\le|Q|;
\]
令 $N\to\infty$ 得 $\sum_k|Q_k|\le|Q|$。因此盒体积先延拓到有限不交盒环，再由下文扩张为
Lebesgue 测度。乘积测度会重复同一结构。
\end{example}

\begin{definition}[$\pi$-系统与 $\lambda$-系统]\label{def:2-2-1-1}
集合族 $\mathcal P$ 对有限交封闭时称为 $\pi$-系统。集合族 $\mathcal D$ 称为 $\lambda$-系统
或 Dynkin 系统，如果
\begin{enumerate}
 \item $X\in\mathcal D$；
 \item $A\subset B$ 且 $A,B\in\mathcal D$ 时，$B\setminus A\in\mathcal D$；
 \item 两两不交的 $D_k\in\mathcal D$ 满足 $\bigsqcup_kD_k\in\mathcal D$。
\end{enumerate}
\end{definition}

\begin{definition}[单调类]\label{def:2-2-1-2}
集合族若对递增并和递减交封闭，称为集合单调类。一个有界实值函数类若含常数，并且对具有共同
一致界的点态单调函数列之极限封闭，称为函数单调类。
\end{definition}

\begin{definition}[简单函数]\label{def:2-2-1-simple-function}
可测空间 $(X,\mathcal A)$ 上只取有限多个值的可测函数称为简单函数。非负简单函数总可按其非零
水平集写成
\[
 s=\sum_{j=1}^r a_j\1_{A_j},
 \qquad a_j>0,\quad A_j\in\mathcal A,\quad A_j\cap A_k=\varnothing\ (j\ne k).
\]
反过来，任何这样的有限和都是非负可测简单函数。表示未必唯一；需要从表示定义数值时，必须证明
共同有限分割不会改变结果。
\end{definition}

\begin{definition}[生成 $\sigma$-代数]\label{def:2-2-1-3}
$\sigma(\mathcal C)$ 是所有包含 $\mathcal C$ 的 $\sigma$-代数之交。它是包含
$\mathcal C$ 的最小 $\sigma$-代数，而不只是能由某个有限集合公式写出的集合族。
\end{definition}

\begin{theorem}[Dynkin 的 $\pi$--$\lambda$ 定理]\label{theorem:2-2-1-1}
若 $\mathcal P$ 是 $\pi$-系统，则包含 $\mathcal P$ 的最小 $\lambda$-系统等于
$\sigma(\mathcal P)$。
\end{theorem}

\begin{proof}
令 $\mathcal D$ 为包含 $\mathcal P$ 的最小 $\lambda$-系统。只需证明 $\mathcal D$ 对有限交封闭，
因为此时补集和任意可数并均可由 $\lambda$-运算得到。

对 $A\subset X$ 令
$\mathcal D_A=\{B\in\mathcal D:A\cap B\in\mathcal D\}$。固定
$A\in\mathcal P$。由于 $\mathcal P$ 对交封闭，$\mathcal P\subset\mathcal D_A$。而且
$A\cap X=A\in\mathcal D$；若 $B\subset C$ 且 $B,C\in\mathcal D_A$，则
\[
 A\cap(C\setminus B)=(A\cap C)\setminus(A\cap B)\in\mathcal D;
\]
若 $(B_k)$ 在 $\mathcal D_A$ 中两两不交，则
$A\cap\bigsqcup_kB_k=\bigsqcup_k(A\cap B_k)\in\mathcal D$。因此
$\mathcal D_A$ 是 $\lambda$-系统，故最小性给
$\mathcal D_A=\mathcal D$。因此 $A\cap B\in\mathcal D$ 对所有
$A\in\mathcal P,B\in\mathcal D$ 成立。

现在固定 $B\in\mathcal D$，令
$\mathcal E_B=\{A\in\mathcal D:A\cap B\in\mathcal D\}$。上一段给
$\mathcal P\subset\mathcal E_B$。用恒等式
$(C\setminus A)\cap B=(C\cap B)\setminus(A\cap B)$ 与
$(\bigsqcup_kA_k)\cap B=\bigsqcup_k(A_k\cap B)$，可以逐条验证
$\mathcal E_B$ 为 $\lambda$-系统，所以
$\mathcal E_B=\mathcal D$。于是 $\mathcal D$ 对有限交封闭，证明完成。
\end{proof}

\begin{theorem}[测度唯一性]\label{theorem:2-2-1-2}
设有限测度 $\mu,\nu$ 在 $\pi$-系统 $\mathcal P$ 上相等，并且
$\mu(X)=\nu(X)<\infty$，则二者在 $\sigma(\mathcal P)$ 上相等。

若存在共同耗尽 $P_n\in\mathcal P$，满足 $P_n\uparrow X$ 且
$\mu(P_n)=\nu(P_n)<\infty$，则同一结论对 $\sigma$-有限测度成立。
\end{theorem}

\begin{proof}
有限情形令 $\mathcal D=\{A:\mu(A)=\nu(A)\}$。总质量相等使 $X\in\mathcal D$；若
$A\subset B$ 且二者在 $\mathcal D$ 中，有限性给
\[
 \mu(B\setminus A)=\mu(B)-\mu(A)=\nu(B)-\nu(A)=\nu(B\setminus A).
\]
对不交可数并则用两测度的可数可加性。因此 $\mathcal D$ 是含 $\mathcal P$ 的 $\lambda$-系统，
由 $\pi$--$\lambda$ 定理包含 $\sigma(\mathcal P)$。

$\sigma$-有限情形固定 $n$，比较有限测度
$A\mapsto\mu(A\cap P_n)$ 与 $A\mapsto\nu(A\cap P_n)$。它们在
$\mathcal P$ 上相等，因为 $A\cap P_n\in\mathcal P$，故在生成 $\sigma$-代数上相等。再令
$n\to\infty$ 并用从下连续性。共同耗尽必须位于测试类中，正是为了保证交后的等式仍在已知范围内。
\end{proof}

\begin{theorem}[\Caratheodory{} 扩张：环版]\label{theorem:2-2-1-0}
设 $\mu_0$ 是环 $\mathcal R$ 上的预测度，并令
$X=\bigcup_{R\in\mathcal R}R$。以可数个 $\mathcal R$ 中集合覆盖并取 $\mu_0$ 成本下确界得到
$\mu^*$。则每个 $R\in\mathcal R$ 都是 $\mu^*$-可测，且
$\mu^*|_{\sigma(\mathcal R)}$ 扩张 $\mu_0$。若存在
$R_n\uparrow X$、$R_n\in\mathcal R$ 且 $\mu_0(R_n)<\infty$，扩张在
$\sigma(\mathcal R)$ 上唯一。
\end{theorem}

\begin{proof}
覆盖下确界是外测度。固定 $A\in\mathcal R$ 和测试集 $E\subset X$。若
$E\subset\bigcup_kR_k$，则
\[
 E\cap A\subset\bigcup_k(R_k\cap A),\qquad
 E\setminus A\subset\bigcup_k(R_k\setminus A).
\]
环对交和差封闭，且预测度的有限可加性给
$\mu_0(R_k)=\mu_0(R_k\cap A)+\mu_0(R_k\setminus A)$。对覆盖成本求和、再对所有覆盖取下确界，
得到
\[
 \mu^*(E)\ge\mu^*(E\cap A)+\mu^*(E\setminus A).
\]
与次可加合并，$A$ 可测。这里没有使用 $A^c\in\mathcal R$。

单集合覆盖给 $\mu^*(A)\le\mu_0(A)$。反过来，若 $A\subset\bigcup_kR_k$，把
$R_k$ 显式不交化为
\[
 D_1=R_1,\qquad D_k=R_k\setminus\bigcup_{j<k}R_j\quad(k\ge2).
\]
环对有限并与差封闭，故 $D_k,A\cap D_k\in\mathcal R$；这些集合两两不交，而且
$A=\bigsqcup_k(A\cap D_k)$。预测度的可数可加性与由有限可加性推出的单调性给
\[
 \mu_0(A)=\sum_k\mu_0(A\cap D_k)
 \le\sum_k\mu_0(D_k)
 \le\sum_k\mu_0(R_k).
\]
对所有覆盖取下确界得反向不等式。所以限制确实扩张 $\mu_0$。
存在性由 \Caratheodory{} 定理完成；唯一性在共同有限耗尽上应用上一测度唯一性定理。
\end{proof}

\begin{theorem}[集合与函数单调类定理]\label{theorem:2-2-1-3}
包含代数 $\mathcal A$ 的最小集合单调类是 $\sigma(\mathcal A)$。若有界实函数的线性类
$\mathcal H$ 含常数，对有共同一致界的点态单调极限封闭，并且
$\1_A\in\mathcal H$ 对每个 $A\in\mathcal A$ 成立，则 $\mathcal H$ 包含所有有界
$\sigma(\mathcal A)$-可测实函数。复值版分别处理实部、虚部。
\end{theorem}

\begin{proof}
令 $\mathcal M$ 为含 $\mathcal A$ 的最小单调类。固定 $B\in\mathcal A$，考察满足
$A\cap B,A\setminus B,B\setminus A\in\mathcal M$ 的 $A\in\mathcal M$。对递增列
$A_n\uparrow A$，三个相应序列分别为
\[
 A_n\cap B\uparrow A\cap B,\qquad
 A_n\setminus B\uparrow A\setminus B,\qquad
 B\setminus A_n\downarrow B\setminus A;
\]
对递减列，三个单调方向全部反转。因此该类是含 $\mathcal A$ 的单调类，必等于
$\mathcal M$。这说明：对每个 $B\in\mathcal A$ 和 $A\in\mathcal M$，上述三个集合都属于
$\mathcal M$。

现在固定 $B\in\mathcal M$，并对 $A$ 重复同一个单调类论证。上一段保证这个新类包含
$\mathcal A$，故它也等于 $\mathcal M$。从而 $\mathcal M$ 对任意两个成员的交与差封闭。特别地，
$X\setminus A\in\mathcal M$，有限并也属于 $\mathcal M$；任意可数并是其有限部分并的递增极限。所以
$\mathcal M$ 是 $\sigma$-代数，即 $\mathcal M=\sigma(\mathcal A)$。

对函数版，令 $\mathcal G=\{A:\1_A\in\mathcal H\}$。常数给
$X\in\mathcal G$；线性和有界单调极限使 $\mathcal G$ 成为含 $\mathcal A$ 的单调类。因此
$\1_A\in\mathcal H$ 对所有 $A\in\sigma(\mathcal A)$ 成立。线性性给所有有界简单函数。
若 $0\le f\le M$ 可测，dyadic 下逼近
\[
 s_n=2^{-n}\lfloor2^nf\rfloor
\]
是简单函数、递增至 $f$ 且共同由 $M$ 控制，故 $f\in\mathcal H$。一般有界实函数拆成正负部。
\end{proof}

\begin{proposition}[分布由半直线确定]\label{proposition:2-2-1-4}
有限 Borel 测度 $\mu$ 在 $\R$ 上由
$F(x)=\mu(( -\infty,x])$ 唯一决定。
\end{proposition}

\begin{proof}
集合族 $\{(-\infty,q]:q\in\Q\}$ 对有限交封闭。对任意实数 $a$，
\[
 (-\infty,a)=\bigcup_{\substack{q\in\Q\\q<a}}(-\infty,q],\qquad
 (-\infty,a]=\bigcap_{\substack{q\in\Q\\q>a}}(-\infty,q].
\]
因此每个开区间可由两个这样的半直线作差得到，该族生成 $\Borel(\R)$。两个具有同一 $F$ 的测度在
这个 $\pi$-系统上相等；测度的从下连续性还给出
\[
 \mu(\R)=\lim_{n\to\infty}F(n)=\nu(\R).
\]
总质量相等且有限，故测度唯一性定理可以应用并给出结论。
\end{proof}

\begin{example}[矩形与柱集]\label{ex:2-2-1-2}
可测矩形 $A\times B$ 对有限交封闭并生成乘积 $\sigma$-代数。对坐标空间
$\prod_{i\in I}X_i$，只限制有限个坐标的柱集也是 $\pi$-系统；这些有限观测生成柱
$\sigma$-代数。确切地，若
\[
 C=\bigcap_{i\in F}\pi_i^{-1}(A_i),\qquad
 D=\bigcap_{i\in G}\pi_i^{-1}(B_i),
\]
其中 \(F,G\subset I\) 有限，则 \(C\cap D\) 仍只限制 \(F\cup G\) 中的有限多个坐标。后续乘积
测度与过程的有限维分布都由此被唯一识别。
\end{example}

\begin{example}[相等类恰好需要 Dynkin 闭包]\label{ex:2-2-1-3}
令 \(X=\{1,2,3,4\}\)，取均匀概率 \(\nu\)，并令
\[
 \mu(\{1\})=\mu(\{4\})=\frac38,\qquad
 \mu(\{2\})=\mu(\{3\})=\frac18.
\]
对 \(A=\{1,2\}\) 和 \(B=\{1,3\}\)，有
\[
 \mu(A)=\nu(A)=\mu(B)=\nu(B)=\frac12,
\]
但
\[
 \mu(A\cap B)=\mu(\{1\})=\frac38
 \ne\frac14=\nu(\{1\}).
\]
因此相等类 \(\mathcal D=\{E:\mu(E)=\nu(E)\}\) 虽是 \(\lambda\)-系统，却不必对交封闭；唯一性
证明必须从一个本身对交封闭的 \(\pi\)-系统起步。
\end{example}

\begin{example}[区间生成 Borel]\label{ex:2-2-1-1}
上一命题的证明同时说明，有理端点半直线已经生成 $\Borel(\R)$。不可数的所有开集观测因此可由
一个可数测试族编码。
\end{example}

\begin{figure}[htbp]
\centering
\begin{tikzpicture}[node distance=8mm and 12mm]
 \node[book strong node] (p) {$\pi$-系统\\有限交};
 \node[book node,right=of p] (d) {最小 $\lambda$-系统\\补差、不交可数并};
 \node[book strong node,right=of d] (s) {$\sigma(\mathcal P)$};
 \node[book warm node,below=of d] (m) {函数单调类\\简单函数 $\to$ 有界可测函数};
 \draw[book arrow] (p)--(d);
 \draw[book arrow] (d)--node[above]{双重闭包}(s);
 \draw[book dashed arrow] (p)--(m);
 \draw[book arrow] (m)--(s);
\end{tikzpicture}
\caption{集合的 $\pi$--$\lambda$ 闭包与函数单调闭包。两条路线都把简单观测扩展到生成对象。}
\label{fig:pi-lambda-monotone-closures}
\end{figure}

\begin{remark}
$\sigma$-有限唯一性不是把“有限”两个字删除。必须给出共同有限耗尽，而且耗尽集合要落在测试类中。
函数单调类也不是对任意点态极限封闭：共同一致界保证极限仍属于有界函数的论域。
\end{remark}

这些闭包机制在后文有三类应用：独立性可先在生成 $\pi$-系统上检验；乘积测度与随机过程的有限维
分布可先在矩形或柱集上识别；第~3~章定义积分后，函数单调类将把简单函数等式推进到有界可测函数。

\begin{exercise}
证明有理端点半开区间生成 $\Borel(\R)$。
\end{exercise}
\begin{exercise}
证明两个概率测度若在一个生成 $\pi$-系统上相等，则处处相等。
\end{exercise}
\begin{exercise}
在三点集合上构造一个 $\lambda$-系统，它不是 $\sigma$-代数。
\end{exercise}
\begin{exercise}
从半开区间半环构造有限不交并环，并证明长度函数从上连续于空集。
\end{exercise}
\begin{exercise}
逐层复现函数单调类证明，从指标函数推进到一般有界可测函数。
\end{exercise}

唯一性回答了“简单测试能否识别测度”；接下来要问的是，已经识别出的 Borel 测度能否由开集、紧集与
紧支撑连续函数可靠地逼近。

\subsection{Radon 测度、支撑与一致紧性}\label{subsec:2-2-2}
\index{Radon 测度}\index{支撑}\index{一致紧性}

Borel $\sigma$-代数告诉我们哪些集合可以测，却没有保证一个 Borel 集能由几何上可见的开集和紧集
逼近。分析中常用的测试对象有紧支撑，概率极限也必须阻止质量逃向无穷；两件事都要求测度与拓扑
之间有更强的配合。

\begin{definition}[Dirac 测度]\label{def:dirac-measure}
设 $(X,\mathcal A)$ 为可测空间且 $x\in X$。定义
\[
 \delta_x(A)=\1_A(x)=
 \begin{cases}1,&x\in A,\\0,&x\notin A.
 \end{cases}
\]
称为点 $x$ 处的 Dirac 测度。这个定义不要求 $\{x\}\in\mathcal A$；只有当我们把单点本身作为
事件或讨论原子集合时，才需要单点可测。
\end{definition}

\begin{definition}[原子]\label{def:measure-atom}
设 $\mu$ 为测度。若可测集 $A$ 满足 $\mu(A)>0$，并且每个可测子集 $B\subset A$ 都满足
$\mu(B)=0$ 或 $\mu(A\setminus B)=0$，则称 $A$ 为 $\mu$ 的原子。若单点可测且
$\mu(\{x\})>0$，则 $\{x\}$ 是原子，并简称 $x$ 为点原子。Dirac 测度 $\delta_x$ 的全部质量
集中在点原子 $x$ 上。
\end{definition}

\begin{example}[质量逃向无穷]\label{ex:2-2-2-3}
在 $\R$ 上取 $\delta_n$。任何紧集 $K$ 有界，故当 $n$ 足够大时
$\delta_n(\R\setminus K)=1$。所以没有一个紧集能统一包含这些概率的大部分质量。另一方面，若
$f$ 连续且在某个紧集外为零，则 $f(n)=0$ 最终成立；于是只观察这类函数时，$\delta_n$ 看起来趋于
零，总质量却逸向无穷。这个差别将在模糊收敛（vague convergence）与概率弱收敛之间再次出现。
\end{example}

\begin{definition}[Radon 测度]\label{def:2-2-2-1}
测度称为局部有限，如果每点有有限测度的开邻域。局部紧 Hausdorff 空间 $X$ 上的 Borel 测度
$\mu$ 称为 Radon 测度，如果
\begin{enumerate}
 \item 每个紧集具有有限测度；
 \item 每个 Borel 集外正则；
 \item 每个开集内正则。
\end{enumerate}
紧集有限蕴含局部有限，因为每点有相对紧邻域。
\end{definition}

\begin{definition}[支撑]\label{def:2-2-2-2}
拓扑空间上 Borel 测度 $\mu$ 的支撑定义为
\[
 \supp\mu=\{x\in X:\mu(U)>0\text{ 对每个含 }x\text{ 的开集 }U\}.
\]
其补集是所有零测开集之并，因此 $\supp\mu$ 总是闭集。这个任意并是否仍零测，需要额外的可数或
紧性压缩，不能从定义推出。
\end{definition}

\begin{example}[Dirac 与原子测度的支撑]\label{ex:2-2-2-1}
在 Hausdorff 空间的 Borel $\sigma$-代数上，
\(\supp\delta_x=\{x\}\)：含 \(x\) 的开集具有 \(\delta_x\)-质量 \(1\)，而对 \(y\ne x\)，
Hausdorff 性给出 \(y\) 的一个不含 \(x\) 的开邻域。更一般地，若 \(x_1,\dots,x_r\) 两两不同且
\[
 \mu=\sum_{j=1}^r a_j\delta_{x_j},\qquad a_j>0,
\]
则
\[
 \mu(B)=\sum_{\{j:x_j\in B\}}a_j,
 \qquad \supp\mu=\{x_1,\dots,x_r\}.
\]
若改取可数多个正权原子并假设 \(\sum_j a_j<\infty\)，同一邻域判据给出
\(\supp\mu=\overline{\{x_j:j\ge1\}}\)。
\end{example}

\begin{definition}[一致紧测度族]\label{def:2-2-2-3}
概率测度族 $\mathcal P$ 称为\emph{一致紧的}（tight），如果对每个 $\varepsilon>0$，存在紧集
$K\subset X$
使
\[
 \sup_{\mu\in\mathcal P}\mu(X\setminus K)<\varepsilon.
\]
单个 Radon 概率测度构成的一点族是一致紧的，因为 $X$ 本身是开集，开集内正则给出
$\mu(K)>1-\varepsilon$ 的紧集。
\end{definition}

\begin{example}[原子质量逃逸]\label{ex:2-2-2-2}
在 $\R$ 上令 $\mu_n=\delta_n$。对任意紧集 $K$，存在 $R>0$ 使 $K\subset[-R,R]$；当整数
$n>R$ 时，
\[
  \mu_n(\R\setminus K)=1.
\]
因此 $(\mu_n)$ 不是一致紧的，尽管每个 $\mu_n$ 单独都是 Radon 概率测度。这个例子把一致紧性所排除的
“质量逃向无穷远”写成了精确的紧集外质量计算。
\end{example}

\begin{remark}[矩界给出一致紧性]
积分建立后，若一族 $\R^n$ 上的概率测度满足
$\sup_\alpha\int |x|^2\,\mu_\alpha(\dd x)\le C$，则对每个 $R>0$，
\[
  \mu_\alpha(\{|x|>R\})
  \le R^{-2}\int |x|^2\,\mu_\alpha(\dd x)
  \le C/R^2.
\]
所以取闭球 $\overline B(0,R)$ 并令 $R\to\infty$ 即得一致紧性。这是 Chebyshev 不等式的测度形式。
\end{remark}

\begin{lemma}[紧集与开集之间的帽函数]\label{lem:lch-compact-open-bump}
设 $X$ 局部紧 Hausdorff，$K\subset U$，其中 $K$ 紧、$U$ 开。则存在
$f\in C_c(X)$ 使 $0\le f\le1$、$f=1$ 于 $K$，并且 $\supp f\subset U$。
\end{lemma}

\begin{proof}
这正是上一章已经完整证明的局部 Urysohn 引理
\ref{theorem:1-3-1-2}；该定理先由相对紧邻域引理得到
$K\subset V\Subset U$，再在一点紧化中对 $K$ 与补集作连续分离。因此它给出的函数同时满足
$f|_K=1$ 与严格的支撑包含 $\supp f\subset U$，而不只是 $f=0$ 于 $X\setminus U$。
\end{proof}

\begin{lemma}[Radon 测度的开集逼近]\label{lemma:2-2-2-1}
若 $U$ 开且 $\mu$ Radon，则集合层已有
\[
 \mu(U)=\sup\{\mu(K):K\subset U,\ K\text{ 紧}\}.
 \label{eq:radon-open-set-skeleton}
\]
\end{lemma}

\begin{proof}
这正是 Radon 测度对开集的内正则性。
\end{proof}

\begin{remark}[连续函数形式]
上一帽函数引理把每个紧核 \(K\subset U\) 夹在一个函数
\(f\in C_c(X)\) 的水平集之间。第~3~章定义积分后，这一事实将给出
\[
 \mu(U)=\sup\left\{\int f\,\dd\mu:
 f\in C_c(X),\ 0\le f\le1,\ \supp f\subset U\right\}.
 \label{eq:radon-open-function-future}
\]
\end{remark}

\begin{lemma}[可数相对紧基]\label{lem:radon-relative-compact-base}
第二可数局部紧 Hausdorff 空间有一个可数基 $(V_j)$，每个 $\overline{V_j}$ 紧；还可选递增紧集
$K_j$ 使 $K_j\subset\interior K_{j+1}$ 且 $\bigcup_jK_j=X$。
\end{lemma}

\begin{proof}
若 $X=\varnothing$，结论平凡。以下设 $X\ne\varnothing$；若所得基只有有限多个元素，就周期性重复
这些元素，把它枚举为一个序列。
给定 $x\in U$ 且 $U$ 开，上一章的相对紧邻域引理
\ref{proposition:1-3-1-1} 给出 $x\in V$ 且 $\overline V\subset U$、$\overline V$ 紧。
若 $(B_i)$ 是一个可数基，再取 $x\in B_i\subset V$，便有
$\overline{B_i}\subset\overline V\subset U$。因此从该可数基中保留所有闭包紧的基元素后，所得
子族仍能细化任意邻域。枚举为 $(V_j)$。取 $K_1=\overline{V_1}$；已得 $K_j$ 后，紧集
$K_j\cup\overline{V_{j+1}}$ 可由有限个相对紧基元
$V_{i_1},\ldots,V_{i_r}$ 覆盖。令
\[
 K_{j+1}=\overline{V_{i_1}}\cup\cdots\cup\overline{V_{i_r}}.
\]
则 $K_{j+1}$ 紧，而开集 $\bigcup_{\ell=1}^rV_{i_\ell}$ 包含
$K_j\cup\overline{V_{j+1}}$ 并包含于 $K_{j+1}$。因此
$K_j\subset\interior K_{j+1}$；每个基元最终包含于某个 $K_j$，所以
$\bigcup_jK_j=X$。
\end{proof}

\begin{theorem}[自动 Radon 判据]\label{theorem:2-2-2-2}
若 $X$ 第二可数、局部紧 Hausdorff，$\mu$ 是在紧集上有限的 Borel 测度，则 $\mu$ 为 Radon
测度。并且每个有限测度 Borel 集都可由紧集从内逼近。
\end{theorem}

\begin{proof}
取上一引理的相对紧基。空开集的内正则性是平凡的。若 $U$ 是非空开集，把满足 $\overline{V_j}\subset U$ 的基元素列为
$W_1,W_2,\dots$，并令 $L_n=\bigcup_{j\le n}\overline{W_j}$。则 $L_n$ 紧、包含于 $U$，而
$\bigcup_nW_n=U$，并且 $\bigcup_{j\le n}W_j\subset L_n$。从下连续性给
\[
 \mu(U)=\lim_n\mu\Bigl(\bigcup_{j\le n}W_j\Bigr)
 \le\sup_{K\subset U,\ K\text{ 紧}}\mu(K)\le\mu(U),
\]
故开集内正则。紧耗尽还说明 $\mu$ 为 $\sigma$-有限。

先记录有限测度空间中的推广步骤。若 $O\subset X$ 是相对紧开集，则限制测度
$\mu_O(A)=\mu(A)$（$A$ 为 $O$ 的 Borel 子集）满足 $\mu_O(O)<\infty$；而 $O$ 的开子集也是
$X$ 的开集，所以上一段给出它们的紧内逼近。令 $\mathcal R_O$ 为这些 Borel 集 $A\subset O$
的集族：对每个 $\varepsilon>0$，存在 $O$ 中的相对闭集 $F\subset A$ 与相对开集 $G\supset A$，
使
\[
 \mu_O(A\setminus F)<\varepsilon,\qquad
 \mu_O(G\setminus A)<\varepsilon.
\]
开集属于 $\mathcal R_O$：上一段的紧内逼近同时给出相对闭内逼近，外逼近可取它本身。由于总测度有限，
若 $A\in\mathcal R_O$，对 $A$ 的闭内逼近取补便给 $O\setminus A$ 的开外逼近，对 $A$ 的开外
逼近取补则给 $O\setminus A$ 的闭内逼近。因此 $\mathcal R_O$ 对补封闭。若
$A=\bigcup_{n\ge1}A_n$ 且 $A_n\in\mathcal R_O$，给每个 $A_n$ 取误差小于
$\varepsilon2^{-n}$ 的开外逼近并取并，得到 $A$ 的开外逼近。再由从下连续性选 $N$，使
\[
 \mu_O\Bigl(A\setminus\bigcup_{n=1}^NA_n\Bigr)<\frac{\varepsilon}{2},
\]
并在前 $N$ 个集合内各取闭集，使误差总和小于 $\varepsilon/2$；这些闭集的有限并给出 $A$ 的闭
内逼近。确切地，若 $F_n\subset A_n$，则
\[
 A\setminus\bigcup_{n=1}^NF_n
 \subset
 \left(A\setminus\bigcup_{n=1}^NA_n\right)
 \cup\bigcup_{n=1}^N(A_n\setminus F_n),
\]
右端测度小于 $\varepsilon$。故 $\mathcal R_O$ 是含所有开集的 $\sigma$-代数，因而 $O$ 中每个
Borel 集都开外正则。

把相对紧基枚举为 $(O_j)_{j\ge1}$；它覆盖 $X$。令
\[
 D_1=O_1,\qquad D_j=O_j\setminus\bigcup_{i<j}O_i\quad(j\ge2),
\]
则 $D_j$ 是两两不交的 Borel 集，$D_j\subset O_j$，且 $\bigcup_jD_j=X$。给定 Borel 集 $E$
与 $\varepsilon>0$，在有限测度空间 $O_j$ 中对 $E\cap D_j$ 取相对开集 $U_j$，使
\[
 E\cap D_j\subset U_j\subset O_j,
 \qquad \mu\bigl(U_j\setminus(E\cap D_j)\bigr)<\varepsilon2^{-j}.
\]
因为 $O_j$ 在 $X$ 中开，每个 $U_j$ 也在 $X$ 中开。于是 $U=\bigcup_jU_j$ 是包含 $E$ 的开集，
并且
\[
 U\setminus E
 \subset\bigcup_{j\ge1}\bigl(U_j\setminus(E\cap D_j)\bigr),
 \qquad
 \mu(U\setminus E)<\varepsilon.
\]
这给出全局外正则。

最后若 $\mu(E)<\infty$，对 $E^c$ 取开集 $U\supset E^c$ 且
$\mu(U\setminus E^c)<\varepsilon$。闭集 $F=X\setminus U$ 包含于 $E$，并且
\[
 E\setminus F=E\cap U=U\setminus E^c,
 \qquad \mu(E\setminus F)<\varepsilon.
\]
由 $E\cap K_n\uparrow E$ 及 $\mu(E)<\infty$，可取足够大的紧耗尽层 $K_n$，使
$\mu(E\setminus K_n)<\varepsilon$。此时 $F\cap K_n$ 是包含于 $E$ 的紧集，而且
\[
 E\setminus(F\cap K_n)=(E\setminus F)\cup(E\setminus K_n),
 \qquad
 \mu(E\setminus(F\cap K_n))<2\varepsilon.
\]
把上述论证中的 $\varepsilon$ 换成任意给定误差的一半，便证明有限
Borel 集内正则，三条 Radon 条件全部成立。
\end{proof}

\begin{proposition}[支撑的最小性]\label{proposition:2-2-2-3}
若 $X$ 第二可数或 $\mu$ Radon，则
$\mu(X\setminus\supp\mu)=0$。若闭集 $F$ 满足 $\mu(X\setminus F)=0$，则
$\supp\mu\subset F$。所以在上述假设下，支撑是最小闭全测集。
\end{proposition}

\begin{proof}
对每个 $x\notin\supp\mu$ 取零测开邻域 $U_x$。若 $X$ 第二可数，开覆盖
$(U_x)$ 有可数子族覆盖其并。更显式地，若 $(B_j)$ 是可数基，令
$J=\{j:\mu(B_j)=0\}$。每个 $x\notin\supp\mu$ 的某个零测开邻域含有一个包含 $x$ 的基元，故
$x\in B_j$ 对某个 $j\in J$ 成立；反向包含由零测基元本身是零测开邻域得到。因此
\[
 X\setminus\supp\mu=\bigcup_{j\in J}B_j,
 \qquad \mu(B_j)=0\quad(j\in J).
\]
故该并零测。若 $\mu$ Radon，记
$U=\bigcup_xU_x$。每个紧集 $K\subset U$ 有有限子覆盖，故 $\mu(K)=0$；开集内正则给
$\mu(U)=0$。两种情形都得到第一式。

若 $F$ 闭且补集零测，任取 $x\notin F$，开邻域 $X\setminus F$ 的测度为零，所以
$x\notin\supp\mu$。故 $\supp\mu\subset F$。
\end{proof}

\begin{figure}[htbp]
\centering
\begin{tikzpicture}[node distance=17mm]
 \node[book strong node,minimum width=42mm,minimum height=18mm] (s) {闭支撑 $\supp\mu$\\每个邻域有正质量};
 \node[book warning node,left=of s] (z) {外部点\\有零测开邻域};
 \node[book warm node,right=of s] (k) {共同紧集 $K$\\包含 $1-\varepsilon$ 质量};
 \draw[book dashed arrow] (z)--node[above]{可数或紧压缩}(s);
 \draw[book accent arrow] (k)--node[above]{一致紧性}(s);
\end{tikzpicture}
\caption{支撑刻画单个测度的零质量外部；一致紧性用同一个紧集控制一族概率测度。}
\label{fig:support-and-tightness}
\end{figure}

\begin{remark}
在这里的第二可数或 Radon 假设下，支撑为空当且仅当测度为零。删除把任意开覆盖压成可数或有限
控制的机制后，零测开集的任意并可能有正测度，因此这句话不再来自定义。
\end{remark}

Radon 正则性把抽象测度与开集、紧集及紧支撑连续函数联系起来。第~7~章的 Riesz 表示定理将反向地
从 \(C_c(X)\) 上的正线性泛函恢复测度；弱收敛与 Prokhorov 定理则进一步使用这里建立的支撑和
一致紧性概念。

\begin{exercise}
设 $X$ 是非空的第二可数 LCH 空间。证明每个非零有限 Borel 测度 $\mu$ 都可唯一写成
$\mu(X)P$，其中 $P$ 是一致紧的 Radon 概率测度；并单独核对零测度。
\end{exercise}
\begin{exercise}
计算有限原子测度以及具有连续正密度的 Borel 测度的支撑。
\end{exercise}
\begin{exercise}
分别用第二可数性和 Radon 内正则性证明支撑补集零测，并指出两种压缩发生在哪一步。
\end{exercise}
\begin{exercise}
给出一族非一致紧的概率测度，并证明任何子列仍不一致紧。
\end{exercise}

正测度只能累积质量。为了处理线性观测中的抵消，下一步必须在保持可数可加性的同时分离正负部分，并给
这种抵消一个不依赖具体分解的绝对大小。

\subsection{有符号测度、复测度与全变差}\label{subsec:2-2-3}
\index{Hahn 分解}\index{Jordan 分解}\index{全变差}

正测度只会累积质量，线性问题却不断产生抵消。在实线的 Borel $\sigma$-代数上，最小的例子是
$\nu=\delta_0-\delta_1$：总质量为零，但若把两个原子分开看，绝对质量应为二。把一个有符号测度
随意写成两个正测度之差并不能解决问题，因为可以同时给两部分加上同一个正测度。我们需要一对
彼此不重叠的最小正负部分，以及一种对所有可测划分稳定的绝对值。

\begin{definition}[有限有符号测度与有限复测度]\label{def:2-2-3-1}
固定可测空间 $(X,\mathcal M)$。本节的有限有符号测度是实值可数可加函数
$\nu:\mathcal M\to\R$；有限复测度是复值可数可加函数
$\nu:\mathcal M\to\C$。这里不允许 $\pm\infty$，所以每次差和每个可数级数都在实数或复数中
解释。复值部分只使用复数的模、实虚部和级数完备性。
\end{definition}

\begin{definition}[正负集合与互相奇异]\label{def:2-2-3-2}
若 $P$ 的每个可测子集 $A$ 都满足 $\nu(A)\ge0$，称 $P$ 为 $\nu$-正集合。若可测集 $N$ 的
每个可测子集 $A$ 都满足 $\nu(A)\le0$，称 $N$ 为 $\nu$-负集合。
正测度 $\mu,\lambda$ 称互相奇异，记作 $\mu\perp\lambda$，如果存在可测 $P$ 使
$\mu(X\setminus P)=0$ 且 $\lambda(P)=0$。
\end{definition}

\begin{definition}[全变差]\label{def:2-2-3-3}
有限复测度 $\nu$ 的全变差定义为
\[
 |\nu|(E)=\sup\left\{\sum_{k=1}^r|\nu(E_k)|:
 E=\bigsqcup_{k=1}^rE_k,\ E_k\in\mathcal M,\ r<\infty\right\}.
\]
总变差范数为 $\|\nu\|_{TV}=|\nu|(X)$。
\end{definition}

\begin{example}[两个 Dirac 的抵消]\label{ex:2-2-3-2}
在 $(\R,\Borel(\R))$ 上，对 $\nu=\delta_0-\delta_1$，总质量 $\nu(\R)=0$。划分
$\{0\},\{1\},\R\setminus\{0,1\}$ 给出变差和 $1+1+0=2$。任一划分中，含 $0$ 的块贡献至多
$1$，含 $1$ 的块贡献至多 $1$，其余块贡献 $0$，故 $|\nu|(\R)=2$。
\end{example}

\begin{example}[复相位]\label{ex:2-2-3-3}
仍在 $(\R,\Borel(\R))$ 上，若 $\nu=i\delta_0-(1+i)\delta_1$，同样的原子划分给
\[
 |\nu|(X)=|i|+|-(1+i)|=1+\sqrt2.
\]
把两个系数除以其模，所得单位相位分别为 $i$ 和 $-(1+i)/\sqrt2$。一般复测度的相位表示要等到
Radon--Nikodym 定理之后才能证明。
\end{example}

\begin{remark}[由密度产生的有符号测度]\label{ex:2-2-3-1}
下一章定义积分后，可积实函数 $f$ 将给出
$\nu(E)=\int_Ef\,\dd\mu$；届时
$\nu^+,\nu^-$ 分别由 $f^+,f^-$ 给出，而
$|\nu|(E)=\int_E|f|\,\dd\mu$。这些公式将在 Radon--Nikodym 理论建立后证明。
\end{remark}

\begin{lemma}[有限有符号测度的值域有界]\label{lem:signed-measure-bounded-range}
若 $\nu$ 为有限有符号测度，则 $\sup_{E\in\mathcal M}|\nu(E)|<\infty$。
\end{lemma}

\begin{proof}
反设某个集合 $R_0=X$ 上的限制值域无界。若 $R$ 上限制无界，可取 $A\subset R$ 使
$|\nu(A)|>|\nu(R)|+1$。由
$\nu(R\setminus A)=\nu(R)-\nu(A)$，
\[
 |\nu(R\setminus A)|
 \ge\bigl||\nu(A)|-|\nu(R)|\bigr|>1,
\]
而 $|\nu(A)|>1$。至少一块上的限制仍无界；否则若两块上的绝对值分别由 $M_1,M_2$ 控制，则任意
$C\subset R$ 满足
\[
 |\nu(C)|
 \le|\nu(C\cap A)|+|\nu(C\cap(R\setminus A))|
 \le M_1+M_2,
\]
与 $R$ 上无界矛盾。保留限制仍无界的一块作下一轮 $R$，把另一块记为 $E_k$。
由此得到两两不交且 $|\nu(E_k)|>1$ 的序列。可数可加性却要求
$\sum_k\nu(E_k)=\nu(\bigsqcup_kE_k)$ 收敛，收敛级数的项必须趋于零，矛盾。
\end{proof}

\begin{lemma}[正子集抽取]\label{lem:hahn-positive-subset}
若 $\nu(A)>0$，则 $A$ 含有一个 $\nu$-正集合 $P$，并且 $\nu(P)>0$。
\end{lemma}

\begin{proof}
置 $R_0=A$。若某一步的 $R_{n-1}$ 已是正集合，便取 $P=R_{n-1}$；下述递推估计给
$\nu(P)\ge\nu(A)>0$，结论已经成立。否则由
\cref{lem:signed-measure-bounded-range}知下式的下确界是有限实数，令
\[
 \alpha_n=\inf\{\nu(B):B\subset R_{n-1},\ B\in\mathcal M\}<0
\]
并选 $B_n\subset R_{n-1}$ 使 $\nu(B_n)<\alpha_n/2<0$，再令
$R_n=R_{n-1}\setminus B_n$。各 $B_n$ 两两不交，且
\[
 \nu(R_n)=\nu(A)-\sum_{j=1}^n\nu(B_j)\ge\nu(A)>0.
\]
若过程不停，令 $P=\bigcap_nR_n$。可数可加性给
$\nu(P)=\nu(A)-\sum_n\nu(B_n)\ge\nu(A)>0$，并且
$\nu(B_n)\to0$。若 $P$ 含子集 $C$ 且 $\nu(C)<0$，则
$C\subset R_{n-1}$ 使 $\alpha_n\le\nu(C)$，从而
$\nu(B_n)<\nu(C)/2<0$ 对每个 $n$ 成立，与 $\nu(B_n)\to0$ 矛盾。因此 $P$ 为正集合。
\end{proof}

\begin{theorem}[Hahn 分解]\label{theorem:2-2-3-1}
存在可测分割 $X=P\sqcup N$，其中 $P$ 是 $\nu$-正集合、$N$ 是 $\nu$-负集合。若
$X=P'\sqcup N'$ 是另一 Hahn 分解，则 $P\mathbin\triangle P'$ 是 $\nu$-零集，即其每个可测
子集的 $\nu$-值均为零。
\end{theorem}

\begin{proof}
令
$\alpha=\sup\{\nu(P):P\text{ 为正集合}\}$；上一有界性引理保证 $\alpha<\infty$。取正集合
$P_n$ 使 $\nu(P_n)>\alpha-2^{-n}$，并置 $P=\bigcup_nP_n$。若 $A\subset P$，把
$A$ 按
\[
 D_1=P_1,\qquad D_n=P_n\setminus\bigcup_{j<n}P_j\quad(n\ge2)
\]
不交化。因 $A\cap D_n\subset P_n$，每项 $\nu(A\cap D_n)\ge0$，可数可加性给
\[
 \nu(A)=\sum_{n=1}^{\infty}\nu(A\cap D_n)\ge0.
\]
故
$P$ 是正集合。它包含每个 $P_n$，所以
$\nu(P)\ge\sup_n\nu(P_n)=\alpha$；按 $\alpha$ 的定义又不超过 $\alpha$，故等号成立。

令 $N=X\setminus P$。若 $N$ 不是负集合，存在 $A\subset N$ 且 $\nu(A)>0$。正子集抽取引理给
$Q\subset A$，其中 $Q$ 正且 $\nu(Q)>0$。于是 $P\sqcup Q$ 正而
$\nu(P\sqcup Q)>\alpha$，矛盾。因此 $N$ 负。

若还有 $P',N'$，则 $P\cap N'$ 同时是正、负集合，它的每个可测子集测度为零；
$P'\cap N$ 也同时包含于正集合 $P'$ 和负集合 $N$，故同样为零。二者之并是
$P\mathbin\triangle P'$，故为 $\nu$-零集。
\end{proof}

\begin{theorem}[Jordan 分解]\label{theorem:2-2-3-2}
存在唯一互相奇异的正测度 $\nu^+,\nu^-$，使
$\nu=\nu^+-\nu^-$。而且
\[
 \nu^+(E)=\sup_{A\subset E}\nu(A),\qquad
 \nu^-(E)=-\inf_{A\subset E}\nu(A),
 \label{eq:jordan-extrema}
\]
上、下确界均取可测子集。
\end{theorem}

\begin{proof}
取 Hahn 分解 $X=P\sqcup N$，定义
\[
 \nu^+(E)=\nu(E\cap P),\qquad
 \nu^-(E)=-\nu(E\cap N).
\]
正负集合性质说明二者为正测度；它们分别集中于 $P,N$，因而互相奇异，且差为 $\nu$。
对 $A\subset E$，
\[
 \nu(A)=\nu(A\cap P)+\nu(A\cap N)
 \le\nu(A\cap P)\le\nu(E\cap P)=\nu^+(E),
\]
其中第一步用 $\nu(A\cap N)\le0$，第二步用正集合 $P$ 上的单调性；取 $A=E\cap P$ 达到上界。
另一方面，
\[
 \nu(A)\ge\nu(A\cap N),
 \qquad
 \nu(E\cap N)=\nu(A\cap N)+\nu((E\setminus A)\cap N)
 \le\nu(A\cap N),
\]
所以 $\nu(A)\ge\nu(E\cap N)=-\nu^-(E)$；取 $A=E\cap N$ 达到下界。这就得到
\eqref{eq:jordan-extrema}。

若 $\nu=\alpha-\beta$，其中 $\alpha,\beta$ 为有限正测度，则对 $A\subset E$ 有
$\nu(A)\le\alpha(A)\le\alpha(E)$，故上确界公式给 $\nu^+(E)\le\alpha(E)$。又有
$-\nu(A)\le\beta(A)\le\beta(E)$，故下确界公式给 $\nu^-(E)\le\beta(E)$。若再有
$\alpha\perp\beta$，取可测集 $S$，使
$\alpha(X\setminus S)=0$、$\beta(S)=0$。对每个 $E$，
\[
 \nu(E\cap S)=\alpha(E\cap S)-\beta(E\cap S)=\alpha(E),
\]
所以上确界公式给 $\nu^+(E)\ge\alpha(E)$；结合上一不等式得到 $\nu^+=\alpha$。同时，
$-\nu(E\setminus S)=\beta(E)$，下确界公式给 $\nu^-(E)\ge\beta(E)$；与已经证明的反向不等式
合并得到 $\nu^-=\beta$。因此互异差分解唯一。
\end{proof}

\begin{remark}[回到两个原子]
对 $\delta_0-\delta_1$ 可取 $P=\{0\}$、$N=X\setminus\{0\}$；零测区域如何分配不影响结果。
于是 $\nu^+=\delta_0,\nu^-=\delta_1$。在 \eqref{eq:jordan-extrema} 中，包含于 $E$ 的集合能获得
的最大净质量恰为 $\1_E(0)$，最小净质量恰为 $-\1_E(1)$。Hahn 分的是空间，Jordan 分的是测度。
\end{remark}

\begin{theorem}[全变差测度]\label{theorem:2-2-3-3}
由有限划分上确界定义的 $|\nu|$ 是有限正测度，并满足
$|\nu(E)|\le|\nu|(E)$。若 $\nu$ 为实值，则
\[
 |\nu|=\nu^++\nu^-.
\]
\end{theorem}

\begin{proof}
先证有限性。实部、虚部分别是有限有符号测度。用已经证明的 Jordan 分解置
\[
 \lambda_R=(\Re\nu)^++(\Re\nu)^-,\qquad
 \lambda_I=(\Im\nu)^++(\Im\nu)^-;
\]
它们是有限正测度。对任意有限划分，
\[
 \sum_k|\nu(E_k)|
 \le\sum_k\bigl(|\Re\nu(E_k)|+|\Im\nu(E_k)|\bigr)
 \le \lambda_R(E)+\lambda_I(E),
\]
故上确界有限。单块划分给 $|\nu(E)|\le|\nu|(E)$。

设 $E_j$ 两两不交，并记 $E=\bigsqcup_{j\ge1}E_j$。固定 $N$。对每个 $1\le j\le N$，在
$E_j$ 内取有限可测划分 $(E_{j,k})_{1\le k\le r_j}$，使
\[
 \sum_{k=1}^{r_j}|\nu(E_{j,k})|>|\nu|(E_j)-\frac{\varepsilon}{N}.
\]
合并这些划分，并另加余块
$E\setminus\bigcup_{j\le N}E_j$，便得到 $E$ 的有限划分。因此
\[
 |\nu|(E)\ge\sum_{j=1}^N|\nu|(E_j)-\varepsilon.
\]
先令 $\varepsilon\downarrow0$，再令 $N\to\infty$，得到“$\ge$”。反过来，任取 $E$ 的有限划分
$E=\bigsqcup_{k=1}^rA_k$，用可数可加性和三角不等式，
\[
 \sum_{k=1}^r|\nu(A_k)|
 \le\sum_{k=1}^r\sum_j|\nu(A_k\cap E_j)|
 \le\sum_j|\nu|(E_j).
\]
对所有划分取上确界得“$\le$”，所以 $|\nu|$ 可数可加。

实值情形，对任意划分 $E=\bigsqcup_kE_k$，由
$\nu(E_k)=\nu^+(E_k)-\nu^-(E_k)$ 得
\[
 \sum_k|\nu(E_k)|
 \le\sum_k\bigl(\nu^+(E_k)+\nu^-(E_k)\bigr)
 =\nu^+(E)+\nu^-(E).
\]
用 Hahn 分区把 $E$ 切成 $E\cap P,E\cap N$，则
\[
 |\nu(E\cap P)|+|\nu(E\cap N)|
 =\nu(E\cap P)-\nu(E\cap N)
 =\nu^+(E)+\nu^-(E),
\]
所以该划分达到上界，故等号成立。
\end{proof}

\begin{theorem}[复测度极分解]\label{theorem:2-2-3-4}
对有限复测度 $\nu$，存在 $|\nu|$-可测函数 $h$，满足
$|h|=1$ 在 $|\nu|$-几乎处处成立，并且
\[
 \nu(E)=\int_Eh\,\dd|\nu|.
\]
\end{theorem}

证明见\cref{theorem:3-3-1-3}；其核心工具是 Radon--Nikodym 定理。

\begin{figure}[htbp]
\centering
\begin{tikzpicture}[node distance=12mm and 14mm]
 \node[book strong node] (x) {空间 $X=P\sqcup N$};
 \node[book node,below left=of x] (p) {$P$：所有子集非负\\$\nu^+$};
 \node[book warning node,below right=of x] (n) {$N$：所有子集非正\\$\nu^-$};
 \node[book warm node,below=22mm of x] (v) {$|\nu|=\nu^++\nu^-$\\把负区翻正后相加};
 \draw[book arrow] (x)--(p); \draw[book arrow] (x)--(n);
 \draw[book arrow] (p)--(v); \draw[book arrow] (n)--(v);
\end{tikzpicture}
\caption{Hahn 分区、Jordan 分量与全变差的关系；复测度的对应结构由全变差与单位相位给出。}
\label{fig:hahn-jordan-variation}
\end{figure}

\begin{remark}
“正集合”要求它的每个可测子集非负，不只是整块质量非负。Jordan 分解的唯一性来自正负部分互相
奇异；若删除互异条件，给两部分同时加任意正测度便得到另一差分解。全变差从有限划分定义，其
可数可加性则是上一定理的结论。
\end{remark}

全变差在第~3~章控制复积分并产生 Radon--Nikodym 极分解；它也把概率测度之间的全变差距离写成
可计算的上确界。第~7~章研究 \(C_0(X)\) 的对偶时，还会把这一结构推广到实、复 Radon 测度。

\begin{exercise}
计算任意有限原子复测度 $\sum_{j=1}^rc_j\delta_{x_j}$ 的全变差。
\end{exercise}
\begin{exercise}
若 $\nu=\alpha-\beta$，其中 $\alpha,\beta$ 为正测度，证明
$\nu^+\le\alpha$、$\nu^-\le\beta$。
\end{exercise}
\begin{exercise}
设 $(\alpha_j)$ 与 $(\beta_k)$ 是有限族或可数族正测度，并且
$\alpha_j\perp\beta_k$ 对每一对 $(j,k)$ 成立。证明
$\sum_j\alpha_j\perp\sum_k\beta_k$；若要求两边仍为有限正测度，写出对可数和所需的总质量条件。
再构造一个例子，说明仅有 $\alpha_j\perp\beta_j$ 并不足够。
\end{exercise}
\begin{exercise}
不使用 Jordan 分解，直接从划分定义证明 $|\nu(E)|\le|\nu|(E)$。
\end{exercise}

测度一旦能够携带正负和复相位，下一步便是研究它怎样沿映射移动、在积空间组合，以及怎样由实线
上的单调函数编码。

\section{传输、乘积与一维分布}\label{sec:02-03}

\subsection{可测映射、推前测度与分布}\label{subsec:2-3-1}
\index{推前测度}\index{分布}

一个测度未必直接写在我们关心的空间上。样本经过统计量变成一个数，Euclidean 点经过坐标投影
变成一个坐标，几何区域经过变换被搬到另一处。目标集合 $B$ 的质量只能通过源中的逆像
$T^{-1}(B)$ 计算；这正是推前的定义。它保留由 $T$ 能观察到的信息，也会不可逆地丢掉同一纤维内的
差别。

\begin{definition}[可测映射]\label{def:2-3-1-1}
映射 $T:(X,\mathcal A)\to(Y,\mathcal B)$ 称为可测映射，如果
$T^{-1}(B)\in\mathcal A$ 对每个 $B\in\mathcal B$ 成立。若 $T$ 为双射且 $T,T^{-1}$ 都可测，
称为可测同构。
\end{definition}

\begin{definition}[推前测度]\label{def:2-3-1-2}
给定 $X$ 上测度 $\mu$ 和可测映射 $T:X\to Y$，定义
\[
 T_*\mu(B)=\mu(T^{-1}(B)),\qquad B\in\mathcal B.
\]
若 $\mu(X)=1$，称 $(X,\mathcal A,\mu)$ 为概率空间，其中的可测映射暂称随机变量，并把
$T_*\mu$ 记作 $\Law(T)$。若可测同构还满足 $T_*\mu=\nu$，称其为测度空间同构。
\end{definition}

逆像保持空集、补集和可数并，因此 $T_*\mu$ 的可数可加性由 $\mu$ 的可数可加性逐项得到。

\begin{proposition}[函子性]\label{proposition:2-3-1-1}
若复合有意义，则
\[
 (S\circ T)_*\mu=S_*(T_*\mu),\qquad (\id_X)_*\mu=\mu.
\]
\end{proposition}

\begin{proof}
对目标可测集 $C$，
\[
 (S\circ T)_*\mu(C)
 =\mu(T^{-1}(S^{-1}C))
 =T_*\mu(S^{-1}C)=S_*(T_*\mu)(C).
\]
恒等映射的逆像不改变集合，第二式随之成立。
\end{proof}

\begin{proposition}[有限简单测试公式]\label{proposition:2-3-1-1a}
若 $B_1,\dots,B_r\in\mathcal B$ 两两不交且 $a_k\ge0$，则
\[
 \sum_{k=1}^ra_k(T_*\mu)(B_k)
 =\sum_{k=1}^ra_k\mu(T^{-1}B_k).
\]
\end{proposition}

\begin{proof}
每一项由推前定义相等；逆像还保持两两不交，所以两边确实是同一个有限分割上的质量和。
\end{proof}

\begin{proposition}[连续推前与支撑]\label{proposition:2-3-1-1b}
设 $X,Y$ 为拓扑空间，$\mu$ 为 $X$ 上的 Borel 测度，且
$\mu(X\setminus\supp\mu)=0$。若 $T:X\to Y$ 连续，则
\[
 \supp(T_*\mu)=\overline{T(\supp\mu)}.
\]
\end{proposition}

\begin{proof}
若 $x\in\supp\mu$，含 $T(x)$ 的每个开集 $V$ 有开逆像 $T^{-1}V$ 含 $x$，故
$T_*\mu(V)=\mu(T^{-1}V)>0$。所以 $T(\supp\mu)\subset\supp(T_*\mu)$；右侧支撑闭，得到一个
包含关系。

若 $y\notin\overline{T(\supp\mu)}$，取开邻域
$V=Y\setminus\overline{T(\supp\mu)}$。它的逆像与 $\supp\mu$ 不交，因此
$T_*\mu(V)\le\mu(X\setminus\supp\mu)=0$，所以 $y$ 不在推前支撑中。
\end{proof}

\begin{theorem}[推前积分公式]\label{theorem:2-3-1-2}
对非负可测函数 $f:Y\to[0,\infty]$，
\[
 \int_Y f\,\dd(T_*\mu)=\int_Xf\circ T\,\dd\mu.
\]
\end{theorem}

非负积分将在下一章由简单函数下逼近定义，届时
\cref{proposition:3-1-1-5} 将证明上式；可积实值和复值版本分别由正负部与实虚部得到。

\begin{proposition}[连续推前保持 Radon]\label{proposition:2-3-1-3}
设 $X,Y$ 为局部紧 Hausdorff 空间，$\mu$ 为有限 Radon 测度，$T:X\to Y$ 连续，则
$T_*\mu$ 为 Radon。若 $\mu$ 不必有限而 $T$ 为固有映射（proper map），即每个紧集的逆像紧，
则结论仍成立。
\end{proposition}

\begin{proof}
先记录有限 Radon 测度的一个推论：它对每个 Borel 集都可由紧集从内逼近。事实上，若
$E\subset X$ 为 Borel 集且 $\varepsilon>0$，由开集 $X$ 的内正则性可取紧集 $K_0\subset X$，使
$\mu(X\setminus K_0)<\varepsilon/2$；又由 $E^c$ 的外正则性可取开集 $O\supset E^c$，使
$\mu(O)<\mu(E^c)+\varepsilon/2$。由于 $\mu$ 有限，
\[
 K:=K_0\setminus O\subset E,\qquad
 \mu(E\setminus K)
 \le \mu(X\setminus K_0)+\mu(O\setminus E^c)<\varepsilon.
\]
这里 $K$ 是紧集。

现在先设 $\mu$ 有限，并记 $\nu=T_*\mu$。若 $B\subset Y$ 为 Borel 集，对 $T^{-1}B$ 应用上面的
紧内逼近，得到紧集 $K\subset T^{-1}B$，使
$\mu(T^{-1}B\setminus K)<\varepsilon$；其像 $T(K)\subset B$ 紧且
\[
 \nu(T(K))=\mu(T^{-1}(T(K)))\ge\mu(K).
\]
从而
\[
 0\le \nu(B)-\nu(T(K))
 \le \mu(T^{-1}B)-\mu(K)
 =\mu(T^{-1}B\setminus K)<\varepsilon.
\]
这就得到 $\nu$ 对 Borel 集的紧内逼近。再对 $B^c$ 作
紧内逼近：若紧集 $L\subset B^c$ 满足 $\nu(B^c\setminus L)<\varepsilon$，则开集
$Y\setminus L$ 包含 $B$，并且
\[
 \nu((Y\setminus L)\setminus B)=\nu(B^c\setminus L)<\varepsilon.
\]
这给出外正则性。最后，目标紧集的质量至多 $\nu(Y)=\mu(X)<\infty$，故 $\nu$ 是 Radon 测度。

再设 $T$ 固有。目标紧集 $L$ 的逆像紧，故
$T_*\mu(L)=\mu(T^{-1}L)<\infty$。还需证明正则性。局部紧 Hausdorff 空间之间的固有连续映射是闭
映射：若 $F$ 闭、$y\notin T(F)$，取含 $y$ 的相对紧开邻域 $V$。集合
$F\cap T^{-1}(\overline V)$ 是紧集，故其像
$T(F\cap T^{-1}(\overline V))$ 是 $\overline V$ 中的闭集，而且不含 $y$。因此可在 $V$ 中取
$y$ 的开邻域 $W$，使它与这个像不交。若 $z\in W\cap T(F)$，则
$z\in\overline V$，其某个 $F$ 中原像会落在 $F\cap T^{-1}(\overline V)$，矛盾。所以
$W\cap T(F)=\varnothing$，从而 $T(F)$ 闭。

仍记 $\nu=T_*\mu$。若 $V\subset Y$ 开，则 $T^{-1}V$ 开。若
$\mu(T^{-1}V)<\infty$，给定 $\varepsilon>0$，源测度的开集内正则性给出紧集
$K\subset T^{-1}V$，使 $\mu(K)>\mu(T^{-1}V)-\varepsilon$；若
$\mu(T^{-1}V)=\infty$，则对每个 $M>0$ 可取这样的紧集使 $\mu(K)>M$。两种情形中
$T(K)\subset V$ 都是紧集，且 $\nu(T(K))\ge\mu(K)$。因此 $\nu$ 对开集内正则。

再证外正则。给定 Borel 集 $B\subset Y$；若 $\nu(B)=\infty$，外正则等式由单调性直接成立。
若 $\nu(B)<\infty$，源测度的外正则性给出开集 $O\supset T^{-1}B$，使
\[
 \mu(O)<\mu(T^{-1}B)+\varepsilon.
\]
闭映射性使
$V=Y\setminus T(X\setminus O)$ 为含 $B$ 的开集，且 $T^{-1}V\subset O$，所以
\[
 \nu(V)=\mu(T^{-1}V)\le\mu(O)<\nu(B)+\varepsilon.
\]
这证明了外正则性；目标紧集的有限性已在上面证明，故 $\nu$ 是 Radon 测度。
\end{proof}

\begin{proposition}[分布的集合层识别]\label{proposition:2-3-1-4}
若随机变量 $\xi,\eta$ 同分布，则
$\Prob(\xi\in B)=\Prob(\eta\in B)$ 对每个目标 Borel 集成立。反过来，若该等式在一个生成
Borel $\sigma$-代数的 $\pi$-系统上成立，则 $\xi,\eta$ 同分布。
\end{proposition}

\begin{proof}
第一句是推前测度相等的展开。反向中两个分布都是概率测度，故在生成 $\pi$-系统上相等后可用
\cref{theorem:2-2-1-2}。
\end{proof}

\begin{example}[经验测度]\label{ex:2-3-1-1}
给定 $x_1,\dots,x_n\in X$，在有限集合 $\{1,\dots,n\}$ 上取均匀概率，令 $T(k)=x_k$，则
\[
 T_*\Bigl(\frac1n\sum_{k=1}^n\delta_k\Bigr)
 =\frac1n\sum_{k=1}^n\delta_{x_k}=:L_n.
\]
对有限简单测试函数，两边都化为有限样本和；经验测度因此把有限数据变成测度对象。
\end{example}

\begin{example}[坐标分布]\label{ex:2-3-1-2}
若 $\mu$ 是 $\R^n$ 上测度，坐标映射 $x\mapsto x_j$ 的推前给出第 $j$ 坐标分布；线性函数
$x\mapsto\langle t,x\rangle$ 的推前记录沿方向 $t$ 的一维观测。例如在 \(\R^2\) 上取
\[
 \mu=\frac12\delta_{(0,1)}+\frac12\delta_{(2,-1)}.
\]
若 \(\pi_1,\pi_2\) 是两个坐标投影，则
\[
 (\pi_1)_*\mu=\frac12\delta_0+\frac12\delta_2,\qquad
 (\pi_2)_*\mu=\frac12\delta_1+\frac12\delta_{-1}.
\]
对 \(t=(1,1)\)，两点的内积分别为 \(1,1\)，所以
\((x\mapsto\langle t,x\rangle)_*\mu=\delta_1\)。\Cramer{}--Wold 方法将研究全部这类一维观测
是否足以识别高维分布。
\end{example}

\begin{example}[非单射丢失信息]\label{ex:2-3-1-3}
$T(x)=x^2$ 满足 $T_*\delta_1=T_*\delta_{-1}=\delta_1$。推前不能区分落在同一纤维中的两个点，
所以没有单射条件便不能反演。
\end{example}

\begin{example}[平方折叠的集合计算]\label{ex:2-3-1-4}
令 $\mu=\tfrac12m|_{[-1,1]}$，$T(x)=x^2$。若 $0\le a<b\le1$，则
\[
 T^{-1}((a,b])=[-\sqrt b,-\sqrt a)\cup(\sqrt a,\sqrt b],
\]
除端点约定外两段长度都为 $\sqrt b-\sqrt a$，所以
\[
 (T_*\mu)((a,b])=\sqrt b-\sqrt a.
\]
第~3~章的换元公式将由这个增量导出推前测度的密度。
\end{example}

\begin{figure}[htbp]
\centering
\begin{tikzpicture}[node distance=17mm and 20mm]
 \node[book strong node,minimum width=38mm] (x) {源空间 $(X,\mu)$\\纤维 $T^{-1}(y)$};
 \node[book strong node,right=of x,minimum width=38mm] (y) {目标空间 $(Y,T_*\mu)$\\集合 $B$};
 \draw[book accent arrow] (x)--node[above]{$T$}(y);
 \draw[book dashed arrow,bend right=22] (y.south)--node[below]{$B\mapsto T^{-1}(B)$}(x.south);
\end{tikzpicture}
\caption{推前沿 $T$ 移动质量，却以逆像计算目标集合。}
\label{fig:pushforward-fibres}
\end{figure}

\begin{remark}
可测映射只要求可测集的逆像可测；可测集的像未必可测。同分布也不表示两个随机变量在同一概率
空间上相等。无限 Radon 测度的推前若缺少固有性，目标紧集的逆像可能携带无穷质量，局部
有限性便会失败。
\end{remark}

对 \(x\mapsto\langle t,x\rangle\) 的推前取复指数积分，便得到特征函数；因此方向投影的一维分布
是 \Cramer{}--Wold 方法的基本数据。Markov 核把“每个起点对应一个目标测度”参数化，其确定性特例
正是 \(x\mapsto\delta_{T(x)}\)。这些概念将在概率论部分定义。

\begin{exercise}
从逆像恒等式证明推前函子性。
\end{exercise}
\begin{exercise}
计算仿射映射对有限原子测度和区间均匀概率的推前。
\end{exercise}
\begin{exercise}
分别写出支撑公式两个包含关系中使用的假设，并给出删除连续性后的反例。
\end{exercise}
\begin{exercise}
补全固有连续映射为闭映射的证明，并据此证明无限 Radon 推前命题。
\end{exercise}
\begin{exercise}
在同一概率空间上构造同分布但不几乎处处相等的两个随机变量。
\end{exercise}

单个映射只搬运一份质量。要同时描述两个观测及其联合事件，还需在矩形上规定质量，并证明这份有限数据
能够扩张到乘积 $\sigma$-代数。

\subsection{乘积 \texorpdfstring{$\sigma$}{sigma}-代数、截面与乘积测度}\label{subsec:2-3-2}
\index{乘积测度}\index{截面}

两个空间中的有限观测组合成矩形 $A\times B$。若规定矩形质量为 $\mu(A)\nu(B)$，存在性需要
把这个规则扩张到可数生成事件，唯一性需要证明矩形足以识别扩张。下面直接在有限简单函数的质量和
上验证预测度；第~3~章的 Tonelli 定理将以乘积测度的存在为基础。

\begin{definition}[乘积 $\sigma$-代数]\label{def:2-3-2-1}
$\mathcal A\otimes\mathcal B$ 是所有可测矩形 $A\times B$ 生成的 $\sigma$-代数。对任意指标族，
柱 $\sigma$-代数由只限制有限多个坐标的坐标逆像生成。
\end{definition}

\begin{definition}[截面]\label{def:2-3-2-2}
对 $E\subset X\times Y$，定义
\[
 E_x=\{y:(x,y)\in E\},\qquad E^y=\{x:(x,y)\in E\}.
\]
对函数 $f$，记 $f_x(y)=f(x,y)$。
\end{definition}

\begin{proposition}[截面可测]\label{proposition:2-3-2-1}
若 $E\in\mathcal A\otimes\mathcal B$，则每个 $E_x\in\mathcal B$、每个
$E^y\in\mathcal A$。
\end{proposition}

\begin{proof}
固定 $x$，令 $\mathcal C_x=\{E:E_x\in\mathcal B\}$。矩形的截面是 $B$ 或空集；截面运算保持
补集和可数并，所以 $\mathcal C_x$ 是含所有矩形的 $\sigma$-代数，遂包含
$\mathcal A\otimes\mathcal B$。固定 $y$ 的证明相同。
\end{proof}

\begin{example}[逐截面可测仍不够]\label{ex:2-3-2-4}
取 $A\subset X$ 为 $\mathcal A$-不可测集，取 $B\in\mathcal B$ 且
$0<\nu(B)<\infty$，令 $E=A\times B$。每个竖截面是 $B$ 或空集，因而可测；但若
$y\in B$，水平截面 $E^y=A$ 不可测，所以上一命题的逆否命题给出
$E\notin\mathcal A\otimes\mathcal B$。而且
$x\mapsto\nu(E_x)=\nu(B)\1_A(x)$ 不可测。逐条截面可测与联合可测是两件事。
\end{example}

\begin{definition}[边缘测度]\label{def:2-3-2-4}
若 $\pi$ 是 $(X\times Y,\mathcal A\otimes\mathcal B)$ 上的测度，坐标投影可测；推前
$(\pi_X)_*\pi$ 与 $(\pi_Y)_*\pi$ 称为 $\pi$ 的两个边缘测度。固定边缘寻找联合测度称为耦合
问题。
\end{definition}

\begin{definition}[乘积测度]\label{def:2-3-2-3}
对 $\sigma$-有限测度 $\mu,\nu$，满足
\[
 (\mu\otimes\nu)(A\times B)=\mu(A)\nu(B),
 \qquad A\in\mathcal A, B\in\mathcal B,
\]
的测度称为乘积测度；约定 $0\cdot\infty=0$。
\end{definition}

$\sigma$-有限性对唯一性是必要的。\Cref{ex:2-3-2-5} 将在
$([0,1],\#)\times([0,1],m)$ 上给出两个测度：它们在每个矩形上都等于
$\#(A)m(B)$，在对角线上的质量却分别为 $0$ 与 $1$。因此存在性证明之前必须明确限制唯一性的范围。

\begin{lemma}[有限简单函数的质量和]\label{lem:product-simple-section-functional}
设 $\mu(X)<\infty$。若非负简单函数有不交表示
$g=\sum_{j=1}^rc_j\1_{A_j}$，定义
$I_\mu(g)=\sum_jc_j\mu(A_j)$。该数与不交表示无关。若
$0\le g_n\uparrow g$，所有函数均简单、有共同有限界，则
$I_\mu(g_n)\uparrow I_\mu(g)$。
\end{lemma}

\begin{proof}
两种表示用所有 $A_j\cap B_k$ 作共同有限分割；在每个非空交块上两边系数相同，故质量和相同。
设所有函数由 $M$ 控制。给定 $\varepsilon>0$，集合
$D_n=\{x:g(x)-g_n(x)<\varepsilon\}$ 递增至 $X$；这是因为每一点上的实数列递增至 $g(x)$。
测度从下连续和 $\mu(X)<\infty$ 给 $\mu(X\setminus D_n)\to0$。共同有限分割还给简单质量和的
有限线性：在 $g,g_n$ 的所有水平集交块上逐项相减，得到
$I_\mu(g-g_n)=I_\mu(g)-I_\mu(g_n)$。因而
\[
 0\le I_\mu(g)-I_\mu(g_n)=I_\mu(g-g_n)
 \le\varepsilon\mu(X)+M\mu(X\setminus D_n).
\]
先令 $n\to\infty$，再令 $\varepsilon\downarrow0$，得到所需收敛。
\end{proof}

\begin{theorem}[乘积测度存在唯一性]\label{theorem:2-3-2-2}
若 $\mu,\nu$ 都 $\sigma$-有限，则存在唯一乘积测度 $\mu\otimes\nu$。
\end{theorem}

\begin{proof}
先设两测度有限。可测矩形构成半环，其有限不交并构成环 $\mathcal R$。若
$R\in\mathcal R$，则 $R_x$ 是有限个可测集合之并，函数
$g_R(x)=\nu(R_x)$ 是非负简单函数。确切地，若
$R=\bigcup_{l=1}^r(A_l\times B_l)$，则由 $(A_l)_{l=1}^r$ 生成的至多 $2^r$ 个胞腔
\[
 C_J=\Bigl(\bigcap_{l\in J}A_l\Bigr)
      \cap\Bigl(\bigcap_{l\notin J}A_l^c\Bigr),
 \qquad J\subset\{1,\dots,r\},
\]
使 $R_x=\bigcup_{l\in J}B_l$ 在每个非空 $C_J$ 上固定。因此 $g_R$ 至多取 $2^r$ 个值。定义
\[
 \rho_0(R)=I_\mu(g_R).
 \label{eq:product-premeasure}
\]
共同矩形分割和上一引理说明定义与 $R$ 的表示无关，并且
$\rho_0(A\times B)=\mu(A)\nu(B)$。

若 $R_1,\dots,R_N$ 两两不交，则截面也两两不交，所以对每个 $x$，
\[
 g_{\bigsqcup_{k=1}^NR_k}(x)
 =\nu\!\left(\bigsqcup_{k=1}^N(R_k)_x\right)
 =\sum_{k=1}^N\nu((R_k)_x)
 =\sum_{k=1}^Ng_{R_k}(x).
\]
在这些简单函数水平集的共同有限分割上，$I_\mu$ 的有限线性性给
\[
 \rho_0\!\left(\bigsqcup_{k=1}^NR_k\right)
 =\sum_{k=1}^N\rho_0(R_k).
\]

若 $R=\bigsqcup_{k\ge1}R_k$ 且各集合都在环中，令
$S_N=\bigsqcup_{k\le N}R_k$。逐个 $x$ 有
$\nu((S_N)_x)\uparrow\nu(R_x)$；极限仍是有界简单函数，因为 $R$ 属于矩形环且 $\nu(Y)<\infty$。
有限简单单调引理给
\[
 \rho_0(R)=\lim_N\rho_0(S_N)=\sum_{k=1}^{\infty}\rho_0(R_k).
\]
所以 $\rho_0$ 是预测度。环版 \Caratheodory{} 扩张给
$\mathcal A\otimes\mathcal B$ 上的测度；有限测度唯一性给唯一。

一般情形取两两不交的可测层 $(X_i)_{i\ge1}$、$(Y_j)_{j\ge1}$，使
\[
 X=\bigsqcup_iX_i,\qquad Y=\bigsqcup_jY_j,\qquad
 \mu(X_i)<\infty,\quad \nu(Y_j)<\infty.
\]
在 $X_i\times Y_j$ 上对限制测度应用有限情形，所得乘积记为 $\rho_{i,j}$。对
$E\in\mathcal A\otimes\mathcal B$ 定义
\[
 \rho(E)=\sum_{i,j\ge1}\rho_{i,j}\bigl(E\cap(X_i\times Y_j)\bigr).
 \label{eq:sigma-finite-product-patching}
\]
若 $E_n$ 两两不交，则每个块内的限制仍两两不交；先用 $\rho_{i,j}$ 的可数可加性，再对非负三重
级数换序，得到
\[
 \rho\Bigl(\bigsqcup_nE_n\Bigr)
 =\sum_{i,j,n}\rho_{i,j}\bigl(E_n\cap(X_i\times Y_j)\bigr)
 =\sum_n\rho(E_n).
\]
所以 $\rho$ 是测度。对矩形 $A\times B$，
\[
\begin{aligned}
 \rho(A\times B)
 &=\sum_{i,j}\mu(A\cap X_i)\nu(B\cap Y_j)\\
 &=\left(\sum_i\mu(A\cap X_i)\right)
   \left(\sum_j\nu(B\cap Y_j)\right)
 =\mu(A)\nu(B),
\end{aligned}
\]
其中各项非负，$0\cdot\infty=0$ 的情形也由逐项为零直接得到。这证明全局存在性。

这里块限制的可测性可直接核对。固定 $i,j$，令 $\mathcal T_{i,j}$ 为所有满足
\[
 E\cap(X_i\times Y_j)
 \in(\mathcal A|_{X_i})\otimes(\mathcal B|_{Y_j})
\]
的 $E\subset X\times Y$。该类对补集和可数并封闭；并且
\[
 (A\times B)\cap(X_i\times Y_j)
 =(A\cap X_i)\times(B\cap Y_j),
\]
所以它包含所有可测矩形，进而包含 $\mathcal A\otimes\mathcal B$。因此
\eqref{eq:sigma-finite-product-patching} 的每一项都确实有定义。反过来，所有形如
$E\cap(X_i\times Y_j)$、$E\in\mathcal A\otimes\mathcal B$ 的集合构成块上的 $\sigma$-代数，
并包含每个 $(A\cap X_i)\times(B\cap Y_j)$。由两边的最小性，实际得到迹等式
\[
 (\mathcal A\otimes\mathcal B)|_{X_i\times Y_j}
 =(\mathcal A|_{X_i})\otimes(\mathcal B|_{Y_j}).
\]

最后令 $P_N=(\bigcup_{i\le N}X_i)\times(\bigcup_{j\le N}Y_j)$。这些矩形递增覆盖
$X\times Y$，且任一满足矩形公式的候选测度在 $P_N$ 上都有相同有限质量。可测矩形构成
$\pi$-系统并生成 $\mathcal A\otimes\mathcal B$，故
\cref{theorem:2-2-1-2} 的 $\sigma$-有限唯一性部分给出全局唯一性。
\end{proof}

\begin{example}[Lebesgue 乘积]\label{ex:2-3-2-1}
先在 Borel 定义域上陈述。若 $m_{d,B}$ 表示 $\R^d$ 上的 Borel Lebesgue 测度，则对半开盒
$A\subset\R^n$、$B\subset\R^k$，$m_{n,B}\otimes m_{k,B}$ 与 $m_{n+k,B}$ 在
$A\times B$ 上都取值 $m_{n,B}(A)m_{k,B}(B)$。有理端点开盒构成 Euclidean 拓扑的可数基，而
每个这样的开盒都是递增半开盒的可数并，故半开乘积盒生成 $\Borel(\R^{n+k})$。它们还给出共同
有限耗尽 $[-N,N)^n\times[-N,N)^k$。唯一性因此给出
\[
 m_{n,B}\otimes m_{k,B}=m_{n+k,B}
 \quad\text{在 }\Borel(\R^n)\otimes\Borel(\R^k)
 =\Borel(\R^{n+k})\text{ 上}.
\]
完成这一个 Borel 乘积测度便得到 $m_{n+k}$。这里没有把
$\mathcal L(\R^n)\otimes\mathcal L(\R^k)$ 与高维 Lebesgue 完成化擅自视为同一个
$\sigma$-代数；完成化与乘积定义域的进一步比较留到 Tonelli 理论建立之后。
\end{example}

\begin{example}[乘积平面中的对角线]\label{ex:2-3-2-2}
对角线 $\Delta=\{(x,x):x\in\R\}$ 闭，因而 Borel。固定 $M>0$，沿对角线用至多
$\lceil2M/\varepsilon\rceil+1$ 个边长 $2\varepsilon$ 的方盒覆盖
$\Delta\cap[-M,M]^2$。总面积至多
\[
 4\varepsilon^2\left(\frac{2M}{\varepsilon}+2\right)
 =8M\varepsilon+8\varepsilon^2.
\]
令 $\varepsilon\downarrow0$，该有界部分零测；再对 $M\in\N$ 取可数并，得到
$(m\otimes m)(\Delta)=0$。
\end{example}

\begin{example}[不可数乘积中的两种 $\sigma$-代数]\label{ex:2-3-2-3}
令 $I$ 不可数，$X=\{0,1\}^I$。柱 $\sigma$-代数中的每个集合只依赖某个可数坐标集：生成柱只
依赖有限坐标，而“依赖至多可数坐标”的集合类对补集和可数并封闭。全零点
$0^I$ 的单点集却在乘积拓扑中闭，因为
\[
 \{0^I\}=\bigcap_{i\in I}\{x:x_i=0\}.
\]
它不依赖任何可数坐标集：固定可数 $J\subset I$ 后，在 $I\setminus J$ 的一个坐标改成 $1$，
所得点与全零点在 $J$ 上相同。因此该闭集不在柱 $\sigma$-代数，乘积拓扑的 Borel
$\sigma$-代数严格更大。
\end{example}

\begin{example}[非 $\sigma$-有限时矩形规则不唯一]\label{ex:2-3-2-5}
令 $X=[0,1]$，第一因子取计数测度 $\#$，第二因子取 Lebesgue 测度 $m$。对 Borel
$E\subset X^2$，置
\[
 \lambda(E)=\sum_{x\in X}m(E_x),\qquad
 \eta(E)=m\{x:(x,x)\in E\},
\]
其中不可数非负和定义为所有有限部分和的上确界。对两两不交的 $E_n$，每个截面也不交；交换两个
非负求和次序得到
\[
\begin{aligned}
 \lambda\!\left(\bigsqcup_nE_n\right)
 &=\sum_{x\in X}m\!\left(\bigsqcup_n(E_n)_x\right)
  =\sum_{x\in X}\sum_nm((E_n)_x)\\
 &=\sum_n\sum_{x\in X}m((E_n)_x)
  =\sum_n\lambda(E_n).
\end{aligned}
\]
这里两个迭代和都等于所有有限集合
$F\subset X$ 与有限指标集 $G\subset\N$ 的矩形部分和
$\sum_{x\in F,n\in G}m((E_n)_x)$ 的上确界，故换序不需要预设不可数级数的条件收敛。
$\eta$ 是对角映射
$x\mapsto(x,x)$ 对 $m$ 的推前，因而也是测度。

矩形上
$\lambda(A\times B)=\#A\,m(B)$。而
$\eta(A\times B)=m(A\cap B)$。若 $m(B)=0$，则两式中的附加项为零，且约定
$\#A\,m(B)=0$；若 $m(B)>0$ 而 $A$ 至多可数，则 $m(A\cap B)=0$；若 $m(B)>0$ 且 $A$
不可数，则 $\#A\,m(B)=\infty$，加上不超过 $1$ 的 $m(A\cap B)$ 仍为无穷。因此
$\lambda$ 与 $\lambda+\eta$ 都满足同一矩形
规则。可是
\[
 \lambda(\Delta)=0,\qquad (\lambda+\eta)(\Delta)=1.
\]
计数测度不 $\sigma$-有限，正是唯一性失效的位置。
\end{example}

\begin{proposition}[矩形近似]\label{proposition:2-3-2-3}
若 $E\in\mathcal A\otimes\mathcal B$ 且
$(\mu\otimes\nu)(E)<\infty$，则对每个 $\varepsilon>0$，存在有限个可测矩形的并 $R$ 使
\[
 (\mu\otimes\nu)(E\mathbin\triangle R)<\varepsilon.
\]
\end{proposition}

\begin{proof}
先在有限乘积空间中证明。令 $\mathcal C$ 为可被矩形代数按对称差测度任意逼近的集合类。它含矩形
代数。若 $E_n\uparrow E$ 且 $E_n\in\mathcal C$，从下连续性使
$(\mu\otimes\nu)(E\setminus E_N)$ 对某个 $N$ 小于 $\varepsilon/2$，再逼近 $E_N$。递减情形由有限
总质量和补集化为递增情形：矩形代数含全空间且对补封闭，并且若 $R$ 逼近 $E$，则
\[
 E^c\mathbin\triangle R^c=E\mathbin\triangle R.
\]
故 $\mathcal C$ 对补集封闭，递减列取补后即化为递增列。因此 $\mathcal C$ 是含矩形代数的单调类，集合单调类定理给
$\mathcal C=\mathcal A\otimes\mathcal B$。

一般 $\sigma$-有限情形，先取共同有限矩形耗尽 $Q_n$，使
$(\mu\otimes\nu)(E\setminus Q_n)<\varepsilon/2$，再在 $Q_n$ 内使用有限情形。
\end{proof}

\begin{proposition}[拓扑 Borel 一致性]\label{proposition:2-3-2-4}
若 $X,Y$ 第二可数，则
\[
 \Borel(X\times Y)=\Borel(X)\otimes\Borel(Y).
\]
一般总有右侧包含于左侧。
\end{proposition}

\begin{proof}
坐标投影连续，所以可测矩形为 Borel 集，给出一般包含。第二可数时，取可数基
$(U_i),(V_j)$。每个乘积开集是它所含基矩形 $U_i\times V_j$ 的并；这样的基矩形只有可数多个，
所以开集属于右侧，给出反向包含。
\end{proof}

\begin{figure}[htbp]
\centering
\begin{tikzpicture}[scale=.9]
 \draw[book line] (-3,-1.8) rectangle (3,1.8);
 \fill[BookAccent!8] (-2.4,-1.3) rectangle (-.5,.8);
 \draw[BookAccent,semithick] (-2.4,-1.3) rectangle (-.5,.8);
 \node at (-1.45,-.25) {$A\times B$};
 \draw[BookWarning,thick] (1.2,-1.5) .. controls (2,-.8) and (.4,.2) .. (1.8,1.4);
 \draw[dashed] (1.25,-1.8)--(1.25,1.8) node[above] {$x$};
 \fill[BookWarning] (1.25,-.91) circle (1.5pt);
 \fill[BookWarning] (1.25,.42) circle (1.5pt);
 \node[right] at (1.35,-.2) {$E_x$};
 \draw[book accent arrow] (3.4,0)--(4.7,0) node[midway,above]{扩张};
 \node[book strong node] at (6.1,0) {$\mu\otimes\nu$\\全局测度};
\end{tikzpicture}
\caption{矩形规则先在环上形成预测度；截面把积空间集合重新变成参数化的单空间集合。}
\label{fig:product-sections-extension}
\end{figure}

\begin{remark}
对 $E\in\mathcal A\otimes\mathcal B$，每个截面可测；但函数
$x\mapsto\nu(E_x)$ 的可测性以及
$(\mu\otimes\nu)(E)$ 与其迭代积分的关系要到 Tonelli 定理才建立。不可数乘积中也必须区分柱
$\sigma$-代数与乘积拓扑的 Borel $\sigma$-代数。
\end{remark}

在概率论中，独立随机变量的联合分布由边缘分布的乘积刻画；把加法映射作用于乘积测度会产生卷积。
Kolmogorov 扩张定理则从相容的有限维分布构造随机过程的分布。两因子乘积的存在与矩形唯一性，是这些
后续构造以及 Tonelli 定理的共同起点。

\begin{exercise}
证明可测集合的截面可测，并说明反命题为什么失败。
\end{exercise}
\begin{exercise}
仅用矩形唯一性证明 $m_n\otimes m_k=m_{n+k}$。
\end{exercise}
\begin{exercise}
证明第二可数空间的乘积 Borel 等于 Borel 乘积。
\end{exercise}
\begin{exercise}
逐项核对非 $\sigma$-有限例中 $\lambda,\eta$ 的可数可加性及所有矩形情形。
\end{exercise}

乘积测度把多坐标观测组合起来；在实线上，次序还允许用半直线质量把整份测度压缩成一个单调函数。

\subsection{分布函数与 Lebesgue--Stieltjes 测度}\label{subsec:2-3-3}
\index{Lebesgue--Stieltjes 测度}\index{分布函数}

实线有一个其他空间没有的便利：集合可以按从左到右的次序累积。有限测度 $\mu$ 的半直线质量
$F(x)=\mu(( -\infty,x])$ 是单调函数。原子表现为跳跃，没有质量的开区间表现为平台。反过来，
若给定右连续非减函数，区间增量 $F(b)-F(a)$ 应当生成测度；关键在于证明这些有限增量对可数分解
连续。

历史上，Stieltjes 的这类增量构造与连分数及矩问题相伴发展；在这里，它的作用是把单调函数编码成
测度。第~4~章引入有界变差与 Riemann--Stieltjes 积分后，我们再比较“由增量生成测度”与“对函数
作 Stieltjes 积分”这两个方向。

\begin{definition}[分布函数]\label{def:2-3-3-1}
实线上有限 Borel 测度 $\mu$ 的分布函数为
\[
 F_\mu(x)=\mu(( -\infty,x]).
\]
若 $\mu$ 是概率，则 $F_\mu(-\infty)=0$、$F_\mu(+\infty)=1$。
\end{definition}

\begin{definition}[Stieltjes 增量]\label{def:2-3-3-2}
设 $F:\R\to\R$ 右连续且非减。在半开区间上置
\[
 \mu_F((a,b])=F(b)-F(a),\qquad a<b.
\]
无限端点只在相应单调极限存在且不产生 $\infty-\infty$ 时使用。下文假设
$F(-\infty)=\lim_{x\to-\infty}F(x)$ 有限。
\end{definition}

\begin{definition}[跳跃]\label{def:2-3-3-3}
非减函数在 $x$ 的跳跃为
\[
 \Delta F(x)=F(x)-F(x-),\qquad F(x-)=\lim_{t\uparrow x}F(t).
\]
\end{definition}

\begin{proposition}[单调函数的跳跃至多可数]\label{proposition:2-3-3-1}
实线上的单调函数只有至多可数多个跳跃点。
\end{proposition}

\begin{proof}
以非减函数为例。固定 $M,n\in\N$，在 $[-M,M]$ 中跳跃至少 $1/n$ 的不同点至多
\[
 n\bigl(F(M)-F(-M-)\bigr)+1
\]
个。事实上，若这样的点为 $-M\le x_1<\cdots<x_r\le M$，则对任意
$\varepsilon>0$，忽略最左边的一个点，并对 $j=1,\ldots,r-1$ 选取
$t_j\in(x_j,x_{j+1})$，使
\[
 x_1<t_1<x_2<t_2<x_3<\cdots<t_{r-1}<x_r
\]
并且 $F(x_{j+1})-F(t_j)>1/n-\varepsilon/r$；这里用的是
$F(t)\to F(x_{j+1}-)$ 当 $t\uparrow x_{j+1}$。这些实轴区间可取为两两不交，单调性于是给
\[
 \sum_{j=1}^{r-1}\bigl(F(x_{j+1})-F(t_j)\bigr)
 \le F(M)-F(-M),
\]
令 $\varepsilon\downarrow0$，得到
$(r-1)/n\le F(M)-F(-M)\le F(M)-F(-M-)$，从而得到所写的上界。
所有正跳跃点是这些有限集合对 $M,n$ 的可数并。非增情形对 $-F$ 使用同一论证。
\end{proof}

\begin{lemma}[Stieltjes 环预测度从上连续]\label{lem:stieltjes-continuity-at-empty}
在有界半开区间及其有限不交并组成的环上，以区间增量定义的有限可加函数 $\mu_0$ 满足：若
$E_n\downarrow\varnothing$，则 $\mu_0(E_n)\downarrow0$。因此 $\mu_0$ 是预测度。
\end{lemma}

\begin{proof}
若 $E=\bigsqcup_{i=1}^r(a_i,b_i]$，先置
\[
 \mu_0(E)=\sum_{i=1}^r\bigl(F(b_i)-F(a_i)\bigr).
\]
若 $E$ 还有另一有限不交区间表示，收集两种表示的全部端点并细分；在每个分量上恒等式
\[
 F(b)-F(a)=\sum_{l=1}^N\bigl(F(c_l)-F(c_{l-1})\bigr),
 \qquad a=c_0<\cdots<c_N=b,
\]
说明两种总和都等于同一组小区间增量之和。因此 $\mu_0$ 良定义，且同一共同细分立即给出有限
可加性。设
$E_n\downarrow\varnothing$，并反设
$\inf_n\mu_0(E_n)=\eta>0$。把每个 $E_n$ 写成有限个不交区间 $(a,b]$。由 $F$ 在每个左端点
右连续，若第 $n$ 轮有 $r_n$ 个非空分量 $(a_{n,l},b_{n,l}]$，可在每个分量中选
$c_{n,l}\in(a_{n,l},b_{n,l})$，使
\[
 F(c_{n,l})-F(a_{n,l})<\frac{\eta2^{-n-2}}{r_n}.
\]
故各左端薄条 $(a_{n,l},c_{n,l}]$ 的增量总和小于 $\eta2^{-n-2}$。令 $K_n$ 为相应闭区间
$[c_{n,l},b_{n,l}]$ 的有限并，令 $H_n$ 为这些薄条的有限并。
于是 $K_n\subset E_n$ 是紧集，$H_n$ 属于该环，且 $H_n$ 覆盖 $E_n\setminus K_n$，并有
$\mu_0(H_n)<\eta2^{-n-2}$。这里没有给未必属于半开区间环的差集
$E_n\setminus K_n$ 直接赋值。

令 $L_n=\bigcap_{j=1}^nK_j$。若某个 $L_n$ 为空，则每个 $x\in E_n$ 至少被前 $n$ 轮中的一个
薄条删去，所以 $E_n\subset\bigcup_{j\le n}H_j$。有限次可加与次可加给
\[
 \mu_0(E_n)\le\sum_{j=1}^n\mu_0(H_j)<\eta/2,
\]
其中第一个不等式只用环上的非负有限可加性。确切地，令
\[
 H'_1=H_1,
 \qquad H'_j=H_j\setminus\bigcup_{i<j}H_i\quad(2\le j\le n).
\]
则 $(E_n\cap H'_j)_{j\le n}$ 两两不交并覆盖 $E_n$，所以
\[
 \mu_0(E_n)=\sum_{j=1}^n\mu_0(E_n\cap H'_j)
 \le\sum_{j=1}^n\mu_0(H'_j)
 \le\sum_{j=1}^n\mu_0(H_j).
\]
与定义 $\mu_0(E_n)\ge\eta$ 矛盾。因此每个 $L_n$ 非空。它们是递减非空紧集，并包含于有界紧集
$L_1$，有限交性质给
$\bigcap_nL_n\ne\varnothing$；但该交包含于 $\bigcap_nE_n=\varnothing$，再次矛盾。

最后说明预测度。若环中 $A=\bigsqcup_{k\ge1}A_k$，令
$R_N=A\setminus\bigsqcup_{k\le N}A_k$。则 $R_N\downarrow\varnothing$，有限可加性给
$\mu_0(A)=\sum_{k\le N}\mu_0(A_k)+\mu_0(R_N)$。令 $N\to\infty$ 即得可数可加。
\end{proof}

\begin{theorem}[Lebesgue--Stieltjes 构造]\label{theorem:2-3-3-2}
每个右连续非减函数 $F:\R\to\R$，若 $F(-\infty)$ 有限，则唯一确定局部有限 Borel 测度
$\mu_F$，满足
\[
 \mu_F((a,b])=F(b)-F(a),\qquad
 \mu_F(( -\infty,b])=F(b)-F(-\infty).
\]
给 $F$ 加常数不改变 $\mu_F$。
\end{theorem}

\begin{proof}
上一引理给出半开有界区间环上的预测度。环版 \Caratheodory{} 扩张产生 Borel 测度；每个有界
区间的质量有限，所以它局部有限；每个紧集包含于有界区间，故紧集有限。实线第二可数且局部紧，
\cref{theorem:2-2-2-2} 还说明该测度为 Radon。半开区间在耗尽 $(-n,n]$ 上形成共同有限的 $\pi$-系统，
\cref{theorem:2-2-1-2} 给出唯一性。对
$(a,b]$ 令 $a\to-\infty$，测度从下连续给第二个公式。所有增量只含函数值之差，故加常数不改变
测度。
\end{proof}

\begin{corollary}[紧区间常值延拓版]\label{cor:stieltjes-compact-interval}
设 $G:[a,b]\to\R$ 右连续且非减，给定 $\alpha\ge0$。定义
\[
 F(x)=
 \begin{cases}
 G(a)-\alpha,&x<a,\\
 G(x),&a\le x\le b,\\
 G(b),&x>b.
 \end{cases}
\]
则 $F$ 生成支撑包含于 $[a,b]$ 的有限测度，并且
\[
 \mu_F(\{a\})=\alpha,
 \qquad \mu_F((s,t])=G(t)-G(s)\quad(a\le s<t\le b).
\]
取 $\alpha=0$ 就是两端常值延拓且左端无原子；若要保留指定的左端原子，必须把左侧常值降低
$\alpha$。这个端点约定在以后把区间上的 BV 函数变成 Stieltjes 测度时保持不变。
\end{corollary}

\begin{proof}
分段函数右连续且非减，并在两端外常值，故上一构造适用。由测度从上连续和构造定理的区间公式，
\[
 \mu_F(\{a\})
 =\lim_{n\to\infty}\mu_F((a-1/n,a])
 =\lim_{n\to\infty}\bigl(F(a)-F(a-1/n)\bigr)=\alpha.
\]
内部半开区间公式就是 Stieltjes 增量。补集
$\R\setminus[a,b]$ 的两部分可逐项计算：
\[
 \mu_F(( -\infty,a))
 =\lim_{n\to\infty}\bigl(F(a-1/n)-F(-\infty)\bigr)=0,
\]
并且由从下连续性
\[
 \mu_F((b,\infty))
 =\lim_{n\to\infty}\mu_F((b,n])
 =\lim_{n\to\infty}\bigl(F(n)-F(b)\bigr)=0.
\]
所以测度集中于 $[a,b]$。
\end{proof}

\begin{proposition}[跳跃就是原子]\label{prop:stieltjes-jump-atom}
对 Stieltjes 测度，
\[
 \mu_F(\{x\})=\Delta F(x).
\]
\end{proposition}

\begin{proof}
集合 $(x-1/n,x]$ 递减至 $\{x\}$，且第一项质量有限。测度从上连续给
\[
 \mu_F(\{x\})=\lim_n\bigl(F(x)-F(x-1/n)\bigr)=F(x)-F(x-).
\]
\end{proof}

\begin{example}[纯原子分布]\label{ex:2-3-3-1}
若 $\mu=\sum_kp_k\delta_{x_k}$、$p_k>0$ 且 $\sum_kp_k<\infty$，则
\[
 F_\mu(x)=\sum_{k:x_k\le x}p_k.
\]
点 $x$ 的跳跃为所有满足 $x_k=x$ 的权重之和。若把有理数枚举为 $(q_k)$ 并取
$p_k=2^{-k}$，跳跃点在每个区间稠密，$F$ 没有非平凡平台，支撑为整条实线。纯原子因此不等于
支撑离散。
\end{example}

\begin{example}[均匀分布]\label{ex:2-3-3-2}
$[0,1]$ 上均匀概率的分布函数是
\[
 F(x)=\begin{cases}0,&x<0,\\x,&0\le x\le1,\\1,&x>1.
 \end{cases}
\]
它连续，且除两个端点外导数为 $\1_{(0,1)}$；导数不是在整条实线上恒为 $1$。
\end{example}

\begin{example}[Cantor 分布]\label{ex:2-3-3-3}
标准 Cantor 函数的逐层构造及连续性、单调性证明明确推迟到
\cref{ex:2-4-3-1}。使用那里将证明的性质：$C$ 连续非减，从 $0$ 增长到 $1$，并在 Cantor 集补集的每个三分开区间上
恒定。它的 Stieltjes 测度没有原子，因为 $C$ 没有跳跃；补集是可数个恒定区间之并，每个区间
质量为零，所以全部质量集中在零 Lebesgue 测度的 Cantor 集上。$C$ 在该补集上局部恒定，因而
$C'=0$ 在补集每点成立；Cantor 集零测，所以导数几乎处处为零，而总增量仍为一。导数遗漏了这份
奇异连续质量；相应分解在后章建立。
\end{example}

\begin{theorem}[测度—分布函数对应]\label{theorem:2-3-3-3}
有限 Borel 测度 $\mu$ 与满足右连续、非减、有界且 $F(-\infty)=0$ 的函数一一对应；总质量为
$F(+\infty)$。
\end{theorem}

\begin{proof}
给定有限 $\mu$，半直线嵌套给 $F_\mu$ 非减。若 $x_j\downarrow x$，则
$(-\infty,x_j]\downarrow(-\infty,x]$；有限测度从上连续给右连续。令
$x\to-\infty$、$x\to+\infty$，分别用从上连续于空集和从下连续于全空间，得到端点值
$0,\mu(\R)$。区间差给
\[
 \mu((a,b])=F_\mu(b)-F_\mu(a).
\]
反过来，对规范化 $F$ 使用 Stieltjes 构造。所得测度满足
$\mu_F(( -\infty,b])=F(b)-F(-\infty)=F(b)$，所以其分布函数就是 $F$。若两个有限 Borel 测度
有同一分布函数，则它们在所有 $(a,b]$ 上相等，而且由从下连续性
\[
 \mu(\R)=\lim_{n\to\infty}F(n)=\nu(\R).
\]
半开区间生成 $\Borel(\R)$，有限测度唯一性因而给出 $\mu=\nu$。所以两次操作互逆。
\end{proof}

\begin{proposition}[连续点上的分布函数收敛]\label{proposition:2-3-3-4}
设 $\mu_n,\mu$ 是 $\R$ 上的有限正 Borel 测度，分布函数分别为 $F_n,F$。若
$F_n(x)\to F(x)$ 对 $F$ 的每个连续点成立，则
\[
 \int f\,\dd\mu_n\longrightarrow\int f\,\dd\mu
 \qquad(f\in C_c(\R)).
\]
若还有 $\mu_n(\R)\to\mu(\R)$，则同一结论对每个 $f\in C_b(\R)$ 成立。因此，总质量条件
对紧支撑测试可以删除，但这类测试可能看不见逃向无穷的质量。
\end{proposition}

证明见\cref{corollary:7-3-2-finite-measures}。前面的 \(\delta_n\) 说明，只有紧支撑测试函数时，总质量可能逸向
无穷。

\begin{proposition}[平台读出支撑]\label{proposition:2-3-3-5}
对有限 Borel 测度 $\mu$，
\[
 \R\setminus\supp\mu
 =\bigcup\{(a,b):F_\mu\text{ 在 }(a,b)\text{ 上恒定}\}.
\]
等价地，$x\in\supp\mu$ 当且仅当 $F_\mu$ 在 $x$ 的每个开邻域内都有正增量。
\end{proposition}

\begin{proof}
若 $F$ 在 $(a,b)$ 恒定，则每个有理端点子区间 $(c,d]\subset(a,b)$ 的质量
$F(d)-F(c)$ 为零。这些子区间可数地覆盖 $(a,b)$，故整段零测，其中每点都不在支撑中。
反过来，若 $x\notin\supp\mu$，存在开区间 $(a,b)\ni x$ 且 $\mu((a,b))=0$。对
$a<c<d<b$，有 $F(d)-F(c)=\mu((c,d])=0$，所以 $F$ 在该区间恒定。
\end{proof}

\begin{figure}[htbp]
\centering
\begin{tikzpicture}[x=1.1cm,y=2.2cm]
 \draw[->] (-.3,0)--(6.3,0) node[right]{$x$};
 \draw[->] (0,-.08)--(0,1.25) node[above]{$F(x)$};
 \draw[BookAccent,thick] (0,.08)--(1.3,.08);
 \draw[BookAccent,thick] (1.3,.08)--(2.8,.55);
 \draw[BookAccent,thick] (2.8,.55)--(3.6,.55);
 \draw[BookAccent,thick] (3.6,.78)--(4.5,.78);
 \draw[BookAccent,thick] (4.5,.78).. controls (5,.9) and (5.4,1.08)..(6,1.12);
 \draw[BookWarning,dashed] (3.6,.55)--(3.6,.78);
 \node[below] at (.65,0){零质量平台};
 \node[above] at (2,.36){连续增长};
 \node[right] at (3.6,.67){原子};
 \node[above] at (5.2,.95){弥散增长};
\end{tikzpicture}
\caption{分布函数中的平台、连续增长和跳跃分别编码零质量区、弥散质量与原子。}
\label{fig:stieltjes-platforms-jumps}
\end{figure}

\begin{remark}
全书固定右连续函数和 $(a,b]$ 增量；换成左连续函数时半开方向也必须同时更换。$F$ 与 $F+C$
生成同一测度，只有规定 $F(-\infty)=0$ 才与有限测度一一对应。几乎处处导数不能恢复 Cantor
分布的奇异连续部分。
\end{remark}

分布函数还产生分位数：对给定概率水平寻找最小的 \(x\)，使 \(F(x)\) 达到该水平。它也把矩问题变成
“若干积分是否足以识别测度”的问题；Cantor 例说明，几乎处处导数不能恢复全部测度。分位数、矩与
弱收敛将在积分理论建立后继续讨论。

\begin{exercise}
证明单调函数的跳跃点至多可数，并给出跳跃点稠密的纯原子分布。
\end{exercise}
\begin{exercise}
计算有限离散分布和一个离散—均匀混合分布的分布函数。
\end{exercise}
\begin{exercise}
从给定右连续阶梯函数构造测度，并核对所有半开区间增量。
\end{exercise}
\begin{exercise}
用从上连续性证明 $\mu_F(\{x\})=F(x)-F(x-)$。
\end{exercise}

至此，测度已经能够从有限集合数据被构造、由生成类识别，并沿映射和乘积传输。下一节回到
Euclidean 几何，建立第 3 章换元公式所需的局部线性模型和零集控制。

\section{\texorpdfstring{$\R^n$}{R n} 的局部线性化与体积几何}\label{sec:02-04}

\subsection{\Frechet{} 微分、链式法则与逆映射}\label{subsec:2-4-1}
\index{Frechet@\Frechet{} 微分}\index{逆函数定理}\index{隐函数定理}

一元函数在一点的导数是一个数；在多元问题中，候选者变成了线性映射。真正需要小心的不是把
偏导数排成矩阵，而是矩阵是否同时控制所有趋近方向。下面这个函数把三种常被混在一起的说法
彻底分开。

\begin{example}[方向导数齐全，函数却不连续]\label{ex:2-4-1-1}
定义
\[
 f(0,0)=0,\qquad
 f(x,y)=\frac{x^3y}{x^6+y^2}\quad ((x,y)\ne(0,0)).
\]
固定 $(a,b)\in\R^2$。若 $b\ne0$，则
\[
 f(ta,tb)=\frac{t^2a^3b}{t^4a^6+b^2},\qquad
 \frac{f(ta,tb)-f(0,0)}{t}\longrightarrow0;
\]
若 $b=0$，上述商恒为零。因此 $f$ 在原点沿每条直线的方向导数都存在且等于零，特别地两个
偏导数均为零。可是沿 $y=x^3$ 有
\[
 f(x,x^3)=\frac{x^6}{2x^6}=\frac12\qquad(x\ne0),
\]
所以 $f$ 在原点不连续。逐方向的极限允许误差随方向改变；一个统一的一阶近似不能允许这种逃逸。
\end{example}

三个具体问题说明统一线性近似所增加的内容。连续偏导能否把坐标方向的资料
合并成同一个线性映射？极坐标映射
$\Phi(r,\theta)=(r\cos\theta,r\sin\theta)$ 在何处有局部逆，而周期性为何仍阻止全局反演？
圆周方程 $x^2+y^2-1=0$ 在何处能解成 $y=g(x)$，又为何在左右端点必须交换坐标？
链式法则、逆函数定理与隐函数定理将依次回答这些问题。

\begin{definition}[\Frechet{} 微分]\label{def:2-4-1-1}
设 $U\subset\R^n$ 为开集，$f:U\to\R^m$，且 $x\in U$。若存在线性映射
$A:\R^n\to\R^m$ 使
\[
 \lim_{\substack{h\to0\\x+h\in U}}
 \frac{|f(x+h)-f(x)-Ah|}{|h|}=0,
 \tag{2.4.1}\label{eq:frechet-definition}
\]
则称 $f$ 在 $x$ \emph{\Frechet{} 可微}，并写 $Df(x)=A$。在标准基下，$Df(x)$ 的矩阵称为
Jacobian 矩阵。
\end{definition}

\begin{lemma}[微分的唯一性与连续性]\label{lemma:2-4-1-derivative-unique-continuous}
满足 \eqref{eq:frechet-definition} 的线性映射至多一个；而且，函数在一点 \Frechet{} 可微便在该点
连续。
\end{lemma}

\begin{proof}
若 $A,B$ 都满足定义，则对每个单位向量 $v$ 以及充分小的非零 $t$，点 $x+tv$ 属于 $U$，并且
\[
 |(A-B)v|
 \le \frac{|f(x+tv)-f(x)-Atv|}{|t|}
    +\frac{|f(x+tv)-f(x)-Btv|}{|t|}.
\]
令 $t\to0$ 得 $(A-B)v=0$。任意非零向量都是单位向量的数乘，故 $A=B$。

再将余项记作 $r(h)=f(x+h)-f(x)-Df(x)h$。由
\[
 |f(x+h)-f(x)|\le \|Df(x)\|\,|h|+|r(h)|
\]
及 $|r(h)|/|h|\to0$，右端随 $h\to0$ 趋于零。
\end{proof}

\begin{definition}[$C^1$ 映射]\label{def:2-4-1-2}
若 $f$ 在开集 $U$ 的每一点 \Frechet{} 可微，且 $x\mapsto Df(x)$ 关于算子范数连续，则写
$f\in C^1(U,\R^m)$。
\end{definition}

\begin{definition}[局部微分同胚]\label{def:2-4-1-3}
若存在 $x$ 的开邻域 $V$ 与 $f(x)$ 的开邻域 $W$，使得 $f:V\to W$ 双射且 $f$、$f^{-1}$
均为 $C^1$，则称 $f$ 在 $x$ 附近给出一个局部 $C^1$ 微分同胚。
\end{definition}

定义中的余项量词正好足以在复合时传递。

\begin{theorem}[链式法则]\label{theorem:2-4-1-1}
设 $U\subset\R^n$ 与 $V\subset\R^m$ 为开集，$f:U\to V$ 在 $x\in U$ 可微，且
$g:V\to\R^p$ 在 $f(x)$ 可微。则 $g\circ f$ 在 $x$ 可微，且
\[
 D(g\circ f)(x)=Dg(f(x))Df(x).
\]
\end{theorem}

\begin{proof}
记 $A=Df(x)$、$B=Dg(f(x))$，并写
\[
 f(x+h)-f(x)=Ah+r(h),\qquad \frac{|r(h)|}{|h|}\to0.
\]
于是存在 $h=0$ 的一个邻域，使
$|f(x+h)-f(x)|\le(\|A\|+1)|h|$。令 $k=f(x+h)-f(x)$。在 $g$ 的可微性公式中写
\[
 g(f(x)+k)-g(f(x))=Bk+s(k),\qquad \frac{|s(k)|}{|k|}\to0.
\]
当 $k=0$ 时取 $s(k)=0$。合并两式得到
\[
 g(f(x+h))-g(f(x))-BAh=Br(h)+s(k).
\]
第一项除以 $|h|$ 后趋于零；第二项满足
如下无除零问题的估计。给定 $\varepsilon>0$，当 $k$ 足够小时（包括 $k=0$，此时 $s(0)=0$）有
$|s(k)|\le\varepsilon|k|$；而 $|k|\le(\|A\|+1)|h|$。所以
\[
 \frac{|s(k)|}{|h|}\le\varepsilon(\|A\|+1).
\]
先令 $h\to0$，再令 $\varepsilon\downarrow0$，得到该商趋于零。这便是以 $BA$ 为微分的
\eqref{eq:frechet-definition}。
\end{proof}

\begin{proposition}[$C^1$ 判据与均值估计]\label{proposition:2-4-1-2}
设 $U\subset\R^n$ 开，$f:U\to\R^m$ 的全部一阶偏导数存在并连续，则 $f\in C^1(U,\R^m)$，
且 $Df(x)$ 是由偏导数组成的 Jacobian 矩阵。若线段 $[x,y]\subset U$，则
\[
 |f(y)-f(x)|
 \le \sup_{z\in[x,y]}\|Df(z)\|\,|y-x|.
 \tag{2.4.2}\label{eq:vector-mean-estimate}
\]
\end{proposition}

\begin{proof}
固定 $x\in U$。取一个以 $x$ 为中心、闭包包含于 $U$ 的开盒。对足够小的
$h=(h_1,\dots,h_n)$，置
\[
 x^{(j)}=x+(h_1,\dots,h_j,0,\dots,0),\qquad x^{(0)}=x.
\]
对 $f$ 的第 $\ell$ 个分量在第 $j$ 段应用一元均值定理，得到位于该坐标线段上的点
$\xi_{j\ell}$，使
\[
 f_\ell(x^{(j)})-f_\ell(x^{(j-1)})
 =\partial_j f_\ell(\xi_{j\ell})h_j.
\]
所有 $\xi_{j\ell}$ 都随 $h\to0$ 趋于 $x$。将坐标段相加并减去
$\sum_j\partial_jf_\ell(x)h_j$。对充分小的 $r>0$ 定义
\[
 \omega_x(r)=
 \max_{\substack{1\le j\le n\\1\le\ell\le m}}
 \sup_{|z-x|\le r}
 |\partial_jf_\ell(z)-\partial_jf_\ell(x)|.
\]
闭球取在前述盒内；有限多个偏导在 $x$ 连续，所以 $\omega_x(r)\to0$。于是
\[
 \left|f_\ell(x+h)-f_\ell(x)-\sum_{j=1}^n\partial_jf_\ell(x)h_j\right|
 \le \omega_x(|h|)\sum_{j=1}^n|h_j|,
\]
由 $\sum_j|h_j|\le\sqrt n\,|h|$，再对 $m$ 个分量使用 Euclidean 范数，余项至多为
$\sqrt{mn}\,\omega_x(|h|)|h|=o(|h|)$，从而得到
\Frechet{} 余项。偏导连续又说明所得矩阵随 $x$ 连续。

现在设 $[x,y]\subset U$。若 $f(y)=f(x)$，估计成立。否则令
\[
 u=\frac{f(y)-f(x)}{|f(y)-f(x)|},\qquad
 \varphi(t)=\langle f(x+t(y-x)),u\rangle .
\]
标量函数 $\varphi$ 在 $[0,1]$ 上连续、在 $(0,1)$ 上可微。一元均值定理给出
$\theta\in(0,1)$，使
\[
 |f(y)-f(x)|=\varphi(1)-\varphi(0)
 =\langle Df(x+\theta(y-x))(y-x),u\rangle.
\]
Cauchy--Schwarz 不等式随即给出 \eqref{eq:vector-mean-estimate}。
\end{proof}

逆函数定理的存在部分归结为一个完备性论证。把它表述在一般完备度量空间中，可以看清欧氏闭球
只是应用场景，而不是原理本身。

\begin{lemma}[压缩映射原理]\label{lemma:2-4-1-2a}
设 $(X,d)$ 是非空完备度量空间，$T:X\to X$ 满足
\[
 d(Tx,Ty)\le qd(x,y)\qquad(x,y\in X)
\]
其中 $0\le q<1$。则 $T$ 有唯一不动点。特别地，完备度量空间中的非空闭子集，或 Banach
空间中被 $T$ 映入自身的非空闭球，都可以作为这里的 $X$。
\end{lemma}

\begin{proof}
任取 $x_0\in X$ 并递归定义 $x_{k+1}=Tx_k$。归纳得到
$d(x_{k+1},x_k)\le q^kd(x_1,x_0)$。当 $p\ge1$ 时，三角不等式给出
\[
 d(x_{k+p},x_k)
 \le\sum_{j=k}^{k+p-1}d(x_{j+1},x_j)
 \le\frac{q^k}{1-q}d(x_1,x_0).
\]
因此 $(x_k)$ 是 Cauchy 列。完备性给出 $x_k\to x_*\in X$。压缩估计还说明 $T$ 连续，故
$Tx_*=\lim_kTx_k=\lim_kx_{k+1}=x_*$。若 $y_*$ 也是不动点，则
$d(x_*,y_*)\le qd(x_*,y_*)$，从 $q<1$ 得 $x_*=y_*$。

若 $F$ 是完备空间的闭子集，则 $F$ 完备；闭球是闭子集。这证明最后一项说明。
\end{proof}

\begin{example}[可算的非线性反演]\label{ex:2-4-1-4}
考虑
\[
 x+\varepsilon\sin x=y,
 \]
并假设 $|\varepsilon|<1$。对固定的 $y$，令
$T_y(x)=y-\varepsilon\sin x$。区间
$I_y=[y-|\varepsilon|,y+|\varepsilon|]$ 满足
\[
 |T_y(x)-y|\le|\varepsilon|,
\]
所以 $T_y(I_y)\subset I_y$。又由一元均值定理，
\[
 |T_y(x)-T_y(x')|\le|\varepsilon|\,|x-x'|.
\]
任取 $x_0\in I_y$，再迭代 $x_{k+1}=y-\varepsilon\sin x_k$；不变性保证所有 $x_k\in I_y$。
压缩映射原理给出唯一解，并给出误差界
\[
 |x_k-x_*|\le\frac{|\varepsilon|^k}{1-|\varepsilon|}|x_1-x_0|.
\]
闭区间的不变性与小于一的压缩常数缺一不可；前者保证迭代始终有定义，后者保证迭代收敛。
\end{example}

\begin{theorem}[逆函数定理]\label{theorem:2-4-1-3}
设 $U\subset\R^n$ 开，$f\in C^1(U,\R^n)$，且 $Df(x_0)$ 可逆。则存在开邻域
$V\ni x_0$ 与 $W\ni f(x_0)$，使 $f:V\to W$ 为 $C^1$ 微分同胚。对 $x\in V$，
\[
 D(f^{-1})(f(x))=Df(x)^{-1}.
 \tag{2.4.3}\label{eq:inverse-derivative}
\]
\end{theorem}

\begin{proof}
令 $A=Df(x_0)$。经过定义域平移和目标空间中的可逆线性变换，问题化为
\[
 x_0=0,\qquad f(0)=0,\qquad Df(0)=I.
 \tag{2.4.4}\label{eq:inverse-normalization}
\]
确切地说，可对 $u\mapsto A^{-1}(f(x_0+u)-f(x_0))$ 证明结论，再把坐标变换复原。

固定 $q\in(0,1)$。由 $Df$ 在 $0$ 连续，可取 $r>0$，使
$\overline B(0,r)\subset U$ 且
\[
 \|Df(u)-I\|\le q\qquad(|u|\le r).
 \tag{2.4.5}\label{eq:inverse-q-bound}
\]
置 $R(u)=f(u)-u$。闭球是凸集，命题 \ref{proposition:2-4-1-2} 应用于 $R$ 给出
\[
 |R(u)-R(v)|\le q|u-v|\qquad(u,v\in\overline B(0,r)).
 \tag{2.4.6}\label{eq:inverse-R-contraction}
\]

令 $W_0=B(0,(1-q)r)$。对每个 $y\in W_0$，在闭球上定义
\[
 T_y(u)=y-R(u).
\]
由于 $R(0)=0$，对 $|u|\le r$ 有
\[
 |T_y(u)|\le |y|+|R(u)-R(0)|<(1-q)r+qr=r.
\]
所以 $T_y$ 把闭球映入自身；\eqref{eq:inverse-R-contraction} 又表明它是 $q$-压缩。存在唯一
$u\in\overline B(0,r)$ 满足 $T_y(u)=u$，而该等式正是 $f(u)=y$。上面的严格不等式还给出
$u\in B(0,r)$。

另一方面，对 $u,v\in\overline B(0,r)$，
\[
 |f(u)-f(v)|
 =|(u-v)+(R(u)-R(v))|
 \ge(1-q)|u-v|.
 \tag{2.4.7}\label{eq:inverse-lower-lipschitz}
\]
故 $f$ 在闭球上单射。令
\[
 V_0=f^{-1}(W_0)\cap B(0,r).
\]
$V_0$ 是包含 $0$ 的开集，以上存在性和唯一性说明 $f:V_0\to W_0$ 双射，而且
\[
 |f^{-1}(y')-f^{-1}(y)|\le(1-q)^{-1}|y'-y|.
 \tag{2.4.8}\label{eq:inverse-lipschitz}
\]
因此逆映射连续。

还需证明逆映射可微，不能预先对它使用链式法则。取 $y=f(u)\in W_0$，令
$y'=f(u')\to y$。由 \eqref{eq:inverse-lipschitz}，$u'-u=O(|y'-y|)$。在 $u$ 处的可微性给出
\[
 y'-y=Df(u)(u'-u)+\rho(u'-u),
 \qquad \frac{|\rho(h)|}{|h|}\to0.
 \tag{2.4.9}\label{eq:inverse-remainder}
\]
由 \eqref{eq:inverse-q-bound}，
$|Df(u)v|\ge(1-q)|v|$；在有限维空间中这说明 $Df(u)$ 可逆，且
$\|Df(u)^{-1}\|\le(1-q)^{-1}$。因而
\[
 \frac{|u'-u-Df(u)^{-1}(y'-y)|}{|y'-y|}
 \le \frac{1}{1-q}\frac{|\rho(u'-u)|}{|u'-u|}
       \frac{|u'-u|}{|y'-y|}\longrightarrow0.
\]
这证明 \eqref{eq:inverse-derivative}。最后，矩阵求逆在可逆矩阵集上连续：若 $B$ 接近可逆矩阵
$C$，置 $\theta=\|C^{-1}(B-C)\|<1$。压缩映射原理应用于
$v\mapsto C^{-1}w-C^{-1}(B-C)v$ 给出 $B^{-1}w$，并给出局部一致的逆范数界；再由
\[
 \|B^{-1}w\|
 \le\|C^{-1}\|\,\|w\|+\theta\|B^{-1}w\|,
 \qquad
 \|B^{-1}\|\le\frac{\|C^{-1}\|}{1-\theta}.
\]
于是 $B\to C$ 时这些逆范数局部一致有界，再由
\[
 B^{-1}-C^{-1}=B^{-1}(C-B)C^{-1}
\]
得 $\|B^{-1}-C^{-1}\|\to0$。因此
$Df(u)^{-1}$ 随 $u$ 连续，\eqref{eq:inverse-lipschitz} 又给 $u=f^{-1}(y)$ 的连续性，故
$f^{-1}\in C^1(W_0)$。

复原 \eqref{eq:inverse-normalization} 前的两次坐标变换，取
$V=x_0+V_0$ 与 $W=f(x_0)+AW_0$，结论成立。
\end{proof}

\begin{example}[极坐标映射]\label{ex:2-4-1-2}
令
\[
 \Phi(r,\theta)=(r\cos\theta,r\sin\theta).
\]
其导数和行列式为
\[
 D\Phi(r,\theta)=
 \begin{pmatrix}
 \cos\theta&-r\sin\theta\\
 \sin\theta&r\cos\theta
 \end{pmatrix},
 \qquad \det D\Phi(r,\theta)=r.
\]
所以在每个 $r_0>0$ 的点附近，逆函数定理给出局部 $C^1$ 逆。$r=0$ 时行列式消失，所有角度
都被送到原点；即使 $r>0$，若角度不限制在长度小于 $2\pi$ 的区间中，$2\pi$ 周期性也破坏全局
单射。局部结论与全局结论在这里可以直接从公式中区分。
\end{example}

\begin{corollary}[隐函数定理]\label{corollary:2-4-1-4}
设 $F:U\subset\R^p\times\R^q\to\R^q$ 为 $C^1$，
$F(x_0,y_0)=0$，且 $D_yF(x_0,y_0)$ 可逆。则存在 $x_0$ 的开邻域 $X$、$y_0$ 的开邻域
$Y$ 以及唯一的 $C^1$ 映射 $g:X\to Y$，使得
\[
 \{(x,y)\in X\times Y:F(x,y)=0\}
 =\{(x,g(x)):x\in X\}.
\]
并且
\[
 Dg(x)=-[D_yF(x,g(x))]^{-1}D_xF(x,g(x)).
 \tag{2.4.10}\label{eq:implicit-derivative}
\]
\end{corollary}

\begin{proof}
定义
\[
 H(x,y)=(x,F(x,y)).
\]
在 $(x_0,y_0)$ 处，
\[
 DH(x_0,y_0)=
 \begin{pmatrix}
 I_p&0\\
 D_xF(x_0,y_0)&D_yF(x_0,y_0)
 \end{pmatrix}.
\]
给定 $(a,b)\in\R^p\times\R^q$，方程
$DH(x_0,y_0)(u,v)=(a,b)$ 有唯一解
\[
 u=a,\qquad
 v=D_yF(x_0,y_0)^{-1}(b-D_xF(x_0,y_0)a),
\]
故该块矩阵可逆。逆函数定理给出 $H$ 的局部 $C^1$ 逆。因为 $H$ 的第一坐标就是 $x$，逆映射
的邻域可显式缩小如下。设定理给出
$H:V_0\to W_0$ 的微分同胚，且 $(x_0,y_0)\in V_0$、$(x_0,0)\in W_0$。先取开邻域
$X_0,Y$ 使 $X_0\times Y\subset V_0$。由 $H^{-1}$ 在 $(x_0,0)$ 连续，再取开邻域
$X_1\ni x_0$ 与 $Z\ni0$，使
\[
 X_1\times Z\subset W_0,
 \qquad H^{-1}(X_1\times Z)\subset X_0\times Y.
\]
令 $X=X_0\cap X_1$。由 $H$ 的第一坐标恒为输入的第一坐标，逆映射在 $X\times Z$ 上必有形式
\[
 H^{-1}(x,z)=(x,h(x,z)).
\]
令 $g(x)=h(x,0)$，便有 $F(x,g(x))=0$。局部逆的唯一性同时说明，缩小邻域后所有零点都以此
形式出现：若 $(x,y)\in X\times Y$ 且 $F(x,y)=0$，则
\[
 H(x,y)=(x,0)=H(x,g(x));
\]
而两点都在 $V_0$，$H|_{V_0}$ 单射，故 $y=g(x)$。

对恒等式 $F(x,g(x))=0$ 使用链式法则，得到
\[
 D_xF(x,g(x))+D_yF(x,g(x))Dg(x)=0.
\]
连续性保证在邻域缩小后 $D_yF(x,g(x))$ 仍可逆，左乘其逆即得
\eqref{eq:implicit-derivative}。
\end{proof}

\begin{example}[圆周在局部是一张图]\label{ex:2-4-1-3}
取 $F(x,y)=x^2+y^2-1$。在圆周上一点 $(x_0,y_0)$ 若 $y_0\ne0$，则
$D_yF(x_0,y_0)=2y_0\ne0$。隐函数定理给出 $y=g(x)$，并由
\eqref{eq:implicit-derivative} 得
\[
 g'(x)=-\frac{x}{g(x)}.
\]
在上半圆和下半圆中，这分别与 $g(x)=\sqrt{1-x^2}$ 和
$g(x)=-\sqrt{1-x^2}$ 的直接求导一致。若 $y_0=0$，不能在该点把圆周写成 $y$ 关于 $x$ 的图，
但 $D_xF=2x_0\ne0$，交换两个坐标后可以写成 $x$ 关于 $y$ 的图。
\end{example}

\begin{figure}[htbp]
  \centering
  \begin{tikzpicture}[scale=0.92]
    \begin{scope}
      \draw[book line] (-2,-1.4) rectangle (0.2,1.4);
      \draw[BookInk!45] (-1.55,-1.4)--(-1.55,1.4) (-0.7,-1.4)--(-0.7,1.4)
                         (-2,-0.45)--(0.2,-0.45) (-2,0.55)--(0.2,0.55);
      \fill[BookAccent!12,draw=BookAccent,semithick] (-1.55,-0.45) rectangle (-0.7,0.55);
      \node[book label] at (-0.9,-1.72) {定义域中的小盒};
      \node[book point] at (-1.12,0.05) {};
      \node[book label,above] at (-1.12,0.12) {$x_0$};
    \end{scope}
    \draw[book accent arrow] (0.45,0.65) -- node[above,book label] {$Df(x_0)$} (2.85,0.65);
    \draw[book dashed arrow] (0.45,-0.65) to[bend right=12]
      node[below,book label] {$f$} (5.75,-0.65);
    \begin{scope}[xshift=4.6cm]
      \fill[BookWarm!12,draw=BookWarm,semithick]
        (-1.45,-0.62)--(-0.55,-0.9)--(0.05,0.45)--(-0.86,0.72)--cycle;
      \node[book label] at (-0.7,-1.72) {线性平行多面体};
    \end{scope}
    \begin{scope}[xshift=7.5cm]
      \fill[BookAccent!10,draw=BookAccent,semithick]
        (-1.5,-0.58) .. controls (-1.16,-0.88) and (-0.78,-0.93) .. (-0.48,-0.8)
        .. controls (-0.15,-0.35) and (0.05,0.15) .. (0.02,0.48)
        .. controls (-0.42,0.78) and (-0.93,0.78) .. (-1.18,0.66)
        .. controls (-1.42,0.25) and (-1.55,-0.22) .. cycle;
      \node[book label] at (-0.72,-1.72) {非线性像};
    \end{scope}
  \end{tikzpicture}
  \caption{在足够小的尺度上，$f$ 的像由 $Df(x_0)$ 的平行多面体加一个低阶误差控制。逆函数定理
  还需要线性部分可逆以及邻域上一致的小误差。}
  \label{fig:2-4-1-linearization}
\end{figure}

逆函数定理只作局部断言。例如 $t\mapsto t^2$ 在 $t=1$ 处导数非零，却在整条实线上不是单射；
边界点也不在上述开集版本的论域内。另一方面，\eqref{eq:inverse-q-bound}--
\eqref{eq:inverse-lower-lipschitz} 留下了比“存在一个逆”更有用的定量信息。下一小节先把其中的线性
映射单独拿出来，核算它对体积的作用。

在常微分方程中，解对初值的局部流若足够光滑，局部反演可以解释短时间内怎样由终点恢复初值；隐函数
定理则用来追踪由方程定义的约束面。随机微分方程的流还需要概率估计和更强的正则性。

\begin{exercise}[微分的唯一性]\label{exercise:2-4-1-1}
不引用引理 \ref{lemma:2-4-1-derivative-unique-continuous}，直接从定义证明：若
$f(x+h)=f(x)+Ah+o(|h|)=f(x)+Bh+o(|h|)$，则 $A=B$。
\end{exercise}

\begin{exercise}[偏导仍然不够]\label{exercise:2-4-1-2}
定义 $u(0,0)=0$，其余点令 $u(x,y)=x^2y/(x^4+y^2)$。计算原点的两个偏导，并研究
$u(x,x^2)$，据此判断连续性和 \Frechet{} 可微性。
\end{exercise}

\begin{exercise}[局部开映射]\label{exercise:2-4-1-3}
设 $f\in C^1(U,\R^n)$ 且 $Df(x_0)$ 可逆。用定理 \ref{theorem:2-4-1-3} 证明：存在
$x_0$ 的邻域 $V$，使 $f(O)$ 对每个相对开集 $O\subset V$ 都是 $f(V)$ 中的相对开集。
\end{exercise}

\begin{exercise}[隐函数导数]\label{exercise:2-4-1-4}
从 $F(x,g(x))=0$ 出发，逐项写出复合映射的微分，并推导
\eqref{eq:implicit-derivative}。说明为什么矩阵乘法的次序不能交换。
\end{exercise}

\subsection{行列式、定向与线性体积缩放}\label{subsec:2-4-2}
\index{行列式}\index{定向}\index{线性换元}

线性映射对小盒的作用可以拆成三种实验：交换坐标，沿某个坐标伸缩，以及把一个坐标的倍数加到
另一个坐标上。前两种的体积变化可以从边长读出；第三种把矩形推成平行多面体，边长相乘的公式不再
直接适用。行列式为三类变换提供同一个乘法因子，但测度公式还必须回答两个问题：剪切为什么
严格保持体积？退化映射为什么把任意集合映入零测集？

行列式最初也可被看作组织线性消元与定向的代数量；在多元微积分中，它获得了新的几何角色：局部
线性化的行列式正是 Jacobian 体积因子。下面先在全局线性映射上严格建立这一解释。

\begin{definition}[行列式的公理刻画]\label{def:2-4-2-1}
在 $(\R^n)^n$ 上唯一满足下列三条性质的函数记作 $\det$：它对每个变量线性；交换两个变量使符号
反转；并且 $\det(e_1,\dots,e_n)=1$。若线性算子 $A$ 的列向量为 $a_1,\dots,a_n$，则
$\det A:=\det(a_1,\dots,a_n)$。
\end{definition}

把 $v_j=\sum_i v_{ij}e_i$ 代入多重线性展开。含有两个相同基向量的项因交换后等于自身的相反数，
故为零；其余项恰由排列 $\sigma\in S_n$ 编号。交换次数决定符号，于是三条性质迫使
\[
 \det(v_1,\dots,v_n)
 =\sum_{\sigma\in S_n}\sgn(\sigma)
   v_{\sigma(1),1}\cdots v_{\sigma(n),n}.
 \tag{2.4.11}\label{eq:det-leibniz-derived}
\]
右端反过来满足三条性质，所以存在性和唯一性都包含在这个有限展开中。

\begin{definition}[定向]\label{def:2-4-2-2}
两个有序基称为同向，如果从一个基到另一个基的过渡矩阵行列式为正。这个关系把全部有序基分成
两个\emph{定向类}。若一个 $C^1$ 微分同胚在所讨论的连通开集上满足 $\det Df>0$，称它保定向；
若 $\det Df<0$，称它反定向。普通 Lebesgue 体积不记录这两个符号，只记录绝对值。
\end{definition}

\begin{proposition}[行列式乘法性]\label{proposition:2-4-2-1}
对任意 $n\times n$ 实矩阵 $A,B$，
\[
 \det(AB)=\det A\det B.
\]
\end{proposition}

\begin{proof}
固定 $A$，考虑
\[
 \Lambda_A(v_1,\dots,v_n)=\det(Av_1,\dots,Av_n).
\]
$\Lambda_A$ 是交替 $n$-线性函数。由定义 \ref{def:2-4-2-1} 的唯一性，任意交替
$n$-线性函数等于其在 $(e_1,\dots,e_n)$ 处的值乘以 $\det$，故
\[
 \Lambda_A(v_1,\dots,v_n)=\det(Ae_1,\dots,Ae_n)\det(v_1,\dots,v_n)
 =\det A\det(v_1,\dots,v_n).
\]
取 $v_1,\dots,v_n$ 为 $B$ 的列向量，左端是 $\det(AB)$，右端是 $\det A\det B$。
\end{proof}

\begin{definition}[单位盒的线性体积比]\label{def:2-4-2-3}
令 $Q_0=[0,1)^n$。对可逆线性映射 $A$，先定义
\[
 J_m(A)=m(AQ_0).
\]
下面的集合公式将证明 $J_m(A)=|\det A|$，并说明以任意正体积盒代替 $Q_0$ 都给出同一体积比。
\end{definition}

首先处理 Gaussian 消元中的三类基本映射。剪切部分的细网格估计是证明中唯一不再把盒送成盒的
地方。

\begin{lemma}[基本线性映射的体积]\label{lemma:2-4-2-elementary-volume}
设 $E\subset\R^n$ Lebesgue 可测。
\begin{enumerate}
  \item 坐标置换保持 $m(E)$；
  \item 若 $D_\lambda$ 只把第 $i$ 个坐标乘以 $\lambda\ne0$，则
        $m(D_\lambda E)=|\lambda|m(E)$；
  \item 若 $S_c$ 把第 $i$ 个坐标替换成 $x_i+cx_j$（$i\ne j$），其余坐标不变，则
        $m(S_cE)=m(E)$。
\end{enumerate}
这些映射还把 Lebesgue 零集的任意子集送到 Lebesgue 零集。
\end{lemma}

\begin{proof}
坐标置换把半开盒送到同体积半开盒。对任意集合的盒覆盖逐个施加映射，再对逆映射重复同一
操作，可得外测度相等。它们是 Borel 双射，因而把 Borel 集送到 Borel 集；结合完成化的 Borel
代表，便得到第一项以及关于零集子集的断言。

对 $D_\lambda$，在 Borel 集上定义 $\nu(B)=m(D_\lambda B)$。由于 $D_\lambda$ 是 Borel 双射，
$\nu$ 是测度。坐标超平面与 $[-M,M]^n$ 的交可放进一个其余边长为 $2M$、法向厚度任意小的盒
中；先令厚度趋于零，再令 $M\in\N$ 取可数并，得到坐标超平面零测。因此当
$\lambda<0$ 时，端点开闭方向的改变不影响测度。于是对任意半开盒 $Q$，
\[
 \nu(Q)=|\lambda|m(Q).
\]
$\nu$ 与 $|\lambda|m$ 都在整数盒耗尽上有限，测度唯一性定理
\ref{theorem:2-2-1-2} 说明二者在全部 Borel 集上相同。若 $E$ Lebesgue 可测，取 Borel 集
$B,N$ 使 $E\mathbin\triangle B\subset N$ 且 $m(N)=0$。Borel 公式给
$m(D_\lambda N)=0$，从
\[
 D_\lambda E\mathbin\triangle D_\lambda B\subset D_\lambda N
\]
得到 $D_\lambda E$ 可测及所需公式。相同论证处理零集的任意子集。

现在考察剪切。设
$Q=\prod_{\ell=1}^n[a_\ell,b_\ell)$，边长为 $d_\ell=b_\ell-a_\ell$。把第 $j$ 个
坐标区间等分成 $N$ 段，长度 $\delta=d_j/N$；若 $Q=\varnothing$，结论平凡，故以下只需考虑
所有 $d_\ell>0$ 的情形。在每个所得薄盒中，$cx_j$ 的振幅不超过
$|c|\delta$。给定 $\varepsilon>0$，该薄盒的像因而包含于一个轴平行盒；后者在第 $i$ 个方向长度为
$d_i+|c|\delta+\varepsilon$，在第 $j$ 个方向长度为 $\delta$，其他方向长度为 $d_\ell$。
这里的任意小外扩处理了半开端点可能不能由同长度半开区间覆盖的问题。于是这 $N$ 个盒给出
\[
 \begin{aligned}
 m^*(S_cQ)
 &\le N(d_i+|c|\delta+\varepsilon)\delta
       \prod_{\ell\ne i,j}d_\ell\\
 &=m(Q)+(|c|\delta+\varepsilon)d_j\prod_{\ell\ne i,j}d_\ell.
 \end{aligned}
\]
先令 $\varepsilon\downarrow0$，再令 $N\to\infty$，得到 $m^*(S_cQ)\le m(Q)$。若任意 $A\subset\R^n$ 被盒列 $(Q_k)$ 覆盖，
则
\[
 m^*(S_cA)\le\sum_km^*(S_cQ_k)\le\sum_km(Q_k).
\]
对所有覆盖取下确界，$m^*(S_cA)\le m^*(A)$。将同一结论用于逆剪切 $S_{-c}$ 和集合 $S_cA$，
得到反向不等式，故
\[
 m^*(S_cA)=m^*(A)\qquad(A\subset\R^n).
 \tag{2.4.12}\label{eq:shear-outer-measure}
\]
若 $E\mathbin\triangle B$ 包含于 Borel 零集 $N$，则 $S_cB$ 是 Borel 集，
\eqref{eq:shear-outer-measure} 说明 $S_cN$ 零测，于是 $S_cE$ 可测且测度不变。零集子集的结论
也直接来自 \eqref{eq:shear-outer-measure}。
\end{proof}

\begin{example}[剪切]\label{ex:2-4-2-1}
在 $\R^2$ 中令 $A(x,y)=(x+cy,y)$。单位方形被送到由四个顶点
\[
 (0,0),\qquad(1,0),\qquad(c,1),\qquad(1+c,1)
\]
张成的平行四边形。上一证明把它按 $N$ 条水平薄带覆盖，总面积
至多 $1+|c|/N$；逆剪切给反向估计，所以其面积等于 $1$。矩阵
\[
 \begin{pmatrix}1&c\\0&1\end{pmatrix}
\]
的行列式也等于 $1$。
\end{example}

\begin{example}[反射]\label{ex:2-4-2-2}
映射
\[
 A(x_1,x_2,\dots,x_n)=(-x_1,x_2,\dots,x_n)
\]
把每个盒的第一坐标区间翻转，体积不变；其矩阵行列式为 $-1$。这项计算说明无符号体积公式若写成
$\det A$，会把非负测度错误地变成负数。
\end{example}

退化矩阵需要另一项覆盖计算。它不能由“把伸缩因子取成零”直接推出，因为尚未知道测度公式对极限
稳定。

\begin{lemma}[低维线性子空间零测]\label{lemma:2-4-2-lower-dimensional-null}
若 $V\subset\R^n$ 是维数 $k<n$ 的线性子空间，则 $m_n(V)=0$。因此 $V$ 的每个子集都是
Lebesgue 可测的零集。
\end{lemma}

\begin{proof}
\(k=0\) 时 \(V=\{0\}\)，而单点可放入边长任意小的 \(n\) 维盒，结论成立。以下设
\(1\le k<n\)。
从坐标投影中选择 $k$ 个坐标，使相应投影 $\pi:V\to\R^k$ 为单射；这是因为 $V$ 的一组基所成
$n\times k$ 矩阵有一个非零的 $k\times k$ 子式。有限维线性代数给出常数 $C\ge1$，使
\[
 |v|\le C|\pi v|\qquad(v\in V).
 \tag{2.4.13}\label{eq:subspace-coordinate-bound}
\]
固定 $M>0$，考察 $V_M=V\cap[-M,M]^n$。集合 $\pi(V_M)$ 位于 $[-M,M]^k$ 中。用至多
$(2M/\varepsilon+2)^k$ 个边长 $\varepsilon$ 的 $k$ 维盒覆盖这个参数盒。由
\eqref{eq:subspace-coordinate-bound}，每个参数小盒对应的 $V_M$ 部分直径至多
$C\sqrt{k}\,\varepsilon$，因而可放进一个边长
$2C\sqrt{k}\,\varepsilon$ 的 $n$ 维坐标盒。于是存在只依赖 $M,n,C$ 的常数 $C_M$，使
\[
\begin{aligned}
 m_n^*(V_M)
 &\le\left(\frac{2M}{\varepsilon}+2\right)^k
       (2C\sqrt{k}\,\varepsilon)^n\\
 &=(2M+2\varepsilon)^k(2C\sqrt{k})^n
       \varepsilon^{n-k}.
\end{aligned}
\]
令 $\varepsilon\downarrow0$ 得 $m_n^*(V_M)=0$。最后
$V=\bigcup_{M\in\N}V_M$，外测度的可数次可加性给 $m_n^*(V)=0$；Lebesgue 测度的完备性处理
所有子集。
\end{proof}

\begin{example}[奇异映射]\label{ex:2-4-2-3}
若 $A:\R^n\to\R^n$ 的秩为 $k<n$，则 $A(\R^n)$ 是一个 $k$ 维线性子空间。对任意
$E\subset\R^n$，
\[
 AE\subset A(\R^n),
\]
所以上一引理给 $m_n^*(AE)=0$。这也包括 $E$ 无界或不可测的情形。
\end{example}

\begin{theorem}[线性换元的集合公式]\label{theorem:2-4-2-2}
设 $A:\R^n\to\R^n$ 线性，$E\subset\R^n$ Lebesgue 可测，则 $AE$ Lebesgue 可测，且
\[
 m(AE)=|\det A|m(E).
 \tag{2.4.14}\label{eq:linear-change-set}
\]
若 $A$ 不可逆，右端约定为零；特别地，当 $m(E)=\infty$ 时也约定
$0\cdot\infty=0$。
\end{theorem}

\begin{proof}
若 $A$ 奇异，则 $\det A=0$，例 \ref{ex:2-4-2-3} 已给出 $AE$ 可测且零测。

设 $A$ 可逆。Gaussian 消元用有限次行交换、非零行伸缩和把一行的倍数加到另一行，把 $A$ 化为
单位矩阵。等价地，$A$ 是有限个坐标置换、坐标伸缩和剪切矩阵的乘积。引理
\ref{lemma:2-4-2-elementary-volume} 可以依次应用于中间像集，得到
\[
 m(AE)=\left(\prod_{r=1}^N c_r\right)m(E),
\]
其中置换、反射和剪切的 $c_r=1$，坐标伸缩的 $c_r$ 是相应伸缩系数的绝对值。另一方面，命题
\ref{proposition:2-4-2-1} 逐项应用，结合三类基本矩阵的直接行列式计算，给出
\[
 |\det A|=\prod_{r=1}^Nc_r.
\]
两式合并即为 \eqref{eq:linear-change-set}。每一步都保持 Lebesgue 可测性，所以中间像与最终像
均可测。
\end{proof}

特别地，对每个正体积半开盒 $Q$，定理给出
\[
 \frac{m(AQ)}{m(Q)}=|\det A|,
\]
从而定义 \ref{def:2-4-2-3} 中的 $J_m(A)$ 的确等于 $|\det A|$，且不依赖所选盒。

\begin{proposition}[奇异值体积公式]\label{proposition:2-4-2-3}
若 $A=U\Sigma V^T$ 是实矩阵的奇异值分解，奇异值为 $s_1,\dots,s_n\ge0$，则
\[
 |\det A|=\prod_{j=1}^ns_j.
\]
单位球在 $A$ 下的像是主半轴长度为 $s_1,\dots,s_n$ 的椭球；若某个 $s_j=0$，这个椭球退化到
低维子空间。
\end{proposition}

\begin{proof}
正交性给 $U^TU=I$ 与 $V^TV=I$。取行列式并使用乘法性，
\[
 (\det U)^2=(\det V)^2=1,
\]
故 $|\det U|=|\det V|=1$。又
$\Sigma=\operatorname{diag}(s_1,\dots,s_n)$，所以
\[
 |\det A|=|\det U|\,|\det\Sigma|\,|\det V^T|
 =\prod_{j=1}^ns_j.
\]
$V^T$ 和 $U$ 都保持欧氏距离，$\Sigma$ 则把第 $j$ 个坐标方向伸缩 $s_j$ 倍；依次作用于单位球便
得到所述椭球。
\end{proof}

\begin{figure}[htbp]
  \centering
  \begin{tikzpicture}[scale=0.88]
    \begin{scope}
      \draw[book line] (-1,-1) rectangle (1,1);
      \draw[BookInk!35] (-1,0)--(1,0) (0,-1)--(0,1);
      \node[book label] at (0,-1.35) {单位盒};
    \end{scope}
    \draw[book accent arrow] (1.25,0) -- node[above,book label] {旋转} (2.35,0);
    \begin{scope}[xshift=3.65cm,rotate=20]
      \draw[book line] (-1,-1) rectangle (1,1);
    \end{scope}
    \node[book label] at (3.65,-1.35) {$|\det|=1$};
    \draw[book accent arrow] (4.9,0) -- node[above,book label] {轴向伸缩} (6.15,0);
    \begin{scope}[xshift=7.45cm]
      \draw[book line] (-1.35,-0.7) rectangle (1.35,0.7);
      \node[book label] at (0,-1.35) {乘 $s_1s_2$};
    \end{scope}
    \draw[book accent arrow] (8.95,0) -- node[above,book label] {剪切} (10.1,0);
    \begin{scope}[xshift=11.45cm]
      \draw[book line] (-1.55,-0.7)--(1.15,-0.7)--(1.55,0.7)--(-1.15,0.7)--cycle;
      \node[book label] at (0,-1.35) {体积不再改变};
    \end{scope}
  \end{tikzpicture}
  \caption{线性体积变化可以分解为正交变换、轴向伸缩与剪切。只有轴向伸缩贡献绝对行列式的大小。}
  \label{fig:2-4-2-linear-volume}
\end{figure}

\begin{proposition}[Jacobian 的局部体积比]\label{corollary:2-4-2-4}
设 $f$ 在 $x$ 的邻域为 $C^1$，且 $Df(x)$ 可逆。逆函数定理保证 $f$ 在更小邻域内单射。对固定的
有界正体积盒 $Q$，其边界零测，则
\[
 \frac{m(f(x+rQ))}{r^nm(Q)}\longrightarrow |\det Df(x)|
 \qquad(r\downarrow0).
 \tag{2.4.15}\label{eq:future-local-jacobian-volume}
\]
这是一个明确延后证明、且本章后续不调用的前瞻命题。证明见\cref{theorem:3-2-2-2}：第~3~章先从
本节已经证明的全局线性体积公式独立建立近恒等映射的局部体积夹逼，再证明集合换元公式；该证明不使用
本命题，因而不存在循环依赖。集合换元公式随后把分子写成 \(J_f\) 在 \(x+rQ\) 上的积分，而
\(J_f\) 的连续性给出极限。若只假设 \(f\) 在 \(x\) 一点可微，则邻域单射性和像集无重叠性都没有
保证。
\end{proposition}

行列式的符号适合记录定向，测度则只能使用其绝对值。非线性映射还可能在不同点把质量送到同一处，
因此线性集合公式并不自动推出非线性换元公式。下一小节只处理一种跨所有尺度的距离控制，并证明它
足以保住零测集。

Gaussian 分布在线性映射下的协方差会按 $A\Sigma A^{T}$ 变换；在有密度的非退化情形中，归一化常数
相应出现 $|\det A|$。Fourier 变换和卷积的线性缩放也会产生同一行列式因子。期望、密度、Fourier
变换与卷积将在积分理论中定义，并以集合公式 \eqref{eq:linear-change-set} 为线性换元的基础。

\begin{exercise}[剪切的外测度版本]\label{exercise:2-4-2-1}
仿照 \eqref{eq:shear-outer-measure} 的证明，写出三维映射
$(x,y,z)\mapsto(x+ay+bz,y,z)$ 的细盒覆盖，并证明它保持每个集合的 Lebesgue 外测度。
\end{exercise}

\begin{exercise}[仿射子空间]\label{exercise:2-4-2-2}
用引理 \ref{lemma:2-4-2-lower-dimensional-null} 证明：$\R^n$ 中每个维数小于 $n$ 的仿射子空间
都是 Lebesgue 零集。
\end{exercise}

\begin{exercise}[椭球体积]\label{exercise:2-4-2-3}
设 $B=\{x:|x|<1\}$，$A$ 可逆。证明
$m(AB)=|\det A|m(B)$，再用奇异值分解把 $AB$ 写成一个主轴椭球并核对其体积因子。
\end{exercise}

\begin{exercise}[有向面积与普通面积]\label{exercise:2-4-2-4}
计算反射 $(x,y)\mapsto(-x,y)$ 对有序基 $(e_1,e_2)$ 的作用。分别说明它怎样改变平行四边形的
有向面积和 Lebesgue 面积。
\end{exercise}

\subsection{Lipschitz 几何、零测集与可求长曲线}\label{subsec:2-4-3}
\index{Lipschitz 映射}\index{Hausdorff 测度}\index{可求长曲线}

连续性只要求很近的点被送到很近的点，却不规定两种“很近”以什么速度对应。下面两个构造展示同一
缺口的两种强度：第一个把零测集映成一个区间，第二个把一条参数区间映满一个正面积方形。它们都
连续；失败发生在尺度估计上。

\begin{example}[Cantor 函数不是 Lipschitz]\label{ex:2-4-3-1}
从 $C_0=[0,1]$ 开始，每一步从每个闭区间删去开中三分之一，记剩余的 $2^k$ 个闭区间之并为
$C_k$，并令 $C=\bigcap_kC_k$。因为
\[
 m(C_k)=2^k3^{-k}=\left(\frac23\right)^k,
\]
而 $C_k\downarrow C$ 且 $m(C_0)=1<\infty$，测度从上连续给
\[
 m(C)=\lim_{k\to\infty}m(C_k)
 =\lim_{k\to\infty}\left(\frac23\right)^k=0.
\]

每个 $x\in C$ 可以选择一个只含数字 $0,2$ 的三进制展开
$x=\sum_{j\ge1}a_j3^{-j}$。定义
\[
 F(x)=\sum_{j\ge1}\frac{a_j}{2}\,2^{-j}\qquad(x\in C),
 \tag{2.4.16}\label{eq:cantor-digit-map}
\]
并在 $[0,1]\setminus C$ 的每个被删开区间上取与两个端点相同的常值。对有两种三进制展开的
端点，选择只含 $0,2$ 的那一种；每个被删区间的两个端点经 \eqref{eq:cantor-digit-map} 恰好得到
同一数值，因此常值延拓无歧义。

先核对单调性。若 $x<y$ 都在 $C$，在它们只含 $0,2$ 的展开首个不同位置 $l$，必有
$a_l(x)=0$、$a_l(y)=2$。因此
\[
 F(y)-F(x)
 \ge2^{-l}-\sum_{j>l}2^{-j}=0.
\]
等号只可能对应某个被删区间的左右端点；在该区间上作常值延拓仍保持次序。因此延拓后的 $F$ 在
$[0,1]$ 上非减。

第 $k$ 层每个基本闭区间上，前 $k$ 个三进制数字固定，故 $F$ 的振幅至多
$\sum_{j>k}2^{-j}=2^{-k}$；被删区间上的振幅为零。

令 $\alpha=\log2/\log3$。若
$3^{-(k+1)}<|x-y|\le3^{-k}$，不同第 $k$ 层基本区间之间的间隙长度至少为 $3^{-k}$，故
$[x,y]$ 至多穿过两个这一级基本区间的有效部分。因此
\[
 |F(x)-F(y)|\le2\cdot2^{-k}\le4|x-y|^\alpha.
\]
这同时严格证明 $F$ 连续，并给出 \Holder{}-$\alpha$ 估计。另一方面，取
$x_k=0$、$y_k=3^{-k}$。数 $3^{-k}$ 有展开
$0.\underbrace{00\cdots0}_{k\text{ 位}}222\cdots$，因而
\[
 |y_k-x_k|=3^{-k},\qquad
 |F(y_k)-F(x_k)|=2^{-k},\qquad
 \frac{|F(y_k)-F(x_k)|}{|y_k-x_k|}=\left(\frac32\right)^k\to\infty.
\]
不存在统一的线性距离常数。最后，任取 $u\in[0,1]$ 的二进制展开
$u=\sum b_j2^{-j}$，把数字 $b_j$ 换成三进制数字 $2b_j$ 便得到 $x\in C$ 且 $F(x)=u$。
于是 $F(C)=[0,1]$：连续映射确实能把零测集送成正测度集。
\end{example}

第二个构造不借用“存在空间填充曲线”这一结论。我们把四个相似变换写出来，让满射从一致极限中
产生。对 $Q=[0,1]^2$ 定义
\[
 \begin{array}{ll}
 S_0(u,v)=(v/2,u/2),
 &S_1(u,v)=(u/2,v/2+1/2),\\[2mm]
 S_2(u,v)=(u/2+1/2,v/2+1/2),
 &S_3(u,v)=(1-v/2,1/2-u/2).
 \end{array}
 \tag{2.4.17}\label{eq:hilbert-similarities}
\]
它们以比例 $1/2$ 把 $Q$ 分别送到左下、左上、右上、右下四个小方形，并安排相邻入口与出口
相接。

\begin{example}[逐层构造的 Peano 曲线]\label{ex:2-4-3-3}
若连续折线 $P:[0,1]\to Q$ 满足 $P(0)=(0,0)$、$P(1)=(1,0)$，定义新折线
\[
 (\mathcal TP)(t)=
 \begin{cases}
 S_0(P(4t)),&0\le t\le1/4,\\
 S_1(P(4t-1)),&1/4\le t\le1/2,\\
 S_2(P(4t-2)),&1/2\le t\le3/4,\\
 S_3(P(4t-3)),&3/4\le t\le1.
 \end{cases}
 \tag{2.4.18}\label{eq:hilbert-operator}
\]
由 \eqref{eq:hilbert-similarities} 直接计算四段的相接值依次为
$(0,0),(0,1/2),(1/2,1/2),(1,1/2),(1,0)$，故 $\mathcal TP$ 连续并保持两个端点。
从 $P_0(t)=(t,0)$ 开始，递归令 $P_{k+1}=\mathcal TP_k$。

每个 $S_i$ 都是相似比 $1/2$ 的映射，所以对任意两条满足端点条件的曲线 $P,R$，
\[
 \|\mathcal TP-\mathcal TR\|_\infty
 \le\frac12\|P-R\|_\infty.
\]
于是
\[
 \|P_{k+1}-P_k\|_\infty
 \le2^{-k}\|P_1-P_0\|_\infty.
\]
对 $m>k$，三角不等式进一步给出
\[
 \|P_m-P_k\|_\infty
 \le\sum_{j=k}^{m-1}\|P_{j+1}-P_j\|_\infty
 \le 2^{1-k}\|P_1-P_0\|_\infty.
\]
故 $(P_k)$ 在一致范数下是 Cauchy 列。$\R^2$ 完备使逐点极限
$P(t)=\lim_kP_k(t)$ 存在；同一尾和与 $k\to\infty$ 给出一致收敛，因此 $P$ 连续。

还需证明极限填满方形。把 $[0,1]$ 按四进制分成 $4^k$ 个闭参数小区间；对每个长度为 $k$ 的词
$w=(i_1,\dots,i_k)\in\{0,1,2,3\}^k$，相应小区间在 $P_k$ 下落入
\[
 Q_w=S_{i_1}\circ\cdots\circ S_{i_k}(Q).
\]
这些 $Q_w$ 恰是 $2^k\times2^k$ dyadic 方格中的全部小方形，每个只出现一次。归纳如下：$k=1$
时这是 \eqref{eq:hilbert-similarities} 的四个像。若长度 $k$ 的词已给出全部父方形，则固定前缀
$w$ 后，四个后继 $wi$ 的像为
\[
 Q_{wi}=S_w(S_i(Q)),\qquad S_w=S_{i_1}\circ\cdots\circ S_{i_k};
\]
它们是 $Q_w$ 的四个 dyadic 子方形。不同前缀属于内部不交的不同父方形，所以不会重复；每个父方形
又被这四个子方形覆盖，归纳完成。四进制参数区间也按同样的前缀规则分成四个嵌套子区间，故词与
参数区间、像方形始终相配。

给定 $z\in Q$，逐层选取包含 $z$ 的嵌套小方形
$Q_{w_k}$，并选取对应的嵌套参数区间 $I_{w_k}$。它们的长度为 $4^{-k}$，故交集只含一点 $t$。
对 $m\ge k$，递归构造保证 $P_m(t)\in Q_{w_k}$。令 $m\to\infty$，再利用
$\diam Q_{w_k}=\sqrt2\,2^{-k}\to0$，得到 $P(t)=z$。因此 $P([0,1])=Q$。

这条曲线不可能满足 $|P(s)-P(t)|\le L|s-t|$。若满足，把 $[0,1]$ 等分成 $N$ 段；每段的像
直径至多 $L/N$。给定 $\eta>0$，每段像可放进边长 $L/N+\eta$ 的半开坐标方形。于是
\[
 1=m(Q)=m^*(P([0,1]))
 \le N\left(\frac LN+\eta\right)^2.
\]
先令 $\eta\downarrow0$ 得 $1\le L^2/N$，再令 $N\to\infty$ 产生矛盾。逐层折线、一致极限与嵌套方格分别保证了可计算性、连续性和满射性。
\end{example}

\begin{definition}[Lipschitz 映射]\label{def:2-4-3-1}
设 $E\subset\R^n$。若 $f:E\to\R^m$ 满足
\[
 |f(x)-f(y)|\le L|x-y|\qquad(x,y\in E),
 \tag{2.4.19}\label{eq:lipschitz-definition}
\]
则称 $f$ 为 $L$-Lipschitz 映射。若定义域的每一点都有一个相对邻域，使 $f$ 在该邻域满足某个
Lipschitz 估计，则称 $f$ 局部 Lipschitz。
\end{definition}

\begin{lemma}[紧集上的统一 Lipschitz 常数]\label{lemma:2-4-3-compact-lipschitz}
若 $f:E\to\R^m$ 局部 Lipschitz，$K\subset E$ 为紧集，则 $f|_K$ 是 Lipschitz 映射。
\end{lemma}

\begin{proof}
若 $K=\varnothing$，结论平凡。以下设 $K\ne\varnothing$。
取有限个相对开集 $U_1,\dots,U_N$ 覆盖 $K$，使 $f$ 在 $U_j$ 上为 $L_j$-Lipschitz。紧度量空间
$K$ 的 Lebesgue 数引理给出 $\delta>0$，使 $K$ 中直径小于 $\delta$ 的集合包含于某个 $U_j$。
因此 $x,y\in K$ 且 $|x-y|<\delta$ 时，
$|f(x)-f(y)|\le L_0|x-y|$，其中 $L_0=\max_jL_j$。局部 Lipschitz 蕴含连续，故
$M=\sup_{x\in K}|f(x)|<\infty$。若 $|x-y|\ge\delta$，则
\[
 |f(x)-f(y)|\le2M\le(2M/\delta)|x-y|.
\]
取 $\max\{L_0,2M/\delta\}$ 即得统一常数。
\end{proof}

线性距离控制还可以跨维数表达。为此只引入本书随后真正会用到的覆盖形式。

\begin{definition}[Hausdorff 内容与测度]\label{def:2-4-3-2}
设 $s\ge0$、$\delta>0$。对度量空间中的 $A$，定义
\[
 \mathcal H^s_\delta(A)=
 \inf\left\{\sum_{i\in I}(\diam U_i)^s:
 A\subset\bigcup_{i\in I}U_i,\ I\text{ 有限或可数},\ \diam U_i\le\delta\right\},
 \tag{2.4.20}\label{eq:hausdorff-delta}
\]
以及
\[
 \mathcal H^s(A)=\lim_{\delta\downarrow0}\mathcal H^s_\delta(A).
\]
删去直径限制得到的下确界记作 $\mathcal H^s_\infty(A)$，称为 Hausdorff 内容。本书采用不含
额外规范化常数的约定。空指标族以成本 $0$ 覆盖空集；$s=0$ 时约定每个非空覆盖块的成本为 $1$、
空覆盖块的成本为 $0$，于是 $\mathcal H^0$ 是计数测度。
\end{definition}

最后一句可由定义直接证明。若 $A$ 含 $N\ge2$ 个不同点 $x_1,\dots,x_N$，令
\[
 d_N=\min_{i\ne j}d(x_i,x_j)>0.
\]
当 $0<\delta<d_N$ 时，每个直径不超过 $\delta$ 的覆盖块至多含一个 $x_i$，所以任一 $A$ 的容许
覆盖至少有 $N$ 个非空块，$\mathcal H^0_\delta(A)\ge N$。若 $A$ 无限，对每个 $N$ 应用此式得到
$\mathcal H^0(A)=\infty$；若 $A$ 有恰好 $N\ge2$ 个点，单点覆盖给反向不等式，故
$\mathcal H^0(A)=N$。单点集的每个覆盖至少含一个非空块且单点覆盖可取，故其值为 $1$；空集由
空覆盖取值 $0$。

\begin{proposition}[Hausdorff 测度是 Borel 测度]\label{proposition:2-4-3-hausdorff-borel}
对每个 $s\ge0$，$\mathcal H^s$ 是度量外测度，因而所有 Borel 集都
$\mathcal H^s$-可测。
\end{proposition}

\begin{proof}
固定 $\delta$，在所有直径不超过 $\delta$ 的集合上取成本 $(\diam U)^s$。覆盖下确界引理
\ref{lemma:2-1-2-1} 说明 $\mathcal H^s_\delta$ 是外测度。随着 $\delta\downarrow0$，这些量单调
增加。对集合列 $(A_j)$，
\[
 \mathcal H^s_\delta\Bigl(\bigcup_jA_j\Bigr)
 \le\sum_j\mathcal H^s_\delta(A_j)
 \le\sum_j\mathcal H^s(A_j).
\]
令 $\delta\downarrow0$，得到 $\mathcal H^s$ 的可数次可加不等式；空集条件与单调性同样从
$\mathcal H^s_\delta$ 传给极限。因此 $\mathcal H^s$ 是外测度。

若 $\dist(A,B)=d>0$，取 $0<\delta<d$。直径不超过 $\delta$ 的覆盖块不可能同时遇到 $A$ 和
$B$。任意一个 $A\cup B$ 的容许覆盖可分成覆盖 $A$ 与覆盖 $B$ 的两个子族，所以
\[
 \mathcal H^s_\delta(A\cup B)
 \ge \mathcal H^s_\delta(A)+\mathcal H^s_\delta(B).
\]
反向不等式来自外测度次可加。令 $\delta\downarrow0$ 得正距离可加性，故
$\mathcal H^s$ 是度量外测度。命题 \ref{proposition:2-1-2-3} 随即给出 Borel
可测性。
\end{proof}

\begin{proposition}[覆盖畸变估计]\label{proposition:2-4-3-1}
若 $f:E\subset\R^n\to\R^m$ 为 $L$-Lipschitz，$A\subset E$，则
\[
 \mathcal H^s(f(A))\le L^s\mathcal H^s(A)\qquad(s>0).
 \tag{2.4.21}\label{eq:hausdorff-lipschitz}
\]
$s=0$ 时，相应结论是像集的点数不超过原集合的点数。若 $m=n$，还有
\[
 m_n^*(f(A))\le n^{n/2}L^n m_n^*(A).
 \tag{2.4.22}\label{eq:lebesgue-lipschitz-distortion}
\]
\end{proposition}

\begin{proof}
对 $U\subset E$，定义直接给出
\[
 \diam f(U)\le L\diam U.
\]
若 $L=0$，像集至多一个点；单点的 $\mathcal H^s$（$s>0$）与 $n$ 维 Lebesgue 外测度都为零，
所以两式均成立。以下设 $L>0$。取 $A$ 的一个直径不超过
$\delta$ 的覆盖；把每个覆盖块与 $E$ 相交后再施加 $f$，得到 $f(A)$ 的直径不超过 $L\delta$ 的
覆盖，且
\[
 \sum_i(\diam f(U_i\cap E))^s
 \le L^s\sum_i(\diam U_i)^s.
\]
先取下确界得到
$\mathcal H^s_{L\delta}(f(A))\le L^s\mathcal H^s_\delta(A)$，再令 $\delta\downarrow0$，得到
\eqref{eq:hausdorff-lipschitz}。$s=0$ 的说法来自函数不能增加集合基数。

证明 \eqref{eq:lebesgue-lipschitz-distortion} 时，若右端无穷则无需再证。否则给定
$\varepsilon>0$，利用命题 \ref{proposition:2-1-3-1a} 取边长为 $r_i$ 的任意细立方体 $Q_i$，使
\[
 A\subset\bigcup_iQ_i,
 \qquad \sum_ir_i^n<m_n^*(A)+\varepsilon.
\]
这里必须考察 $f(Q_i\cap E)$，因为 $f$ 未必定义在整个 $Q_i$ 上。该像集的直径至多
$L\sqrt n\,r_i$。每个坐标函数在该像集上的振幅不超过这个直径；将各坐标值域稍作任意小的
扩大即可得到半开坐标盒。对可数族必须分配可求和误差：再给定 $\eta>0$，对第 $i$ 个非空像集
选盒 $R_i$，使
\[
 |R_i|<(L\sqrt n\,r_i)^n+\eta2^{-i}.
\]
于是
\[
 m_n^*(f(A))
 \le\sum_i|R_i|
 \le n^{n/2}L^n(m_n^*(A)+\varepsilon)+\eta.
\]
先令 $\eta\downarrow0$，再令 $\varepsilon\downarrow0$。
\end{proof}

\begin{theorem}[零测集保持]\label{theorem:2-4-3-2}
若 $f:E\subset\R^n\to\R^n$ 为 Lipschitz 映射，且 $N\subset E$ 满足
$m_n^*(N)=0$，则
\[
 m_n^*(f(N))=0.
\]
特别地，$f(N)$ 自动 Lebesgue 可测且测度为零。
\end{theorem}

\begin{proof}
将 $A=N$ 代入 \eqref{eq:lebesgue-lipschitz-distortion}，得到
\[
 0\le m_n^*(f(N))\le n^{n/2}L^n m_n^*(N)=0.
\]
令 $S=f(N)$。对任意 $A\subset\R^n$，由 $m_n^*(A\cap S)=0$、次可加性和单调性，
\[
 m_n^*(A)
 \le m_n^*(A\setminus S)+m_n^*(A\cap S)
 =m_n^*(A\setminus S)
 \le m_n^*(A).
\]
所以等号成立，$S$ 直接满足 \Caratheodory{} 条件且测度为零。
这个论证没有假设连续像是 Borel 集。
\end{proof}

\begin{example}[Lipschitz 图]\label{ex:2-4-3-2}
设 $n\ge2$，$g:\R^{n-1}\to\R$ 为 $L$-Lipschitz，并记
\[
 \Gamma=\{(x,g(x)):x\in\R^{n-1}\}.
\]
固定 $M$，把 $[-M,M]^{n-1}$ 分成边长不超过 $\delta$ 的小盒；盒数至多
$(2M/\delta+2)^{n-1}$。在每个底盒 $Q$ 中选一点 $x_Q$。若 $x\in Q$，则
\[
 |g(x)-g(x_Q)|\le L\sqrt{n-1}\,\delta.
\]
因此 $Q$ 上方的图可放进一个底边长不超过 $\delta$、高度不超过
$(2L\sqrt{n-1}+1)\delta$ 的 $n$ 维盒。所有这些盒的总体积不超过
\[
 (2M/\delta+2)^{n-1}(2L\sqrt{n-1}+1)\delta^n,
\]
它随 $\delta\downarrow0$ 趋于零，故该图块的 $m_n^*$ 为零。又因 $g$ Lipschitz 从而连续，
\[
 \Gamma=\{(x,t)\in\R^{n-1}\times\R:t-g(x)=0\}
\]
是连续函数 $(x,t)\mapsto t-g(x)$ 的零水平闭集，因而 Borel 可测。故
$m_n(\Gamma\cap([-M,M]^{n-1}\times\R))=0$。再令 $M\in\N$ 并取可数并，得到
$m_n(\Gamma)=0$。

另一方面，参数映射 $x\mapsto(x,g(x))$ 是 $\sqrt{1+L^2}$-Lipschitz；
\eqref{eq:hausdorff-lipschitz} 表明它自然受 $(n-1)$ 维 Hausdorff 测度控制。这不与 $n$ 维零测
矛盾，而是在区分两种维数。
\end{example}

微分提供点态的一阶信息，Lipschitz 条件提供全尺度信息。两者之间的逆向联系由下述经典定理给出。

\begin{remark}[外部背景：Rademacher 定理]\label{theorem:2-4-3-3}
设 $U\subset\R^n$ 开，$f:U\to\R^m$ 局部 Lipschitz，则 $f$ 在 Lebesgue 几乎处处
\Frechet{} 可微。
\end{remark}

Rademacher 定理的高维证明需要 Vitali 型覆盖与 Lebesgue 微分工具，本书略去证明。area 与
coarea 公式还要同时处理 Jacobian 和像点重数，属于几何测度论的进一步内容。这里把
Rademacher 定理作为明确标出的外部背景记录，而不把它纳入本章已经证明的可用结论；本书后文的
任何证明也不以它为前提。

曲线长度可以直接由折线逼近定义。这个定义不依赖积分，并使可求长性成为一个纯粹的几何性质。

\begin{definition}[可求长曲线]\label{def:2-4-3-3}
设 $\gamma:[a,b]\to\R^n$ 连续。对分割
$P=\{a=t_0<t_1<\cdots<t_N=b\}$，令
\[
 L(\gamma,P)=\sum_{j=1}^N|\gamma(t_j)-\gamma(t_{j-1})|,
\]
并定义
\[
 \ell(\gamma)=\sup_P L(\gamma,P).
 \tag{2.4.23}\label{eq:curve-length-partitions}
\]
若 $\ell(\gamma)<\infty$，称 $\gamma$ 可求长。
\end{definition}

\begin{proposition}[弧长参数化]\label{proposition:2-4-3-arclength-param}
若 $\gamma:[a,b]\to\R^n$ 可求长，且 $L=\ell(\gamma)$，则存在
$1$-Lipschitz 映射 $\widetilde\gamma:[0,L]\to\R^n$，使
\[
 \widetilde\gamma([0,L])=\gamma([a,b]).
\]
原曲线停驻的参数区间不会在新参数中制造歧义。
\end{proposition}

\begin{proof}
记 $\ell(\gamma|_{[u,v]})$ 为限制曲线的长度。若 $u\le v\le w$，拼接两侧分割给出
\[
 \ell(\gamma|_{[u,w]})
 \ge\ell(\gamma|_{[u,v]})+\ell(\gamma|_{[v,w]}).
\]
反过来，在 $[u,w]$ 的任意分割中加入 $v$；三角不等式说明加点不减折线和，然后把两侧分别用其
长度控制，故反向不等式也成立。于是长度对相邻子区间可加。

令
\[
 s(t)=\ell(\gamma|_{[a,t]}).
\]
$s$ 单调不减，而且相邻区间可加性给
\[
 s(u)-s(t)=\ell(\gamma|_{[t,u]})\qquad(a\le t\le u\le b).
\]
下面证明它连续。若它在 $t$ 处不右连续，设
$\eta=\lim_{u\downarrow t}\ell(\gamma|_{[t,u]})>0$。递归选择
$u_1>u_2>\cdots>t$。给定 $u_j$，取 $[t,u_j]$ 的分割
$t=t_0<t_1<\cdots<t_N=u_j$，使折线和大于 $5\eta/6$；再利用 $\gamma$ 在 $t$ 连续，选择
$u_{j+1}\in(t,\min\{t_1,t+(u_j-t)/2\})$，使
$|\gamma(u_{j+1})-\gamma(t)|<\eta/6$。删除分割中的 $t$、加入 $u_{j+1}$ 后，三角不等式给
\[
\begin{aligned}
 \ell(\gamma|_{[u_{j+1},u_j]})
 &\ge |\gamma(t_1)-\gamma(u_{j+1})|
      +\sum_{i=2}^N|\gamma(t_i)-\gamma(t_{i-1})|\\
 &\ge \sum_{i=1}^N|\gamma(t_i)-\gamma(t_{i-1})|
      -|\gamma(u_{j+1})-\gamma(t)|
 >\frac{2\eta}{3}.
\end{aligned}
\]
此外 $u_j\downarrow t$，而这些相邻区间互不重叠。反复使用长度可加性，对每个 $N$ 有
\[
 L\ge \ell(\gamma|_{[u_{N+1},u_1]})
 =\sum_{j=1}^N\ell(\gamma|_{[u_{j+1},u_j]})
 >\frac{2N\eta}{3},
\]
这与 $L<\infty$ 矛盾。故 $s$ 右连续；反向参数化同样
证明左连续的细节如下。令
\[
 \gamma^{\leftarrow}(r)=\gamma(a+b-r),\qquad a\le r\le b.
\]
它仍连续且长度为 $L$。把已经证明的右连续结论应用于其长度函数
$s^{\leftarrow}(r)=\ell(\gamma^{\leftarrow}|_{[a,r]})$。若 $u\uparrow t$，则
$a+b-u\downarrow a+b-t$，并且
\[
 s(t)-s(u)
 =\ell(\gamma|_{[u,t]})
 =s^{\leftarrow}(a+b-u)-s^{\leftarrow}(a+b-t)\longrightarrow0.
\]
故 $s$ 左连续。

因此 $s$ 连续，并把 $[a,b]$ 映到 $[0,L]$。若 $s(t_1)=s(t_2)$ 且 $t_1<t_2$，长度可加性给
$\ell(\gamma|_{[t_1,t_2]})=0$，从而
$|\gamma(t_2)-\gamma(t_1)|=0$。所以可对 $u=s(t)$ 定义
\[
 \widetilde\gamma(u)=\gamma(t),
\]
而不依赖 $t$ 的选择。若 $u=s(t)<v=s(r)$，可取 $t<r$，并有
\[
 |\widetilde\gamma(v)-\widetilde\gamma(u)|
 =|\gamma(r)-\gamma(t)|
 \le\ell(\gamma|_{[t,r]})=v-u.
\]
故 $\widetilde\gamma$ 为 $1$-Lipschitz；$s$ 的满射性还说明两条曲线的轨迹相同。
\end{proof}

\begin{example}[折线路径的长度]\label{ex:2-4-3-polygonal}
设 $a=t_0<\cdots<t_N=b$，$\gamma$ 在每段 $[t_{j-1},t_j]$ 上线性，顶点为
$p_j=\gamma(t_j)$。含全部 $t_j$ 的分割给出
\[
 \ell(\gamma)\ge\sum_{j=1}^N|p_j-p_{j-1}|.
\]
对任意分割加入全部 $t_j$。每条线性段上按参数顺序排列的弦同向共线，其长度和等于该段两端距离；
故加细后的折线和不超过右端总和。于是
\[
 \ell(\gamma)=\sum_{j=1}^N|p_j-p_{j-1}|.
\]
弧长参数化就是以单位速度依次走过这些线段；零长度线段被自动略去。
\end{example}

\begin{proposition}[绝对连续曲线的长度]\label{proposition:2-4-3-4}
若 $\gamma:[a,b]\to\R^n$ 绝对连续，则
\[
 \ell(\gamma)=\int_a^b|\gamma'(t)|\,\dd t.
 \tag{2.4.24}\label{eq:future-length-integral}
\]
\end{proposition}

绝对连续性及 Lebesgue 微积分基本定理将在第~4~章建立，公式
\eqref{eq:future-length-integral} 的证明见\cref{proposition:4-2-3-curve-length}。上面的弧长
参数化则只使用分割上确界。

\begin{figure}[htbp]
  \centering
  \begin{tikzpicture}[node distance=7mm and 9mm]
    \node[book strong node] (lip) {Lipschitz\quad $r\mapsto Lr$};
    \node[book node,below=of lip] (lipout) {Hausdorff 维数不增\\同维零集保持};
    \draw[book accent arrow] (lip)--(lipout);

    \node[book warm node,right=of lip] (cantor) {Cantor\quad $r\mapsto r^\alpha$, $\alpha<1$};
    \node[book node,below=of cantor] (cantorout) {$C$ 零测\quad $F(C)=[0,1]$};
    \draw[book arrow] (cantor)--(cantorout);

    \node[book warning node,right=of cantor] (peano) {Peano\quad 四进制层级};
    \node[book node,below=of peano] (peanoout) {$[0,1]$ 的像\quad $[0,1]^2$};
    \draw[book arrow] (peano)--(peanoout);

    \draw[book dashed arrow] (cantor.west) to[bend right=20]
      node[above,book label] {线性尺度失效} (lip.east);
    \draw[book dashed arrow] (peano.west) to[bend right=20]
      node[above,book label] {面积被填满} (cantor.east);
  \end{tikzpicture}
  \caption{三种尺度行为。Hausdorff 覆盖估计精确表达了 Lipschitz 情形的维数控制；另外两列说明
  单独的连续性为何不足。}
  \label{fig:2-4-3-scale-comparison}
\end{figure}

局部 $C^1$ 映射在每个紧盒上导数有界，均值估计与引理
\ref{lemma:2-4-3-compact-lipschitz} 因而给出局部 Lipschitz 控制。特别地，$C^1$ 微分同胚会在
局部把零测异常集送到零测集。Brownian 路径一类远不满足此尺度控制的对象，则需要用 Hausdorff
维数描述其几何大小，而不能依赖可求长参数化。

\begin{exercise}[Lipschitz 图是零集]\label{exercise:2-4-3-1}
补全例 \ref{ex:2-4-3-2} 中底域分割的整数取整细节，并证明局部 Lipschitz 图也在每个有界柱体中
是 $n$ 维零集。
\end{exercise}

\begin{exercise}[Hausdorff 畸变]\label{exercise:2-4-3-2}
从 \eqref{eq:hausdorff-delta} 出发证明内容版本
$\mathcal H^s_\infty(f(A))\le L^s\mathcal H^s_\infty(A)$，并解释为什么证明不要求 $A$ 可测。
\end{exercise}

\begin{exercise}[Lipschitz 曲线的一维大小]\label{exercise:2-4-3-3}
若 $\gamma:[a,b]\to\R^m$ 为 $L$-Lipschitz，用等长小区间覆盖证明
\[
 \mathcal H^1(\gamma([a,b]))\le L(b-a)
\]
（采用定义 \ref{def:2-4-3-2} 的规范化）。
\end{exercise}

\begin{exercise}[空间填充与 \Holder{} 指数]\label{exercise:2-4-3-4}
设 $P:[0,1]\to[0,1]^2$ 满射且满足
$|P(s)-P(t)|\le C|s-t|^\alpha$。把参数区间等分成 $N$ 段并用像集的方形覆盖证明
$\alpha\le1/2$。
\end{exercise}

至此，集合层面的体积、线性缩放和零集控制已经建立。第~3 章将从简单函数开始构造积分，并把推前
测度、局部 Jacobian 估计与零测集保持组合成非线性换元公式。

